\documentclass[11pt,a4paper]{article}
\usepackage[T1]{fontenc}
\usepackage{lmodern}
\usepackage{iftex}
\ifXeTeX
\usepackage{newunicodechar}
\newunicodechar{§}{\S}
\newunicodechar{ö}{\"o}
\newunicodechar{é}{\'e}
\newunicodechar{ł}{\l}
\newunicodechar{î}{\^i}
\newunicodechar{∎}{\ensuremath{\square}}
\newunicodechar{¬}{\ensuremath{\neg}}
\newunicodechar{χ}{\ensuremath{\chi}}
\newunicodechar{∧}{\ensuremath{\wedge}}
\newunicodechar{κ}{\ensuremath{\kappa}}
\newunicodechar{β}{\ensuremath{\beta}}
\newunicodechar{Θ}{\ensuremath{\Theta}}
\newunicodechar{ℓ}{\ensuremath{\ell}}
\newunicodechar{₀}{\ensuremath{{}_0}}
\newunicodechar{≤}{\ensuremath{\le}}
\newunicodechar{≥}{\ensuremath{\ge}}
\newunicodechar{−}{\ensuremath{-}}
\fi
\usepackage[margin=23mm,top=22mm,bottom=23mm]{geometry}
\usepackage{microtype}
\usepackage{amsmath,amssymb,amsthm,mathtools,mathrsfs}
\usepackage{booktabs,longtable,array,calc}
\usepackage{etoolbox,needspace}
\usepackage{xcolor}
\definecolor{linkblue}{RGB}{30,62,95}
\usepackage[colorlinks=true,linkcolor=linkblue,urlcolor=linkblue,citecolor=linkblue]{hyperref}
\hypersetup{pdftitle={The quartic QRH boundary},pdfsubject={Quartic zero-free boundary for Hecke and Dirichlet L-functions},pdfauthor={Hailey Collet}}
\urlstyle{same}
\setlength{\parindent}{1.1em}
\setlength{\parskip}{3pt plus 1pt minus 1pt}
\setlength{\emergencystretch}{2.2em}
\newcommand{\sourcecode}[1]{{\small\ttfamily #1}}
\setcounter{secnumdepth}{-1}
\allowdisplaybreaks[1]
\raggedbottom
\providecommand{\tightlist}{\setlength{\itemsep}{0pt}\setlength{\parskip}{0pt}}
\providecommand{\pandocbounded}[1]{#1}
\pretocmd{\section}{\Needspace{10\baselineskip}}{}{}
\pretocmd{\subsection}{\Needspace{5\baselineskip}}{}{}
\title{\vspace{-1.3em}The quartic QRH boundary}
\author{Hailey Collet\\\small\href{mailto:HaileyCollet@gmail.com}{HaileyCollet@gmail.com}}
\date{October 10, 2026}
\begin{document}
\maketitle
\vspace{-1.5em}
\begin{abstract}
We prove nonvanishing for finite-order Hecke $L$-functions of trivial infinite type over $\mathbf{Q}(\sqrt{-3})$, and for all Dirichlet $L$-functions, in the strict half-plane $\Re s>b_0$, apart from principal poles. Here $b_0=11/12-\ell_0/4=0.8748866470399641\ldots$, where $\ell_0$ is the root of $927X^4-3135X^3+2433X^2+275X-100$ in $[1/6,1/5]$. The corresponding value is $c_0=1-b_0=0.1251133529600359\ldots$. The proof uses a variable-parameter moment estimate, the physical weighted numerator and an exact quartic optimization, with the established $7/8$ theorem as initialization. An ordinary Lean formalization supplies the exact boundary and a rational corollary at $\theta=683505193/781250000$.
\end{abstract}

\hypertarget{t0.-exact-statement-and-a-noncircular-initialization}{%
\section{T0. Exact statement and a noncircular
initialization}\label{t0.-exact-statement-and-a-noncircular-initialization}}

Let \(K=\mathbf Q(\omega)\), \(\omega=e^{2\pi i/3}\), and
\(\mathcal O=\mathbf Z[\omega]\). Let \(\ell_0\) be the unique root in
\[
I=(16712007850680/10^{14},16712007850682/10^{14})
\] of \(p(X)=927X^4-3135X^3+2433X^2+275X-100\). Set \[
b_0=11/12-\ell_0/4,\qquad c_0=1-b_0.
\]

\textbf{Theorem T0.} Every finite-order Hecke character of trivial
infinite type over \(K\), in any nonzero integral modulus, has nonzero
ordinary \(L\)-function for \(b_0<\Re s\), except at its principal pole.
Consequently, for every \(q\ge1\), every
\(\chi\in\operatorname{DirChar}(\mathbf C,\allowbreak{}q)\), and every
\(s\in\mathbf C\), \[
b_0<\Re s,\quad\neg(\chi=1\land s=1)
\quad\Longrightarrow\quad L(s,\chi)\ne0.
\] The boundary is strict; no assertion of nonvanishing on \(\Re s=b_0\)
is made.

We use the ordinary Hecke objects of the formal source. In OpenAI/math
at \sourcecode{fd4aeeb2\allowbreak{}ee4fc729\allowbreak{}c18d9844\allowbreak{}4fed42fd\allowbreak{}0529eeeb},
\sourcecode{OAI.\allowbreak{}Seven\allowbreak{}Eighths.\allowbreak{}Hecke\allowbreak{}Family.\allowbreak{}LFunction\_\allowbreak{}ne\_\allowbreak{}zero\_\allowbreak{}of\_\allowbreak{}seven\_\allowbreak{}eighths\_\allowbreak{}lt\_\allowbreak{}re}
proves nonvanishing for every \sourcecode{Hecke\allowbreak{}Family.\allowbreak{}Character} when
\(7/8<\Re s\), with the principal pole excluded. We use this coarse
theorem to initialize the family supremum.

For exact agreement with
\sourcecode{OAI.\allowbreak{}Seven\allowbreak{}Eighths.\allowbreak{}Hecke\allowbreak{}Zero\allowbreak{}Supremum.\allowbreak{}beta}, take the supremum of
\(\{1/2\}\) together with the real parts of all positive-real-part Hecke
zeros outside the principal pole. The set is nonempty and bounded above.
Primitive reduction in T2 shows that imprimitive characters add no zeros
in \(\Re s>0\). The upstream theorem therefore gives \(\beta_*\le7/8\).
The source's primitive-family supremum and this definition agree
whenever \(\beta_*>b_0\); the sentinel is then irrelevant.

Suppose for contradiction that \(\beta_*>b_0\), and put
\(\kappa=2\beta_*-1\). Then \[
7/10<\kappa_0=2b_0-1<\kappa\le3/4.
\] Only this compact range is needed for the modified positive-slot
moment estimate. No assumption \(\beta_*\le b_0\) occurs. The hypothesis
\(\beta_*\le(1+\kappa)/2\) used to estimate the small prime slots is an
equality by this definition. B3 proves the moment induction with this
hypothesis and a positive comparison margin, not with the desired
improved conclusion.

The only initial zero-free bound used in the argument is
\(\beta_*\le7/8\). The inverse amplification estimates below remain
necessary for the nonprincipal rows.

\hypertarget{t1.-dependency-order-conventions-and-exact-endpoint}{%
\section{T1. Dependency order, conventions and exact
endpoint}\label{t1.-dependency-order-conventions-and-exact-endpoint}}

Read B1 for finite arithmetic; B2 for completed reflection, local
branches and its energy bound; B3 for inverse and plain moments; B4 for
physical Euler factors and coefficient mass. B0 assembles these results
using the detector, signed Gram bound, full numerator inversion and a
uniform improvement of the zero supremum. The forward references are
acyclic.

The normalizations are fixed throughout: ideal columns use one primary
generator, element rows keep all six units, \(q_a=|a|^2\), the additive
character is \(e(z)=\exp(2\pi i\operatorname{Tr}(z/\lambda))\), and the
measure is \(2\,\allowbreak{}dx\,\allowbreak{}dy/\sqrt3\). Symbols at nonunits are zero even for
exponent zero or six. The Hecke ideal-series normalization is ordinary
\(L\), not completed \(L\); upstream \sourcecode{Hecke/\allowbreak{}Ideal\allowbreak{}Bridge.\allowbreak{}lean}
divides its lattice continuation by six. No extra factor six is
introduced into the unique frequency factorization \(H=ur^6\).

The following elementary facts complete the exact numerical part of
B0.6. The formal counterpart is \sourcecode{Cycle25.\allowbreak{}Arithmetic.\allowbreak{}Endpoint},
using exact rational and polynomial identities.

\textbf{T1.1 (root).} Exact substitution gives
\(p(\inf I)<0<p(\sup I)\). On \([1/6,\allowbreak{}1/5]\), \[
p'(x)=3708x^3-9405x^2+4866x+275
\ge-9405/25+4866/6+275=3549/5>0.
\] The intermediate value theorem and strict monotonicity prove
existence and uniqueness of the chosen root. Other roots of the quartic
are irrelevant.

\textbf{T1.2 (geometry without numerical optimization).} In B0.45--B0.49
use \[
\begin{split}
\kappa_0&=(5-3\ell_0)/6,\\
\delta_0&=\frac{25-60\ell_0+27\ell_0^2}{6(5+4\ell_0-3\ell_0^2)},\\
C_0&=\frac{6\kappa_0-1}{15\kappa_0-2},\quad V_0=5/6-\delta_0,\\
R'_0&=-1+\frac{C_0(V_0^2-\delta_0^2C_0)}{2(V_0+\delta_0C_0)^2},\\
h_0&=-\frac{(1+\ell_0)/2+\ell_0/(6\kappa_0)}{3R'_0/4-1/(24\kappa_0)},\quad
x_0=1+\ell_0-h_0,\quad y_0=h_0-2\ell_0.
\end{split}
\] Every denominator is separated from zero by the same rational
interval calculation. Reducing polynomials modulo \(p\) verifies both
binding equations B0.48 and the full two-variable identity B0.49. The
latter identity, with its strictly positive coefficient minorants,
proves the bound on the entire parameter rectangle; a floating-point
grid would not suffice. In particular, \(36-174\delta+347\delta^2>0\)
follows by completing the square, since \(36-174^2/(4\cdot347)>0\). The
nonnegative terms in B0.49 prove \(F^{(4)}_{\kappa_0}\le b_0\), with
equality only at its binding point.

\textbf{T1.3 (uniform perturbation).} On \(0\le p\le1/2\),
\(\kappa\ge7/10\), the counting coefficient of §B0.5 has derivative \[
\partial_\kappa C_\kappa(p)
=\frac{6p(1-p)^2}{[3\kappa(3-p)-2p]^2}
\le\frac{8/9}{[(15/2)(7/10)-1]^2}<1/18.
\] For \(V=5/6-\delta\ge0\),
\(\partial_C R=\delta V^2/[2(V+\delta C)^2]\le5/12\). The exact
intervals give \(h_0<4/5\) and \(0<h_0-4\ell_0<13/100\). Differentiating
B0.47 now gives \[
\partial_\kappa F^{(4)}
\le\frac34\frac45\frac5{12}\frac1{18}
 +\frac5{12}\frac{13/100}{12(7/10)^2}<1/10.
\] As \(\kappa-\kappa_0=2(\beta_*-b_0)\), \[
F^{(4)}_\kappa-\beta_*\le-\frac45(\beta_*-b_0).
\] The other dyad endpoint, floor bin and small rows have respectively
margins exceeding \(.06,\allowbreak{}.01,\allowbreak{}.05\), all checked with rational intervals.
The same exact interval identities check the supply and geometry
inequalities. Thus the strict losses can be chosen uniformly before a
target character is fixed.

\textbf{T1.4 (minimal generalized damping interface).} For
\(\kappa\ge7/10\), \(t,\allowbreak{}z,\allowbreak{}\delta,\allowbreak{}q\ge0\), \(q\le\delta/2\), let
\(z_\kappa(t)=\max(0,\allowbreak{}z-t/(3\kappa))\). Since
\(0\le z-z_\kappa(t)\le t/(3\kappa)\), \[
-\delta t/4-qz_\kappa(t)/2
\le-qz/2-\delta t\bigl(1/4-1/(12\kappa)\bigr)\le-qz/2.
\] This is the variable-\(\kappa\) form of the weaker damping used by
the weighted-numerator Lean route. B0.5 proves its own stronger
\(\delta/2\)-weighted coefficient statement. The physical coefficient
dictionary and combined-conductor mass bound are proved in
\sourcecode{Cycle25.\allowbreak{}Numerator.\allowbreak{}Imported.\allowbreak{}Numerator\allowbreak{}Local\allowbreak{}Coefficients},
\sourcecode{Cycle25.\allowbreak{}Numerator.\allowbreak{}Imported.\allowbreak{}Numerator\allowbreak{}Physical\allowbreak{}Packet\allowbreak{}Mass},
\sourcecode{Cycle25.\allowbreak{}Numerator.\allowbreak{}Reflection\allowbreak{}Mass}, and
\sourcecode{Cycle25.\allowbreak{}Numerator.\allowbreak{}Actual\allowbreak{}Bin}; these supply the scalar-to-source
transfer.

\textbf{T1.5 (rational target).} The root interval implies \[
b_0<\theta=683505193/781250000,
\quad 0<\theta-b_0<4\cdot10^{-14}.
\] Exact cross-multiplication proves \[
\theta<10499/12000
<874957019420098946128604623/10^{27}<7/8.
\] The second intermediate value is the pinned public catalogue
threshold. No finite decimal approximation to \(b_0\) is treated as its
definition. Proving T0 implies the rational target by transitivity of
the strict inequalities.

\hypertarget{analytic-argument}{%
\section{Analytic argument}\label{analytic-argument}}

B0 gives the compensated-probe argument. Appendices B1--B4 supply finite
arithmetic, completed reflection, moments and local coefficient
estimates. The classical theta, sieve and functional-equation inputs are
identified in the references and in the named source declarations below.

\hypertarget{b0.-the-compensated-probe-proof}{%
\section{B0. The compensated-probe
proof}\label{b0.-the-compensated-probe-proof}}

\textbf{Theorem B0 (Hecke boundary).} Let \[
c_0=\frac1{12}+\frac{\ell_0}{4}=0.1251133529600359\ldots,
\quad b_0=1-c_0,
\quad927\ell_0^4-3135\ell_0^3+2433\ell_0^2+275\ell_0-100=0,
\tag{B0.1}
\] where \(0.16712007850680<\ell_0<0.16712007850682\). Every primitive
finite-order Hecke \(L\)-function over \(\mathbb K=\mathbb Q(\omega)\)
is zero-free in \(\Re s>b_0\), with its principal pole allowed. In
particular, \(\zeta(s)\ne0\) for \(\Re s>b_0\), with its pole at one
excluded.

We prove this theorem using the dependency order in T1 and the weaker
initialization in T0. Reflection bounds a compensated probe; Poisson
summation exposes its reciprocal Hecke principal term. Moments and the
exact optimization prevent the remaining rows from cancelling that term.
Appendices B1--B4 supply the finite arithmetic, reflection, moments and
local factors. Notation in these appendices is local: in particular,
\(M_u\) denotes the sextic probe rather than the cubic polynomial in the
main text.

\hypertarget{b0.1.-notation-and-the-four-estimates-used-in-the-proof}{%
\subsection{B0.1. Notation and the four estimates used in the
proof}\label{b0.1.-notation-and-the-four-estimates-used-in-the-proof}}

Let \(\mathcal O=\mathbf Z[\omega]\), \(\omega=e^{2\pi i/3}\),
\(\lambda=1+2\omega\), \(q_a=Na\), \(\alpha(a)=a/|a|\), and \[
e(z)=\exp(2\pi i\operatorname{Tr}_{\mathbb K/\mathbf Q}(z/\lambda)),
\qquad d\mu(z)=2\,dx\,dy/\sqrt3.
\] Fourier transforms use \(e(-zy)\). Ideal sums use the unique primary
generator congruent to one modulo \(3\mathcal O\); element rows retain
six units. The fixed excluded set \(S\) contains primes above six, fixed
conductors and a sufficiently large fixed small-prime cutoff. Columns
avoid \(S\). A star denotes squarefree; \((6)\) denotes row valuations
\(0,\allowbreak{}\ldots,\allowbreak{}5\).

The sextic symbol \(\chi_n\) is multiplicative in its denominator. Every
character power, including zero and six, vanishes at nonunits. Put
\(\psi_u(n)=\chi_n(u)\), and, for squarefree \(n\), \[
\gamma_j(n)=q_n^{-1/2}\sum_{d\bmod n}\chi_n(d)^j e(d/n),
\qquad a_\nu(n)=\overline{\alpha(n)}\gamma_2(n)\nu(n).
\] Appendix B1 proves the finite identities \[
\begin{gathered}
\chi_b(a)=R(a,b)\chi_a(b),\quad R(a,b)=R(b,a)\in\{\pm1\},\quad(a,b)=1,\\
\gamma_2(n)^3=\mu(n)\alpha(n),\quad
\gamma_1(n)\gamma_2(n)=\mu(n)\alpha(n)G(n),\\
\mu(n)\gamma_{-1}(n)=\chi_n(-1)G(n)^{-1}\overline{\alpha(n)}\gamma_2(n),\\
a_\nu(ab)=a_\nu(a)a_\nu(b)\chi_b(a)^4\quad((a,b)=1).
\end{gathered}\tag{B0.2}
\] For odd \(a+b\omega\), one convention is
\(\Gamma_q(a+b\omega)=(1+i^{-b}+i^a+i^{b-a})/2\),
\(G(c)=\overline{\chi_c(4)}\Gamma_q(c)\), and
\(R(v,\allowbreak{}w)=\Gamma_q(vw)/(\Gamma_q(v)\Gamma_q(w))\). Here \(\gamma_{-1}\)
means the inverse-character Gauss sum. Fixed finite phases need not have
order dividing six.

For \(\varphi_u=\nu\psi_u^{\pm1}\), define \[
M_u(D;W)=D^{-1/2}\sum_n\mu(n)\varphi_u(n)W(q_n/D),\qquad
S_u(N;W)=N^{-1/2}\sum_n\varphi_u(n)W(q_n/N).
\tag{B0.3}
\] Profiles are common smooth functions on fixed annuli; bounded real
norm powers and independent height twists have finite seminorm and
polynomial height costs. Moving arithmetic coefficients are never
absorbed into these costs.

For disjoint row-independent prime supports \(\mathcal P_i\), with
\(q_p\asymp U^{z_i}\), set \[
Q_{u,i}=U^{-z_i/2}\sum_{p\in\mathcal P_i}a_i(p)\varphi_u(p)W_i(q_p/U^{z_i}),
\quad Q_{u,I}=\prod_{i\in I}Q_{u,i},\quad z_I=\sum_{i\in I}z_i.
\] The slot count is fixed before \(U\to\infty\). The following
estimates are proved in Appendix B3. Here \(\beta_*\) is the family
supremum in T0.

\begin{itemize}
\tightlist
\item
  \textbf{UM:} \(\sum_{0<q_u\le H}|M_u(D)|^2\ll H(HD)^\epsilon\) for
  \(H\ge D^{1+\vartheta}\), every fixed \(\vartheta,\allowbreak{}\epsilon>0\), and
  fixed polynomial scale ranges. All rows, including sixth powers,
  occur.
\item
  \textbf{MI:} for arbitrary bounded row-independent \(a_i\), \[
  r+2z_I\le1-c_1,\quad2r+8z_I\le3-c_2
  \ \Longrightarrow\ 
  \sum_{0<q_u\ll U}|M_u(U^r)Q_{u,I}|^2\ll U^{1+\epsilon},
  \tag{B0.4}
  \] for fixed positive margins \(c_1,\allowbreak{}c_2\), either common orientation
  and every subcollection. There is no mesh condition.
\item
  \textbf{MP:} the prime sums run over the genuine canonical good-prime
  pools of a fixed ray modulus, with one common smooth annular profile
  \(W\). Each slot coefficient is an individual fixed ray character; its
  bounded real norm power and height may vary by slot. On the original
  natural eligible rows, \[
  2m+6\kappa z_I\le1\ \Longrightarrow\
  \sum_{u\ {\rm eligible}}|S_u(U^m)|^4|Q_{u,I}|^2\ll U^{1+\epsilon}.
  \tag{B0.5}
  \] Assume \(7/10\le\kappa\le3/4\), \(\beta_*\ge51/100\), and
  \(2\beta_*-1\le\kappa\). Eligibility, the row zeros and the common
  puncture are exactly those described in B3.4. The mesh depends only on
  the bounded width, puncture and length ranges and the requested loss;
  it is chosen before the arithmetic family, eventual slot count,
  profiles and target height. Constants and finite seminorm and height
  orders may depend on the subsequently fixed slot system and common
  profile. With no slots, \textbf{MP0} holds for all bounded \(m\),
  excluding only the principal inducing character, without a zero-free
  hypothesis.
\end{itemize}

The contradiction uses \(\kappa=2\beta_*-1\), which satisfies these
hypotheses by T0.

The completed theta sum is \[
\mathcal C_V(N;m,\nu)=\sum_{c\ {\rm sf},n}
\frac{\overline{\alpha(cn^3)}\gamma_2(c)\nu(cn^3)\chi_{cn^3}(m)}
{\sqrt{q_c}\,q_n}V(q_cq_n^3/N).
\tag{B0.6}
\] The indices \(c,\allowbreak{}n\) may share primes. A mark means divisibility of
the \textbf{entire} index \(cn^3\). The unmarked energy is
\((HN)^\epsilon(H+H^2/N)\).

\textbf{RE} is the marked reflection estimate proved in Appendix B2. We
use its consequence (B0.26) below. A mark is imposed on the whole
completed index; its prime support may meet the row. The proof keeps
every cusp source, the fixed ray phases and the natural character zeros.

\hypertarget{b0.2.-finite-reconstruction-and-common-zero-witnesses}{%
\subsection{B0.2. Finite reconstruction and common zero
witnesses}\label{b0.2.-finite-reconstruction-and-common-zero-witnesses}}

The original zero extensions give

\[
\varphi_{ua^6}(n)=\varphi_u(n)\mathbf1_{(n,a)=1}.
\]

Separating the primes of a squarefree column which divide \(a\) gives
the exact finite identity

\[
M_u(D;W)=\sum_{d\mid\operatorname{rad}(a)}
\mu(d)\varphi_u(d)q_d^{-1/2}M_{ua^6}(D/q_d;W).
\tag{B0.7}
\]

There is no requirement \((u,\allowbreak{}a)=1\). The character zeros make (B0.7)
valid also at shared primes. The coefficient mass is at most
\(d(a)\ll q_a^\varepsilon\), where \(d(a)\) counts the ideal divisors of
\(a\).

Differentiating in log scale replaces \(W\) by \(-W/2-yW'\). Sobolev on
unit log intervals, followed by UM for these common profiles, therefore
gives

\[
\sum_{0<q_v\ll H}\sup_{0<E\le D}|M_v(E;W)|^2
\ll H(HD)^\varepsilon,
\qquad H\ge D^{1+\vartheta}.
\tag{B0.8}
\]

The range of subunit scales contains only boundedly many columns and is
empty below a fixed positive scale. The \(O(\log(2+D))\) intervals cost
only a small power. This argument never assigns a moving arithmetic
coefficient to a common-profile moment.

For sixth-power-free \(q_u\asymp U\), choose

\[
H=\max(2U,D^{1+\vartheta}),\qquad P=(H/U)^{1/6}.
\]

There are \(\gg_S P\) primary ideals \(a\), prime to \(S\), with
\(q_a\le P\). The map \((u,\allowbreak{}a)\mapsto ua^6\) is injective: valuations
modulo six recover \(u\), the quotients recover \(a\), and its primary
generator fixes the unit of \(u\). Averaging (B0.7), summing the rows,
and applying (B0.8) gives

\[
\sum_{q_u\asymp U}^{(6)}|M_u(D;W)|^2
\ll U^{1/6}H^{5/6}(UD)^\varepsilon.
\tag{B0.9}
\]

For bounded \(r\) and \(D=U^r\), choose \(\vartheta\) sufficiently small
in terms of the requested final loss. Thus the exponent in the
sixth-power-free inverse moment is \(e(r)=\max\{1,\allowbreak{}(1+5r)/6\}\). Sobolev
in the real power and height variables gives the same bound for separate
rowwise parameters in a fixed real range and a height interval of length
\(T\), at a fixed polynomial cost \((1+T)^A\). Only the smooth tests
vary in that argument.

\hypertarget{buffered-zero-bins}{%
\subsubsection{Buffered zero bins}\label{buffered-zero-bins}}

Fix a finite conjugation-stable group \(\Theta\) of characters supported
on \(S\), containing the numerator, denominator, and fixed phase
presentations needed for the chosen target. For a sixth-power-free row
having a prime factor outside \(S\), define

\[
\mathcal X(u)=\{\nu\psi_u^{\pm1}:\nu\in\Theta\}.
\]

Every member is nonprincipal: a good prime of valuation \(1,\allowbreak{}\ldots,\allowbreak{}5\)
in \(u\) has nontrivial tame inertia. Its primitive conductor norm is
\(O_S(q_u)\).

Let \(a_*=51/100\), take a small buffer \(e>0\), and put
\(I=\lceil1/e\rceil+2\). For a height \(T\), let \(B_j(u)\) be the
maximum of \(a_*\) and the real parts of zeros of members of
\(\mathcal X(u)\) at height at most \(3^jT\). Some \(i<I\) satisfies
\(B_{i+1}-B_i\le e\). Round \(B_i\) down to the \(e\)-grid, obtaining
\(a\), so that

\[
a\le B_i<a+e,\qquad B_{i+1}<a+2e.
\tag{B0.13}
\]

There are only finitely many labels, independently of \(U\). A nonfloor
bin supplies an actual zero \(\rho=\sigma+i\gamma\), with
\(a\le\sigma<a+e\), \(|\gamma|\le3^iT\). Put \(\delta=2a-1\). The floor
bin is counted trivially.

Here is the analytic estimate used in these bins. The functional
equation and Euler convergence give a polynomial conductor-height bound
for \(L_{\mathbb K}\) on a fixed strip. In a zero-free disk about
\(2+it\), continue its Euler logarithm. Borel--Carathéodory bounds the
logarithm by \(O_e(\log(Q(3+|t|)^2))\) on a slightly smaller disk; it is
bounded on the concentric radius \(49/100\) by Euler convergence. Three
circles gives a bound of the form
\(O_e((\log(Q(3+|t|)^2))^{\theta_e})\), with \(\theta_e<1\), on a disk
smaller by another fixed buffer. Exponentiating either sign yields

\[
|L_{\mathbb K}(s,\varphi)|+|L_{\mathbb K}(s,\varphi)^{-1}|
\ll_{e,\varepsilon}\{2Q(3+|t|)^2\}^{\varepsilon}
\tag{B0.14}
\]

on the retained zero-free region. Cauchy's formula gives the
corresponding bounds for finitely many derivatives of
\(\log L_{\mathbb K}\). A principal pole is handled by applying the
logarithm argument to \((s-1)\zeta_{\mathbb K}(s)/(s+1)\).

Deleting primes of radical norm \(O_S(U)\) costs \(U^\varepsilon\) on a
fixed positive half-plane. On \(\Re s\ge-v\), the direct deleted product
costs \(U^{v+\varepsilon}\). These follow primewise from
\(\log|1-z|^{\pm1}=O(|z|)\), treating a fixed finite set separately.

Mellin inversion, (B0.14), and the primitive functional equation give,
for the required bounded real powers and retained heights,

\[
|M_u(U^r)|^2\ll U^{\delta r+\varepsilon_0},\qquad
|S_u(U^m)|^2\ll U^{\delta\min(m,1-m)+\varepsilon_0}.
\tag{B0.15}
\]

For the inverse polynomial, move the retained Mellin segment only to
\(\Re s=a+6e\). For the plain polynomial, use either that line or
\(1-a-6e\); reflection supplies \(U^{a-1/2+O(e)}\) on the second line.
This gives exactly the two exponents in (B0.15). The original character
zeros are restored as deleted Euler factors, not removed from the
polynomials.

For a centrally normalized smooth prime polynomial of length \(P\), the
same logarithm estimate and the Dirichlet series of \(-L'/L\) give

\[
\left|P^{-1/2}\sum_{\mathfrak p}\overline{\psi_u(\mathfrak p)}
w(q_{\mathfrak p}/P)(q_{\mathfrak p}/P)^{it}\right|
\ll P^{\delta/2+O(e)+\varepsilon}U^\varepsilon(1+T)^A.
\tag{B0.16}
\]

Insert the factor \(1/\log q_{\mathfrak p}\) into the smooth test after
using \(-L'/L\). Prime powers contribute \(O(P^{1/2+\varepsilon})\)
before central normalization and are harmless. Primes in the fixed
exclusion are removed; the row's nonunit zeros remain. Formula (B0.16)
needs only the numerator presentation and its conjugate, both included
in \(\mathcal X(u)\).

For precision about heights, first fix the real losses and smooth
orders. On a retained segment allocate a fixed fraction of \(T\) to each
of the finitely many independent height coordinates. Outside it, leave
reciprocal contours on an absolute or globally zero-free line. The
discarded tails are \(O(U^BT^{-N})\), with \(B\) fixed before \(N\).
Taking \(U^{\tau_-}\le T\le U^{\tau_+}\), with fixed sufficiently small
positive exponents, makes the retained bounds and discarded tails
compatible. The final choices in §B0.7 impose these requirements
simultaneously.

\hypertarget{a-common-pair-of-saturated-witnesses}{%
\subsubsection{A common pair of saturated
witnesses}\label{a-common-pair-of-saturated-witnesses}}

\textbf{Lemma 1W (shared zero witnesses).} Choose \(t_0\in[1,\allowbreak{}3/2]\)
before choosing witnesses. Each row in a nonfloor bin has one
presentation, one real power \(\sigma\in[a,\allowbreak{}a+e)\), one common height,
and dyadic exponents \(r,\allowbreak{}m\) such that

\[
r\le t_0+o(1),\qquad r+m\ge t_0-o(1),\qquad
0\le m\le\tfrac12+O(\varepsilon_0),
\tag{B0.17}
\]

\[
|M_u(U^r)|^2\gg U^{\delta r-O(\varepsilon_0)},\qquad
|S_u(U^m)|^2\gg U^{\delta m-O(\varepsilon_0)}.
\tag{B0.18}
\]

Once the presentation and dyads are fixed, the base profiles are common
to the rows, with uniform required seminorms.

\textbf{Proof.} Choose a smooth cutoff \(V_\le\), equal to one on
\([0,\allowbreak{}1]\) and zero on \([2,\allowbreak{}\infty)\), and set

\[
C_\varphi(s)=\sum_l\mu(l)\varphi(l)V_\le(q_l/U^{t_0})q_l^{-s}.
\]

At the buffered zero \(\rho\), consider

\[
\frac1{2\pi i}\int_{(2)}U^{40z}\Gamma(z)
L_{\mathbb K}^S(\rho+z,\varphi)C_\varphi(\rho+z)\,dz.
\tag{B0.19}
\]

Move to \(\Re z=-1/4\). The residue at zero vanishes. A weak polynomial
Hecke bound, ideal counting, and gamma decay make the remaining integral
\(o(1)\); the factor \(U^{-10}\) dominates these fixed powers when the
analysis height exponent is sufficiently small.

On the initial line, expand the two Dirichlet series. The unit
coefficient is \(1+o(1)\), and every coefficient at
\(1<q_{lk}\le U^{t_0}\) is zero by \(\mu*1=\mathbf1_{1}\). Removing the
unit therefore leaves a sum of magnitude \(1+o(1)\), supported smoothly
at product norm at least a fixed multiple of \(U^{t_0}\). A smooth
terminal cutoff at \(U^{41}\) costs \(o(1)\).

Split the two factors into dyads and separate their product weight by
Fourier inversion in log norm. Its common Fourier measure has bounded
mass and arbitrary polynomial decay. Some pair and frequency satisfies

\[
|M_u(U^r)S_u(U^m)|^2
\gg U^{\delta(r+m)-O(\varepsilon_0)}.
\tag{B0.20}
\]

The two factors have the same character, real power, and height.
Comparing (B0.20) with (B0.15) forces \(m\le1/2+O(\varepsilon_0)\),
since \(\delta\ge1/50\). Dividing by each of the two pointwise bounds
proves (B0.18). The inverse profile is a fixed dyadic profile times
\(V_\le(U^r y/U^{t_0})\); derivatives of the latter are supported on a
fixed annulus. The \(O((\log U)^2)\) dyadic choices and the discarded
height tails fit the previously prescribed losses. ∎

This proof permits \(t_0\) to depend on the fixed bin and on a
subsequently fixed amplitude class, provided it is fixed before that
class's witnesses are chosen. It does not select a different truncation
separately for each row inside a moment.

\hypertarget{b0.3.-the-physical-probe-signed-gram-bound-and-compensator}{%
\subsection{B0.3. The physical probe, signed Gram bound and
compensator}\label{b0.3.-the-physical-probe-signed-gram-bound-and-compensator}}

Fix primitive \(\eta\), let \(b_*\) be the product of fixed excluded
prime generators, and choose a primitive residue character
\(\xi\bmod b_*\), locally quadratic at odd primes and cubic at norm
four. Then \(\xi^6\) is principal on units. Put
\(\tau=q_{b_*}^{-1/2}g_\xi(b_*,\allowbreak{}1)\), \(\Xi(A)=\xi(A)\chi_A(b_*)\), and
\(V(r)=e^{-(\log r)^2/4}/(2\sqrt\pi)\), so \(\widehat V(t)=e^{t^2}\).
For nonnegative nonzero smooth annular \(W_0,\allowbreak{}W_1\), define \[
T_{m,a,\eta}(t)=\sum_{c\ {\rm sf},n}
\gamma_2(c)\overline{\alpha(A)G(A)}\eta(A)\Xi(A)^{-1}
R(A,a)\chi_A(m)q_c^{-1/2-t}q_n^{-1-3t},\quad A=cn^3,
\] \[
\begin{gathered}
I_\eta(X,Y,Z)=\frac1Y\sum_a
\frac{W_1(q_a/Y)\chi_a(b_*)}{\tau\xi(a)\sqrt{q_{b_*}q_aX}}\\
\times\sum_m\xi(m)q_a^{-1/2}g_{\chi_a}(a,-m)
W_0\!\left(\frac{q_m}{q_{b_*}q_aX}\right)\\
\times\frac1{2\pi i}\int_{(4)}Z^te^{t^2}T_{m,a,\eta}(t)\,dt.
\end{gathered}
\tag{B0.21}
\] Expanding fixed ray phases and separating the radial weight gives
finitely many expressions
\(Q^{-1/2}\int\widehat W_0(iv)\sum_m\Omega(q_m/Q)\xi(m)(q_m/Q)^{-iv}A_{m,\allowbreak{}\sigma,\allowbreak{}v}(Y)B_{m,\allowbreak{}\sigma}(Z)\,\allowbreak{}dv/(2\pi)\),
where \(Q=q_{b_*}XY\), \(B\) is (B0.6), and \[
\begin{gathered}
A_{m,\sigma,v}=Y^{-3/2}\sum_{a\in\sigma}
\frac{\chi_a(b_*)}{\tau\xi(a)}W_1(q_a/Y)(q_a/Y)^{-1+iv}\\ \qquad{}\times g_{\chi_a}(a,-m).
\end{gathered}
\] The \(Y^{-3/2}\) normalization is essential.

For \(P_a=Y^2/Q\ge1\), the following signed Gram estimate uses the
displayed coefficients of \(A\): \[
\sum_{q_m\ll Q}|A_{m,\sigma,v}|^2
\ll(1+|v|)^J\frac QY\left(1+P_a^{1/6}+P_a^2/Y\right)Z^\epsilon.
\tag{B0.22}
\] \textbf{Proof.} Apply Poisson summation to the physical row after
expanding the square. The complete finite transform of the two
unnormalized Gauss sums is \(q_uq_vF(u,\allowbreak{}v;j)\), where \(F\) is (B1.14).
This cancels the reciprocal modulus \(Q/(q_uq_v)\) from row Poisson and
leaves the exact prefactor \(Q/Y^3\). The zero frequency is
\(\mathbf1_{u=v}\phi(u)\), by (B1.16), and contributes \(Q/Y\).

For nonzero frequency write \(u=Cn_1,\allowbreak{}v=Cn_2,\allowbreak{}j=Ck\). Equation (B1.17)
expresses the correlation through a bounded local coefficient modulo
\(\operatorname{rad}C\). Extend that coefficient by the local table in
Appendix B1 and insert the full projector (B1.18) before separating the
columns. On its \(d\)-term, the natural zeros force \((d,\allowbreak{}Ck)=1\). After
fixed ray splitting, the remaining joint coefficient is periodic in the
two columns modulo an ideal of norm \(\ll q_Cq_k\). Its complete mean is
zero unless every good valuation of \(k\) outside \(C\) is divisible by
six: otherwise one complete nonprincipal character sum vanishes.

Dyadically let \(q_C\asymp C_0\), \(q_d\asymp D_0\),
\(q_Cq_k\asymp R P_a\), and \(N=Y/(C_0D_0)\). For a mean-zero
coefficient, Poisson in the two complex column variables has no zero
mode. Repeated integration by parts in its common smooth profile gives,
for any fixed \(A,\allowbreak{}B>0\), \[
R^{-B}C_0N^2\min\{1,(RP_a/N)^A\}.
\tag{B0.23}
\] The factor \(C_0\) is the local bound \(|L_C|\le q_C\); a nonzero
Fourier mode has size at least the reciprocal modulus. There are
\(O(D_0RP_a)\) choices of \(C,\allowbreak{}d,\allowbreak{}k\) on these dyads, with divisor factors
absorbed by \(Z^\epsilon\). The sum reduces to \[
\sum_{i,j\ge0}2^{-i-j}\min\{1,(2^{i+j}/\Lambda)^A\}
\ll_A\frac{1+\log(2+\Lambda)}\Lambda,\qquad
\Lambda=Y/(RP_a),\quad A>1.
\tag{B0.24}
\] Indeed grouping by \(i+j=n\) leaves \((n+1)2^{-n}\); splitting at
\(2^n=\Lambda\) proves the bound, with the same upper bound when
\(\Lambda<1\). After the prefactor \(Q/Y^3\), this gives
\((Q/Y)P_a^2/Y\). Taking \(B\) large sums the Fourier shells.

In the exceptional case, write \(k=e r^6\), with \(e\) supported on
\(C\) and the fixed primes and having sixth-power-free valuations. There
are \(\ll Z^\epsilon(RP_a/C_0)^{1/6}\) possible roots. The remaining
dyadic weights are \(C_0^{-1/6+\epsilon}D_0^{-1+\epsilon}\); their sums
give \((Q/Y)P_a^{1/6}\). This proves (B0.22), with the height factor
coming from finitely many profile derivatives. The exceptional term is
retained, and the proof does not give the same bound for arbitrary
column coefficients. ∎

Take disjoint prime slots of lengths \(P_i=Z^{\ell_i}\),
\(\ell=\sum\ell_i\), with nonnegative profiles \(w_i\). If
\(I_{\eta;D}\) inserts \(D\mid cn^3\), the exact finite compensator is
\[
I_\eta^{\rm comp}(Z)=\sum_{(p_i)}\prod_iw_i(q_{p_i}/P_i)
\sum_J(-1)^{|J|}q_{p_J}^{-3/2}\overline{\eta(p_{J^c})}
I_{\eta;p_{J^c}}(X/q_{p_J},Y/q_{p_J},Zq_{p_{J^c}}).
\tag{B0.25}
\] All good primes are allowed; \(\overline\eta\), rather than \(\eta\),
is necessary.

The reflected-energy corollary in Appendix B2 says that, for completed
length \(Z^N\), row range \(Z^M\), and product-form prime marks of total
length \(z\), \[
M+2z\le N\quad\Longrightarrow\quad
\sum_{q_m\ll Z^M}|\mathcal T(m)|^2\ll Z^{M+(N+2z-2M)_+/4+\epsilon}.
\tag{B0.26}
\] The proof uses the full reflected-block exponent, so it also covers
collisions between a row and a mark.

Consequently, for \(X=Z^x,\allowbreak{}Y=Z^y,\allowbreak{}M=x+y\), \[
0<\ell<x\le y,\quad M+\ell\le1,\quad y-2x+\ell<0
\quad\Longrightarrow\quad
|I_\eta^{\rm comp}|\ll Z^{x/2+(y-x)/12+(1+3\ell-2M)_+/8+\epsilon}.
\tag{B0.27}
\] A rescaled subset of total length \(d\) has
\(M'=M-2d,\allowbreak{}z'=\ell-d,\allowbreak{}N'=1+\ell-d\). Combining (B0.22)--(B0.26) with its
outside coefficient and tuple count adds \(-d+(1+3\ell-2M+d)_+/8\),
decreasing in \(d\). Its maximum is at zero. Requiring every rescaled
energy separately to have exponent \(M'\) would lose this saving.

\hypertarget{b0.4.-exact-euler-replacement-and-principal-signal}{%
\subsection{B0.4. Exact Euler replacement and principal
signal}\label{b0.4.-exact-euler-replacement-and-principal-signal}}

Write \(\mathscr M(z)=\int_0^\infty\widetilde W_0(r)r^{z-1}dr\) for the
radial Fourier Mellin transform. Poisson in the physical row and
\(s=t+1-z\) give \[
\begin{gathered}
I_\eta=\frac1{(2\pi i)^3}\int_{(2)}\int_{(3)}\int_{(3)}\\
\times X^{1/2-z}Z^{s+z-1}Y^{w-1}e^{(s+z-1)^2}\mathscr M(z)\widehat W_1(w)\\
\times\sum_u^{(6)}q_u^{-z}\overline{\xi(u)}
\frac{\zeta_{\mathbb K}^S(6z)L_{\mathbb K}^S(w,\psi_u)}{L_{\mathbb K}^S(s,\eta\overline\psi_u)}H_{\eta,u}\,dz\,dw\,ds.
\end{gathered}
\tag{B0.28}
\] The fixed primitive factor excludes zero frequency. Writing a nonzero
frequency as \(ua^6\) gives \(\zeta_{\mathbb K}^S(6z)\); no unit row is
discarded.

The complete local table is as follows. For \(p\notin S\), set \[
Q=q_p,\ j=v_p(u),\ \rho=\chi_p(u/p^j),\
\mathfrak a=\overline{\alpha(p)}^3\eta(p)^3,\
\mathfrak b=\overline{\alpha(p)}^2\eta(p)^2\overline{\chi_p(4)},\quad\mathfrak b^3=\mathfrak a^2,
\] \[
V=Q^{-6z},\quad R=\mathfrak a^2Q^{4-6s-6z},\quad W^\circ=\rho Q^{-w},
\quad W=\chi_p(u)Q^{-w},\quad D=\eta(p)\overline{\chi_p(u)}Q^{-s}.
\] In particular \(W=D=0\) if \(j>0\). Then \[
P_p^*=\frac1{1-R}\left\{\frac{R(1-Q^{-1})-\eta(p)(Q-1)Q^{-s-w}V^{\mathbf1_{j\le1}}}{1-V}+J_j\right\},
\] \[
P_p=(1-V)^{-1}+\mathbf1_{j=0}W^\circ/(1-W^\circ)+P_p^*,\qquad
H_p=P_p(1-V)(1-W)/(1-D),\quad H_{\eta,u}=\prod_{p\notin S}H_p.
\tag{B0.29}
\]

\begin{longtable}[]{@{}
  >{\raggedleft\arraybackslash}p{(\columnwidth - 2\tabcolsep) * \real{0.5714}}
  >{\raggedright\arraybackslash}p{(\columnwidth - 2\tabcolsep) * \real{0.4286}}@{}}
\toprule\noalign{}
\begin{minipage}[b]{\linewidth}\raggedleft
\(j\)
\end{minipage} & \begin{minipage}[b]{\linewidth}\raggedright
\(J_j\)
\end{minipage} \\
\midrule\noalign{}
\endhead
\bottomrule\noalign{}
\endlastfoot
0 & \(-\eta(p)\rho^{-1}Q^{-s}+W^\circ R\) \\
1 & \(\eta(p)Q^{-s-w}\) \\
2 & \(\mathfrak a\rho^{-3}Q^{3/2-3s}\) \\
3 &
\(-\eta(p)\mathfrak b\rho^{-2}Q^{2-3s-w}+\mathfrak b^2\rho^{-4}Q^{2-4s}\) \\
4 & \(-\eta(p)\mathfrak a\rho^{-3}Q^{5/2-4s-w}\) \\
5 & \(-\mathfrak a^2Q^{3-6s}\) \\
\end{longtable}

The exponent \(V^{\mathbf1_{j\le1}}\) is \(V\) at \(j=0,\allowbreak{}1\), and one
otherwise. Appendix B4 derives this table by full-modulus Poisson
summation and CRT. For reference, its finite local sum is \[
C_p(t,k,j')=\sum_{\substack{d\bmod p^k\\p^{j'}\equiv p^td\pmod{p^k}}}
\chi_p(d)^k g_{\chi_p^t}\!\left(p^t,\frac{p^{j'}-p^td}{p^k}\right),\qquad t=e+3l,
\tag{B0.30}
\] with \(e=0,\allowbreak{}1\). The appendix evaluates every case, restores the unit
phases, and sums the resulting geometric series. The local table
therefore includes all shared primes and imprimitive rows.

Set \(c_p=\overline{\eta(p)}Q^s\), \(h_p=Q^{-w}\), and define, without
division by \(P_p\) or \(H_p\), \[
G_p=\frac{(c_p-h_p)(1-V)(1-W)P_p^*-h_p(1-VW)}{1-D}.
\tag{B0.31}
\] Then
\(\overline{\eta(p)}Q^{s+z-1}P_p^*-Q^{z-w-1}P_p =Q^{z-1}(1-D)G_p/((1-V)(1-W))\).
Thus compensation replaces \(H_{\eta,\allowbreak{}u}\) in (B0.28) by \[
\mathcal H_{\eta,u,Z}=\sum_{(p_i)}\prod_iw_i(Q_i/P_i)Q_i^{z-1}G_{p_i}
\prod_{p\notin S\cup\{p_i\}}H_p.
\tag{B0.32}
\] For \(p\nmid u\), put \(v=\chi_p(u)\), \(E_p=P_p^*+D\). The phase
identity \(c_pD=v^{-1}\) gives \[
G_p+v^{-1}H_p=\frac{(1-V)(1-W)}{1-D}
\left[(v^{-1}-h_p)\left(\frac V{1-V}-D\right)
+(c_p+v^{-1}-h_p)E_p\right].
\tag{B0.33}
\] Define \(\Delta_p=G_p+\overline{\psi_u(p)}H_p\) also at row primes.
Expanding \(G_p=-\overline{\psi_u(p)}H_p+\Delta_p\) factors every main
slot into \(-\sum_pw_i(Q/P_i)Q^{z-1}\overline{\psi_u(p)}\), while error
slots replace the corresponding \(H_p\) by \(\Delta_p\). This exact
expansion keeps row collisions and never assumes local nonvanishing.

The local products converge normally and satisfy
\(H_{\eta,\allowbreak{}u}\ll q_u^\epsilon\) on \[
\Re s\ge.51,\quad\Re z\ge.34,\quad\Re w\ge-.01,\quad\Re(s+w)\ge1+\epsilon_1.
\tag{B0.34}
\] For the principal row they also converge on
\(\Re s\ge.87,\allowbreak{}\Re w\ge3/4,\allowbreak{}\Re z\ge1/8\). The entire correction has a
fixed raw polynomial bound on every contour box. On
\((\beta_*+e,\allowbreak{}1/2,\allowbreak{}.34)\) and \((2,\allowbreak{}2,\allowbreak{}v)\), respectively, it is
\(\ll q_u^\epsilon Z^{.34\ell}\) and \(\ll q_u^\epsilon Z^{v\ell}\).

Only \(u=1\) has principal numerator. A nonunit sixth-power-free row
coprime to \(S\) has nontrivial tame inertia. For a unit \(u\ne1\),
choose a primary \(n\equiv-2-3\omega\pmod{36\mathcal O}\), coprime to
the remaining fixed primes by CRT. Then \(q_n\equiv7\pmod{36}\), so
\(\chi_n(u)=u\ne1\); this proves that its numerator character is also
nonprincipal. At \(w=1,\allowbreak{}z=1/6\), write \(t=Q^{-1}\), \(D=\eta(p)Q^{-s}\),
\(R=\overline{\alpha(p)}^6\eta(p)^6Q^{3-6s}\). Then \[
H_{\eta,p}(s)=\frac{1-t^2}{1-R}\left(1+\frac{t(D-R)}{1-D}\right).
\tag{B0.35}
\] Enlarging the fixed cutoff gives \(|H_\eta(s)-1|\le1/2\) for
\(\Re s\ge.87\). At this residue, \(G_p/H_p=-1+O(Q^{-.87})\). With
\(r_{\rm slot}\) fixed prime slots, define \[
c_S=\frac{(\operatorname{Res}_{s=1}\zeta_{\mathbb K}^S(s))^2}{6}\widehat W_1(1)\mathscr M(1/6)>0,
\quad A(Z)=(-1)^{r_{\rm slot}}Z^{-\ell/6}\prod_i\sum_pw_i(Q/P_i)Q^{-5/6}.
\] Writing the radial Riesz kernel as a positive integral of Gaussians
and Fourier-transforming gives, with \(\kappa_{\rm rad}=2\pi/\sqrt3\),
\[
\mathscr M(\sigma)=\kappa_{\rm rad}^{\,1-2\sigma}
\frac{\Gamma(\sigma)}{\Gamma(1-\sigma)}
\int_0^\infty W_0(r)r^{-\sigma}\,dr>0,\qquad 0<\sigma<1.
\] In particular, \(\mathscr M(1/6)>0\). The prime ideal theorem gives
\(|A(Z)|\asymp(\log Z)^{-r_{\rm slot}}\), and
\(\mathcal H_{\eta,\allowbreak{}1,\allowbreak{}Z}(s,\allowbreak{}1,\allowbreak{}1/6)=H_\eta(s)Z^{\ell/6}A(Z)(1+O(Z^{-\mu_P}))\)
for any \(0<\mu_P<.87\min\ell_i\). Put \[
h=1+\ell-x,\quad k=5/6-x/3-\ell/6,\quad J_\eta=I_\eta^{\rm comp}/(c_SA(Z)),
\] \[
f_\eta(Z)=\frac1{2\pi i}\int_{(2)}Z^{s-k}e^{(s-5/6)^2}\frac{H_\eta(s)}{L_{\mathbb K}^S(s,\eta)}\,ds.
\tag{B0.36}
\] The main term of the double residue is \(f_\eta\), with the
prime-product approximation error estimated in §B0.7; the power
\(Z^{\ell/6}\), rather than the logarithmic normalizer, is the
compensator's quantitative gain.

\hypertarget{b0.5.-amplitude-sensitive-witness-and-numerator-estimates}{%
\subsection{B0.5. Amplitude-sensitive witness and numerator
estimates}\label{b0.5.-amplitude-sensitive-witness-and-numerator-estimates}}

For the contradiction argument, suppose \(\beta_*>b_0\). For a bin
\(q_u\asymp U=Z^d\), let \(\kappa=2\beta_*-1\in[7/10,\allowbreak{}3/4]\),
\(L=\ell/d\). Partition physical main-slot amplitudes into finite
classes with exponents \(g_i\in[0,\allowbreak{}\delta/2]\); error slots and small
amplitudes have \(g_i=0\). Put \(q=L^{-1}\sum_i(\ell_i/d)g_i=\delta p\),
\(0\le p\le1/2\). Greedy filling in decreasing \(g_i\) supplies, at
capacity \(z\le L\), a genuine subcollection with \[
\prod_{i\in I}|Q_{u,i}|^2\ge U^{2qz-\delta\max_i(\ell_i/d)-o(1)}.
\tag{B0.37}
\] Discarding one fractional slot accounts for the mesh loss. Adjusting
its coefficient by the inverse fixed character, or conjugating the whole
witness, matches the physical \(\overline{\psi_u}\) orientation without
introducing row-dependent coefficients.

For \(r\le1\), use inverse capacity \((1-r)/2\); for \(m\le1/2\), use
plain capacity \((1-2m)/(6\kappa)\). MI, MP and (B0.18) give count lines
\[
A_I(r)=1-q-(\delta-q)r,\qquad
A_P(m)=1-q/(3\kappa)-2(\delta-q/(3\kappa))m.
\tag{B0.38}
\] \textbf{Proof of the row count.} The short-witness crossing is
\(r_*(t)=[2(1-p/(3\kappa))t+p/(3\kappa)-p]/[3-p-2p/(3\kappa)]\), with
\(r_*\ge13/21\), \(1/3\le t-r_*\le1/2\). Thus inverse and plain
capacities are at most \(4/21\) and \(1/12\). If \(L>4/21\), the supply
suffices. On the inverse branch, decrease capacity to
\(z=(1-r)/2-\nu_0\). Then \[
1-r-2z=2\nu_0,\qquad 3-2r-8z=8\nu_0+(2r-1)\ge8\nu_0+5/21.
\tag{B0.39}
\] Thus both MI inequalities have positive margins. Smaller capacities
use the unmarked estimate with an \(O(\nu_0)\) loss; decreasing the
plain capacity is handled in the same way. Set \[
C_\kappa(p)=\frac{2(1-p)(1-p/(3\kappa))}{3-p-2p/(3\kappa)},\quad V=5/6-\delta.
\] The short envelope is \(1-\delta+\delta C_\kappa(p)(3/2-t)\); the
long inverse envelope from (B0.9) is \(1-\delta+V(t-1)\). Balancing at
\(t=1+\delta C_\kappa(p)/(2(V+\delta C_\kappa(p)))\) gives \[
R_\kappa(\delta,p)=1-\delta+
\frac{\delta C_\kappa(p)V}{2(V+\delta C_\kappa(p))},\qquad
\#\text{rows}\ll U^{R_\kappa+\epsilon}(1+T)^A.
\tag{B0.40}
\] The endpoint \(V=0\) uses \(t=3/2\). Without marks, the same balance
gives \(R_0(\delta)=1-10\delta/[3(5-2\delta)]\), \(t=5/(5-2\delta)\).

The physical numerator needs its \textbf{own} moment. For one main/error
term of (B0.32), define \[
\mathcal N_u(Y;s,z)=\frac1{2\pi i}\int_{(3)}Y^{w-1}\widehat W_1(w)L_{\mathbb K}^S(w,\psi_u)\mathcal H^{\rm term}_{\eta,u,Z}(s,w,z)\,dw.
\] At \(\Re s=a+16e,\allowbreak{}\Re z=z_0=.34\), with \(m=\log Y/\log U>1/2\), set
\(z_A=[2\min(m,\allowbreak{}1)-1]/(6\kappa)\le L\). Then \[
\sum_{u\text{ in class}}|\mathcal N_u|
\ll Y^{-1/2}Z^{\ell(z_0-1/2)+\ell q}
U^{1/4+3R/4-qz_A/2+\epsilon}(1+T)^A.
\tag{B0.41}
\] The reciprocal denominator is independent of \(w\) and is estimated
\textbf{after} inversion.

\textbf{Proof of the numerator estimate.} Appendix B4 proves the
following coefficient bounds directly from the six-row table. The
factors are polynomials in \(Q^{-w}\): \(\deg H_p\le2\),
\(\deg\Delta_p\le3\), with degrees one smaller at row primes. Their
denominators are independent of \(w\). Off the row, direct expansion of
(B0.29) gives \[
|H_{p,0}-1|+\sum_{j=1}^2|H_{p,j}|Q^{j(a-1)}\ll Q^{-1-\epsilon_e},\quad
\sum_{j=0}^3|\Delta_{p,j}|Q^{j(a-1)}\ll Q^{-a}.
\tag{B0.42}
\] At a row prime of valuation \(j\), allocate \(\min(k,\allowbreak{}j-1)\) powers of
a monomial shift \(Q^k\) to
\(h_r\mid\Lambda_u=q_u/q_{\operatorname{rad}(u)}\); the rest belongs to
a free shift \(h_0\). The exact table yields \[
\sum_k|H_{p,k}|Q^{-k/2}Q^{\delta(k-j+1)_+/2}\le1+CQ^{-c_e},\quad
\sum_k|\Delta_{p,k}|Q^{-(k+1)/2}Q^{\delta(k-j+1)_+/2}\ll1.
\tag{B0.43}
\] The sharp error at \(j\ge2\) is \(\Delta_{p,\allowbreak{}1}\asymp Q\): central
normalization cancels it exactly, and its shift is absorbed by
\(\Lambda_u\). The quadratic error at \(j=2\) contributes
\(Q^{1-a}Q^{-3/2}Q^{\delta/2}=Q^{-1}\). No valuation is charged twice.

On the original annular support, \(W_1(q_nh/Y)\ne0\) forces
\(h\le C Y\). Retain this cutoff in the coefficient expansion: it gives
\(\log_Uh_0\le m+O(1/\log U)\), so all later moment lengths stay in
fixed bounded ranges. Inverting a monomial with shift \(h=h_0h_r\) gives
\(Y^{-1/2}h^{-1/2}S_{\psi_u}(Y/h)\). Primitive functional-equation
reflection preserves central amplitude and gives
\(\varepsilon_uS_{\overline\psi_u}(\mathfrak c_uh/Y;W_1^\vee)\), where
\(\mathfrak c_u=3q_{\mathfrak f_u}/(2\pi)^2\asymp_Sq_{\operatorname{rad}(u)}\)
and \(\widehat W_1^\vee(v)=\widehat W_1(1-v)\Gamma(v)/\Gamma(1-v)\).
Thus the dual scale is \(\ll_S(U/Y)h_0\). Before reflection, delete the
fixed Euler factors by their finite expansion; restore their natural
zeros afterward. The restoring series has mass bounded by a fixed
product of \((1-Q^{-1/2})^{-1}\) and only shortens scales. The reflected
profiles are common and rapidly decreasing at infinity. Moving their
Mellin line just left of zero crosses only the simple \(\Gamma(v)\) pole
there, giving \(W_1^\vee(x)=O(1)\) as \(x\downarrow0\). Hence annular
decomposition has summable normalized coefficients. The varying
conductor changes their scale, not the arithmetic coefficients to which
MP is applied. Equations (B0.42)--(B0.43) give the rowwise coefficient
mass below. Indeed, after its slot normalizer is removed, an unramified
error slot has weighted mass
\(P^{-1/2}\sum_{q_p\asymp P}q_p^{-a}\ll P^{1/2-a}\le1\); ramified
choices cost a divisor power. The disjoint slots prevent double
counting. Hence \[
\sum_{h_0,h_r}|c_u(h_0,h_r)|h_0^{\delta/2}\ll U^\epsilon.
\tag{B0.44}
\] Replace remaining plain values by a common-profile scale supremum
\textbf{before} applying MP. For \(h_0\asymp U^t\), the allowed capacity
is \(z(t)=\max\{0,\allowbreak{}\min(1,\allowbreak{}2m-1-2t)/(6\kappa)\}\). The chosen prime
product is at least \(U^{qz(t)-o(1)}\) throughout its amplitude class.
Hölder and MP therefore give \(U^{1/4+3R/4-qz(t)/2+\epsilon}\); the
fourth root accounts for the gain \(qz(t)/2\). Equation (B0.44) pays
\(U^{-\delta t/2}\). Positive capacity uses MP and zero capacity uses
MP0. Reserve small capacity decrements for the reflected-scale
enlargement and the two finite meshes before taking the limit. Since
\(-\delta t/2-qz(t)/2\le-qz_A/2\), summing these dyads proves (B0.41).
Subunit reflected scales are handled by the original Schwartz profiles
and assigned nonnegative length zero; capacities are capped at
\(1/(6\kappa)\), never computed from a negative length.

\hypertarget{b0.6.-geometry-and-an-exact-exponent-certificate}{%
\subsection{B0.6. Geometry and an exact exponent
certificate}\label{b0.6.-geometry-and-an-exact-exponent-certificate}}

Let \[
\kappa_0=(5-3\ell_0)/6,\qquad
\delta_0=\frac{25-60\ell_0+27\ell_0^2}{6(5+4\ell_0-3\ell_0^2)}.
\] With \(C=(6\kappa_0-1)/(15\kappa_0-2)\), \(V=5/6-\delta_0\), put \[
R'=-1+\frac{C(V^2-\delta_0^2C)}{2(V+\delta_0C)^2},\quad
h_0=-\frac{(1+\ell_0)/2+\ell_0/(6\kappa_0)}{3R'/4-1/(24\kappa_0)},
\quad x_0=1+\ell_0-h_0,\quad y_0=h_0-2\ell_0.
\tag{B0.45}
\] Numerically
\[
\begin{gathered}
(\delta_0,h_0,x_0,y_0)=\\
(.469344858698143,.793947597388136,.373172481118674,.459707440374516).
\end{gathered}
\]
Then \(M=1-\ell_0\), \(\ell_0<x_0<y_0<2x_0\), \(y_0<1/2<h_0<2y_0\), and
(B0.27)'s physical boundary is \(b_0=11/12-\ell_0/4\).

For general nearby geometry with \(h=1+\ell-x\), the dyadic numerator
boundary is \[
F_\kappa(d)=a+h(z_0-1/6)-y/2-\ell/2+\ell q
+d(1/4+3R_\kappa/4-z_0)-\frac q{12\kappa}\min(d,2y-d).
\tag{B0.46}
\] It is convex in \(d\); checking endpoints suffices. At \(d=h_0\), \[
F^{(4)}_\kappa(\delta,p)=\frac{1+\ell_0}{2}+\frac\delta2
+h_0\left(\frac{3R_\kappa(\delta,p)}4-\frac5{12}\right)
+\delta p\left(\ell_0-\frac{h_0-4\ell_0}{12\kappa}\right).
\tag{B0.47}
\] The binding point is \((\delta,\allowbreak{}p,\allowbreak{}\kappa)=(\delta_0,\allowbreak{}1/2,\allowbreak{}\kappa_0)\).
The equations behind the quartic are \[
(324\kappa^2-3)\delta^2+(-756\kappa^2+54\kappa+5)\delta+300\kappa^2-40\kappa=0,
\quad \delta=\frac{\kappa(5-9\ell)}{6\kappa+(6\kappa+2)\ell}.
\tag{B0.48}
\] The first is equivalent to
\(R_\kappa(\delta,\allowbreak{}1/2)=5/9+\delta/(18\kappa)\), eliminating the \(h\)
coefficient; (B0.45) enforces stationarity in \(\delta\).

For an exact global certificate, suppress subscripts and let
\(u=1/2-p\). Set \[
\begin{gathered}
j_{00}=300\kappa-40,\ j_{01}=120\kappa+80,\ j_{10}=24-216\kappa,\ j_{11}=144\kappa-96,\\
B_0=24\kappa-4,\ B_1=48\kappa,\ B_2=16,\quad
J=j_{00}+j_{01}u+\delta(j_{10}+j_{11}u+96u^2),\\
T=\ell-(h-4\ell)/(12\kappa),\quad
k_0=b-(1+\ell)/2-h/3,\\ k_1=-1/2+3h/4-T/2,\quad
K_*=8j_{10}k_1+18hB_0.
\end{gathered}
\] Then \(69<J\le270\), and \[
8J\{b-F^{(4)}_\kappa(\delta,1/2-u)\}
=K_*(\delta-\delta_0)^2+uA_1(\delta)+u^2A_2(\delta)+768Tu^3\delta^2,
\tag{B0.49}
\] where \[
\begin{aligned}
A_1={}&8j_{01}k_0+[8(j_{00}T+j_{01}k_1+j_{11}k_0)-15hB_1]\delta
+[8(j_{10}T+j_{11}k_1)+18hB_1]\delta^2,\\
A_2={}&[8(j_{01}T+96k_0)-15hB_2]\delta
+[8(j_{11}T+96k_1)+18hB_2]\delta^2.
\end{aligned}
\] This follows by substituting
\(R=1-\delta+\delta(B_0+B_1u+B_2u^2)(5-6\delta)/(2J)\). The zero-\(u\)
coefficient has a double zero by (B0.48) and stationarity. Rational
interval substitution from (B0.1) gives \[
A_1>36-174\delta+347\delta^2>0,\quad
A_2>38\delta+257\delta^2,\quad179<K_*<180,\quad117<768T<119.
\] Thus \(F^{(4)}_{\kappa_0}\le b_0\), with equality only at the binding
point, without relying on a floating-point grid.

All capacities fit: \(\ell_0/h_0>1/5>4/21\), and the physical numerator
requires only \(dz_A=\min(d,\allowbreak{}2y_0-d)/(6\kappa)\le\ell_0\), since
\(y_0<4\ell_0\). Also \(\partial_\kappa C\le1/18\),
\(\partial_C R\le5/12\), hence \(\partial_\kappa F^{(4)}<1/10\). If
\(\beta_*>b_0\), choose \(\kappa=2\beta_*-1\); the MP hypothesis holds
by equality and \(F^{(4)}_\kappa-\beta_*\le-4(\beta_*-b_0)/5\). The
other endpoint in (B0.46) is at most
\(11/12-49h_0/150+11\ell_0/12<b_0-.06\). The floor
\((a,\allowbreak{}R,\allowbreak{}q)=(.51,\allowbreak{}1,\allowbreak{}q\le.01)\) has endpoint at most
\(.51+h_0/3+.51\ell_0<b_0-.01\).

\hypertarget{b0.7.-completing-the-zero-free-argument}{%
\subsection{B0.7. Completing the zero-free
argument}\label{b0.7.-completing-the-zero-free-argument}}

Retain the assumption \(\beta_*>b_0\), and put \(\mu=\beta_*-b_0>0\).
The geometry is (B0.45), and \(k\) is defined in (B0.36). All real power
margins in this section are common to the family; constants, the fixed
exclusion \(S\), and sufficiently small analysis-height exponents may
depend on the fixed target \(\eta\).

\hypertarget{the-principal-term}{%
\subsubsection{The principal term}\label{the-principal-term}}

(B0.27), the normalization in (B0.36), and (B0.45) give

\[
J_\eta(Z)\ll Z^{b_0-k+\varepsilon}.
\tag{B0.50}
\]

We will prove, for some fixed \(\Delta>0\) independent of \(\eta\),

\[
J_\eta(Z)-f_\eta(Z)\ll_\eta Z^{\beta_*-k-\Delta}.
\tag{B0.51}
\]

Consider first \(u=1\) in the compensated version of (B0.28). Move to
\((\Re s,\allowbreak{}\Re w,\allowbreak{}\Re z)=(\beta_*+e,\allowbreak{}1+e,\allowbreak{}1/6+e)\). The global line for \(s\)
is zero-free by definition and (B0.14). Move \(w\) to \(3/4\), crossing
only the simple numerator pole at one. In its residue, move \(z\) to
\(1/8\), crossing only the simple pole at \(1/6\). The local product and
replacement are holomorphic on the full principal box following (B0.34).
The Mellin transform \(\mathscr M\) has no pole in these positive
strips.

Equations (B0.35)--(B0.36) identify the double residue, including both
residue scalars. Substitute the prime-product approximation following
(B0.35) on the global \(s\)-line before replacing its main integral by
the line \(\Re s=2\). The latter move crosses no singularity. Its main
term is exactly \(f_\eta\). The two contour remainders and the error in
the residue approximation are respectively

\[
O_\eta\!\left(
 Z^{\beta_*-k-y/4+O(e)+\varepsilon}
+Z^{\beta_*-k-h/24+O(e)+\varepsilon}
+Z^{\beta_*-k-\mu_P+O(e)+\varepsilon}
\right).
\tag{B0.52}
\]

For the first two bounds use the quotient-free correction on the
principal box, with size \(O(Z^{\ell\Re z+\varepsilon})\); the remaining
geometric exponents are \(y(\Re w-1)\) and \(h(\Re z-1/6)\). The third
bound uses the prime-product approximation following (B0.35) only on its
stated residue locus. In particular a \(Q^{-7/8}\) approximation has not
been used below its permissible line.

\hypertarget{retained-nonprincipal-rows}{%
\subsubsection{Retained nonprincipal
rows}\label{retained-nonprincipal-rows}}

Fix \(d_{\min}=1/200\) and a sufficiently small positive extension
\(\zeta_0\). For \(U=Z^d\), with \(d_{\min}\le d\le h+\zeta_0\),
partition nonprincipal rows into the buffered bins of §B0.2.

For each bin move \(z\) to \(z_0=17/50\), keeping \(s=2,\allowbreak{}w=3\). Next move
\(s\) to \(\beta_*+20e\), still keeping \(w=3\). On retained heights
move its segment to \(a+16e\). The denominator's zero-free buffer
permits this last move for every presentation in the bin; outside those
heights leave the reciprocal on the global line. The entire deformation
lies in (B0.34). There is no nonprincipal numerator pole, and no
\(z\)-pole for \(\Re z\ge z_0\).

Use the exact expansion (B0.33), fix a main-prime amplitude class, and
now perform the full \(w\)-inversion on its original absolute line. The
reciprocal is independent of \(w\). (B0.40) and (B0.41), and the
buffered subpower reciprocal bound, give

\[
\sum_{u\ \text{ in the class}}^{(6)}
\left|\frac{q_u^{-z}\mathcal N_u(Y;s,z)}
{L_{\mathbb K}^S(s,\eta\overline{\psi_u})}\right|
\ll (1+T)^A Y^{-1/2}
Z^{\ell(z_0-1/2)+\ell q+\varepsilon_2}
U^{1/4+3R_\kappa/4-z_0-qz_A(d)/2+\varepsilon_2},
\tag{B0.53}
\]

where

\[
z_A(d)=\frac{2\min(y/d,1)-1}{6\kappa},\qquad
d z_A(d)=\frac{\min(d,2y-d)}{6\kappa}.
\tag{B0.54}
\]

For sufficiently small \(\zeta_0\), this capacity is positive and
available. The independent heights of the physical prime sums and the
detector witnesses are covered by the moment seminorm bounds; no common
physical/witness height is asserted.

Multiplying (B0.53) by the remaining factors in (B0.28), and subtracting
\(\beta_*-k\), gives exactly

\[
a-\beta_*+h(z_0-1/6)-y/2-\ell/2+\ell q
+d(1/4+3R_\kappa/4-z_0)
-\frac q{12\kappa}\min(d,2y-d).
\tag{B0.55}
\]

At \(d=h\) this is \(F^{(4)}_\kappa(\delta,\allowbreak{}p)-\beta_*\), with
\(F^{(4)}\) as in (B0.47). Its convexity in \(d\), the endpoint bounds
in §B0.6, and the floor estimate therefore supply a fixed negative
margin on the whole retained interval. The extension above \(h\) costs
at most a fixed multiple of \(\zeta_0\). Dyadic sums, amplitude classes,
and finite profile-height costs fit inside a reserved fraction of this
margin.

All witness capacities are available because \(\ell/(h+\zeta_0)>4/21\).
The new numerator capacity is available because
\(\min(d,\allowbreak{}2y-d)\le y<4\ell\le6\kappa\ell\). If \(\ell_{\max}\) is the
largest base-\(Z\) slot length, its largest relative length on the
retained range is at most \(\ell_{\max}/d_{\min}\). Finitely many
sufficiently short equal slots, on disjoint norm subintervals, therefore
meet both the MP mesh and the greedy rounding condition uniformly,
including the smallest dyads.

\hypertarget{the-remaining-rows-and-heights}{%
\subsubsection{The remaining rows and
heights}\label{the-remaining-rows-and-heights}}

The remaining units and nonunit rows \(q_u\le Z^{d_{\min}}\) do not need
a detector. Use the global lines \[
(\Re s,\Re w,\Re z)=(\beta_*+e,1/2,z_0).
\] For their nonprincipal numerators the ordinary functional equation
and convexity give the more than sufficient bound \[
L_{\mathbb K}^S(1/2+it,\psi_u)\ll U^{3/5+\varepsilon}(1+|t|)^A.
\] Together with the first of the raw bounds following (B0.34), ideal
counting, and \(q_u^{-z_0}\), this makes their exponent relative to
\(\beta_*-k\) at most

\[
h(z_0-1/6)-y/2+\frac{63}{50}d_{\min}+O(e)+\varepsilon.
\tag{B0.56}
\]

For (B0.45) the expression before losses is less than \(-1/20\), as
follows from the quartic isolating interval (B0.1). The finite unit rows
are included by replacing \(U\) with a fixed lower bound. No
nonprincipal pole is crossed.

For \(q_u>Z^{h+\zeta_0}\), take \((\Re s,\allowbreak{}\Re w,\allowbreak{}\Re z)=(2,\allowbreak{}2,\allowbreak{}v)\), with
fixed \(v>2\). The second of the raw bounds following (B0.34), the
absolute \(L\)-bounds, and the row count give for a dyad

\[
O\!\left(Z^{B_0+hv+\varepsilon}U^{1+\varepsilon-v}\right),
\qquad B_0=x/2+1+y.
\]

The dyadic sum is

\[
O\!\left(
 Z^{B_0+(h+\zeta_0)(1+\varepsilon)-\zeta_0v+\varepsilon}
\right).
\tag{B0.57}
\]

Choose \(v\) after \(\zeta_0\), as a sufficiently large fixed number.
This term then has any required fixed negative power relative to the
principal scale. It does not use the central factorization or a
large-row zero count.

We spell out the height control because it is needed for a common-family
argument. The three independent frequency coordinates can be taken to be
\[
\Im(s+z),\qquad\Im z,\qquad\Im w.
\] The Gaussian decays in the first, while the two Mellin transforms
have arbitrary polynomial decay in the latter two on every fixed
positive real strip used above. A fixed fraction of the bin's height
buffer is allocated to each coordinate. The joining segments leave every
reciprocal either within its retained zero-free buffer or on a global
zero-free line. The unseparated correction has raw bound
\(Z^B U^\varepsilon\) on all these fixed real boxes, by (B0.31); \(B\)
is fixed before the integration-by-parts order. Off the retained region
the elementary conductor-height bounds and this raw estimate thus give

\[
O_N(Z^{B'}T^{-N})
\tag{B0.58}
\]

with \(B'\) fixed before \(N\). For the retained numerator estimate the
\(w\)-integral was performed exactly on \(\Re w=3\); its common
reflected-profile tails were handled before MP in the proof of (B0.41).
The remaining deformations use only the buffered \(s\) coordinate and
the external \(z\) weight. The same bound also applies to the tails in
Lemma 1W, after the finite smooth orders required by the moments have
been fixed. In particular rapid decay is applied before any
tail-dependent separation that could alter \(B'\).

Here is a noncircular order of choices. First fix the contradiction gap
\(\mu\), a fraction of the high margin, and a small row extension
\(\zeta_0\). Fix the moment loss, capacity decrement, and amplitude
width in terms of this margin. Obtain the uniform MP mesh, choose a
finite slot count satisfying that mesh on \(d\ge d_{\min}\), and freeze
the slot lengths. This fixes a positive \(\mu_P\) in the prime-product
approximation following (B0.35). Now choose the real contour and
detector buffers smaller than the high margin and the principal margins
\(y/4,\allowbreak{}h/24,\allowbreak{}\mu_P\). The moment loss was already fixed; shrinking these
buffers does not demand a new mesh. Losses used just for the principal
product or the physical estimate can be chosen independently of the
marked-moment loss.

Next fix the target, its finite character group and cutoff, and all
finite internal Sobolev and moment orders. Choose a positive,
sufficiently small \(\tau\), and take \(T=3Z^\tau\). It is small enough
both in terms of \(d_{\min}\) and every positive \(\ell_i\) to meet all
retained-height restrictions and to absorb \((1+T)^A\) into the reserved
margin. It is also bounded below by a fixed positive power of \(U\) on
the retained row interval. Finally take the external tail order \(N\)
sufficiently large in (B0.58). Enlarging this last order affects only
the external Schwartz weights. No slot count, real margin, or raw power
is selected after the asymptotic limit.

All exponent gaps before the final choice of \(\tau\) are common to the
family. A target-dependent \(\tau\) can absorb its target-dependent
polynomial height costs into the same reserved gap. Taking a fixed
sufficiently small \(\Delta>0\), less than fixed fractions of
\(\mu,\allowbreak{}y/4,\allowbreak{}h/24,\allowbreak{}\mu_P\) and the gap in (B0.56), proves (B0.51).

\hypertarget{mellin-continuation}{%
\subsubsection{Mellin continuation}\label{mellin-continuation}}

Choose the \(\varepsilon\) in (B0.50) smaller than \(\mu/4\), and reduce
\(\Delta\), if necessary, so that \(\Delta<\mu/4\). Equations
(B0.50)--(B0.51) imply

\[
f_\eta(Z)\ll_\eta Z^{\beta_*-k-\Delta}\qquad(Z\longrightarrow\infty).
\tag{B0.59}
\]

For \(0<Z\le1\), move the defining line in (B0.36) to any fixed \(A>2\).
Absolute Euler convergence, the Gaussian, and (B0.35) give \[
f_\eta(Z)=O_{\eta,A}(Z^{A-k}).
\tag{B0.60}
\] Consequently

\[
\mathfrak F_\eta(s)=\int_0^\infty f_\eta(Z)Z^{-s+k}\,\frac{dZ}{Z}
\tag{B0.61}
\]

is holomorphic on \(\Re s>\beta_*-\Delta\). On \(\Re s>1\), Fourier
inversion in \(\log Z\), first on a fixed vertical line, identifies it
with

\[
\mathfrak F_\eta(s)=
e^{(s-5/6)^2}\frac{H_\eta(s)}{L_{\mathbb K}^S(s,\eta)}.
\tag{B0.62}
\]

The Fourier density on that line is integrable by Gaussian decay. Its
inverse transform is integrable by (B0.59)--(B0.60), which justifies the
identification as an ordinary Mellin integral. Identity on the full
half-plane \(\Re s>1\) also follows by analytic continuation from any
intervening strip of absolute integrability.

Choose \(\Delta\) additionally so that \(\beta_*-\Delta>b_0>87/100\).
Formula (B0.35) and the nonvanishing Gaussian show that the right side
of (B0.62) would have a pole at every zero of
\(L_{\mathbb K}^S(s,\allowbreak{}\eta)\) in this half-plane. Its holomorphic
continuation (B0.61) excludes all such zeros. This conclusion holds for
each primitive finite-order target, with the same real boundary
\(\beta_*-\Delta\). Constants and thresholds may depend on the target;
the excluded strip does not. It contradicts the definition of
\(\beta_*\). Thus \(\beta_*\le b_0\).

The classical nonvanishing on \(\Re s=1\) and the Euler product on
\(\Re s>1\) complete the asserted half-plane statement. Finally

\[
\zeta_{\mathbb K}(s)=\zeta(s)L(s,\chi_{-3}).
\tag{B0.63}
\]

The second factor is holomorphic in the relevant strip. A zeta zero
there would be a zero of the principal Hecke function, which has just
been excluded. Reflection by the zeta functional equation gives the
stated symmetric location of nontrivial zeros. This proves Theorem B0.

\hypertarget{appendix-b1.-finite-arithmetic-used-by-both-expansions}{%
\section{Appendix B1. Finite arithmetic used by both
expansions}\label{appendix-b1.-finite-arithmetic-used-by-both-expansions}}

This appendix fixes the phases and proves the finite identities used in
the probe, reflection and moment arguments. All moduli are primary and
prime to the fixed set \(S\). Character powers are extended by zero on
nonunits, including the zeroth and sixth powers.

\hypertarget{b1.1.-the-fixed-ray-phase-and-sextic-reciprocity}{%
\subsection{B1.1. The fixed ray phase and sextic
reciprocity}\label{b1.1.-the-fixed-ray-phase-and-sextic-reciprocity}}

For an odd \(c=a+b\omega\), put \[
\Gamma_q(c)=\frac{1+i^{-b}+i^a+i^{b-a}}2,
\qquad G(c)=\overline{\chi_c(4)}\Gamma_q(c),
\qquad R(v,w)=\frac{\Gamma_q(vw)}{\Gamma_q(v)\Gamma_q(w)}.
\tag{B1.1}
\] These formulas apply to arbitrary odd moduli, not just squarefree
ones. In fact \[
\Gamma_q(c)=q_c^{-1/2}\sum_{x\bmod c}e(x^2/c)
=\frac12\sum_{y\bmod2\mathcal O}e(-cy^2/4).
\tag{B1.2}
\] To prove the second equality without a squarefreeness restriction,
apply lattice Poisson summation to \(e(z^2/c)e^{-\pi\delta|z|^2}\). With
the self-dual measure, its transform is \[
\frac{|c|}{2\sqrt r}\exp\!\left(-\frac{\pi\delta q_c|y|^2}{4r}\right)e(-cy^2/(4r)),
\qquad r=1+3q_c\delta^2/16.
\] Divide the identity by \(2/(\sqrt3\delta)\) and let
\(\delta\downarrow0\). The original periodic mean is
\(q_c^{-1}\sum_{x\bmod c}e(x^2/c)\); the transformed mean is
\((2|c|)^{-1}\sum_{y\bmod2\mathcal O}e(-cy^2/4)\). Replacing \(r\) by
one produces an \(O_c(1)\) error before division: its phase error is
bounded by \(O_c(\delta^2\sum_y|y|^2e^{-c_1\delta|y|^2})\), and the
amplitude and damping errors are smaller. The four representatives
\(0,\allowbreak{}1,\allowbreak{}\omega,\allowbreak{}1+\omega\) give (B1.1).

On \((\mathcal O/4\mathcal O)^\times\) modulo squares, the
representatives \(1,\allowbreak{}-1,\allowbreak{}\lambda,\allowbreak{}-\lambda\) have values \(1,\allowbreak{}1,\allowbreak{}i,\allowbreak{}-i\). Thus
\(R\) is a symmetric sign-valued bicharacter, explicitly \[
R((-1)^e\lambda^f,(-1)^g\lambda^h)=(-1)^{eh+fg+fh}.
\tag{B1.3}
\] Cubic reciprocity gives \(\chi_c(4)=(c/(-2))_3\). Hence \(G\) and
\(R\) depend only on the fixed ray class modulo \(36\mathcal O\), and \[
G(vw)=G(v)G(w)R(v,w),\quad
G(p^3)=\gamma_3(p),\quad R(p,p)=\gamma_3(p)^2=\chi_p(-1).
\tag{B1.4}
\] Indeed \(p^2\) is a square unit modulo four, so
\(\Gamma_q(p^3)=\Gamma_q(p)\), while \(\chi_{p^3}(4)=1\). At a prime,
the quadratic exponential sum equals the quadratic-character Gauss sum;
its square is \(\chi_p(-1)\).

For distinct primary primes \(v,\allowbreak{}w\), cubic reciprocity makes
\((\chi_v(w)/\chi_w(v))^2=1\). The quadratic Gauss sum and CRT identify
its cube with \(R(v,\allowbreak{}w)\). A sign equals its cube. Multiplicativity and
(B1.3) therefore give the full sextic reciprocity law \[
\chi_v(w)=R(v,w)\chi_w(v),\qquad (v,w)=1,
\tag{B1.5}
\] including nonsquarefree \(v,\allowbreak{}w\).

\hypertarget{b1.2.-prime-gauss-sums-and-squarefree-coefficients}{%
\subsection{B1.2. Prime Gauss sums and squarefree
coefficients}\label{b1.2.-prime-gauss-sums-and-squarefree-coefficients}}

Write \(g_\chi(p,\allowbreak{}h)=\sum_{x\bmod p}\chi(x)e(hx/p)\), and \(Q=q_p\). The
finite Gauss--Jacobi identity, obtained by grouping the double Gauss sum
according to \(x+y\), is \[
g_A(p,1)g_B(p,1)=J(A,B)g_{AB}(p,1),\quad
J(A,B)=\sum_xA(x)B(1-x),
\tag{B1.6}
\] when \(A,\allowbreak{}B,\allowbreak{}AB\) are nonprincipal. Also \(g_Ag_{\bar A}=A(-1)Q\).

For \(\psi=\chi_p^2\), reduction of \(J(\psi,\allowbreak{}\psi)\) modulo \(p\) gives
the sum of \(x^{(Q-1)/3}(1-x)^{(Q-1)/3}\), a polynomial of degree less
than \(Q-1\), so \(p\mid J\). Its absolute value is \(\sqrt Q\). The
order-three map \(x\mapsto1/(1-x)\) preserves the summand; its two fixed
points have \(x(1-x)=1\), and every other orbit has size three.
Therefore \(J\equiv-1\pmod{3\mathcal O}\). Since \(p\) is primary, the
unit \(J/p\) is \(-1\), and \(J=-p\). Applying (B1.6) gives \[
\gamma_2(p)^3=-\alpha(p).
\tag{B1.7}
\] The changes of variable \(x=(1+y)/2\) and \(y\mapsto y/(y-1)\), with
zero extensions at the omitted points, give \[
J(\chi^2,\chi^2)=\overline{\chi(4)}^{\,2}J(\chi^3,\chi^2)
=\overline{\chi(4)}^{\,2}J(\chi,\chi^2),\qquad \chi=\chi_p.
\] Since \(\chi(4)^3=1\), (B1.6) now yields \[
\gamma_1(p)\gamma_2(p)=-\alpha(p)G(p).
\tag{B1.8}
\] CRT gives, for coprime squarefree \(a,\allowbreak{}b\), \[
\gamma_j(ab)=\gamma_j(a)\gamma_j(b)\chi_a(b)^j\chi_b(a)^j.
\tag{B1.9}
\] Cubing this formula, and using (B1.4)--(B1.8), proves \[
\gamma_2(c)^3=\mu(c)\alpha(c),\qquad
\gamma_1(c)\gamma_2(c)=\mu(c)\alpha(c)G(c),
\] \[
\mu(c)\gamma_{-1}(c)=\chi_c(-1)G(c)^{-1}\overline{\alpha(c)}\gamma_2(c).
\tag{B1.10}
\] The last formula also follows directly from
\(\gamma_1(c)\gamma_{-1}(c)=\chi_c(-1)\). Finally, for
\(a_\nu(c)=\overline{\alpha(c)}\gamma_2(c)\nu(c)\), reciprocity in
(B1.9) gives \[
a_\nu(ab)=a_\nu(a)a_\nu(b)\chi_b(a)^4.
\tag{B1.11}
\]

\hypertarget{b1.3.-imprimitive-gauss-sums-and-the-complete-correlation}{%
\subsection{B1.3. Imprimitive Gauss sums and the complete
correlation}\label{b1.3.-imprimitive-gauss-sums-and-the-complete-correlation}}

The two-variable notation \[
\mathcal G(a,h)=q_a^{-1/2}\sum_{x\bmod a}\chi_a(x)e(hx/a),\qquad \mathcal G(1,h)=1,
\tag{B1.12}
\] is distinct from the fixed ray function \(G(a)\). At a prime power
\(p^j\), with \(j\ge1\), lifting \(x=x_0+py\) first forces
\(p^{j-1}\mid h\). The remaining prime sum is nonprincipal if
\(6\nmid j\), and is a Ramanujan sum if \(6\mid j\). Consequently \[
|\mathcal G(p^j,h)|=Q^{(j-1)/2}\mathbf1_{v_p(h)=j-1}\quad(6\nmid j),
\] \[
\mathcal G(p^j,h)=Q^{j/2}\mathbf1_{p^j\mid h}-Q^{j/2-1}\mathbf1_{p^{j-1}\mid h}\quad(6\mid j).
\tag{B1.13}
\] In particular the zeros at nonunits survive every imprimitive
reduction.

For arbitrary \(u,\allowbreak{}v\), define \[
F(u,v;j)=\sum_{\substack{x\bmod u,\ y\bmod v\\vx-uy\equiv j\pmod{uv}}}\chi_u(x)\overline{\chi_v(y)}.
\tag{B1.14}
\] Expanding the Gauss sums and using additive orthogonality gives \[
F(u,v;j)=\frac1{\sqrt{q_uq_v}}\sum_{h\bmod uv}\mathcal G(u,h)\overline{\mathcal G(v,h)}e(-jh/(uv)).
\tag{B1.15}
\] At \(j=0\), a nonzero summand requires \(u\mid v\) and \(v\mid u\),
since \(x,\allowbreak{}y\) must be units. Thus \[
F(u,v;0)=\mathbf1_{u=v}\phi(u).
\tag{B1.16}
\]

To expose the common support, write \(u=Cn_1,\allowbreak{}v=Cn_2\), where \(C=(u,\allowbreak{}v)\)
and \((n_1,\allowbreak{}n_2)=1\). The congruence forces \(C\mid j\). For \(j=Ck\),
solve it first modulo \(n_1,\allowbreak{}n_2\) and use (B1.5). The exact result is \[
F(Cn_1,Cn_2;Ck)=\chi_{n_1}(k)\overline{\chi_{n_2}(-k)}R(n_1,n_2)L_C(n_1,n_2;k),
\tag{B1.17}
\] where \[
L_C(n_1,n_2;k)=\sum_{\substack{x,y\bmod C\\n_2x-n_1y\equiv k\pmod C}}\chi_C(x)\overline{\chi_C(y)}.
\] Here \(C\) may meet either residual column. The full local table, for
\(p^c\Vert C\), is

\begin{longtable}[]{@{}
  >{\raggedright\arraybackslash}p{(\columnwidth - 2\tabcolsep) * \real{0.5000}}
  >{\raggedright\arraybackslash}p{(\columnwidth - 2\tabcolsep) * \real{0.5000}}@{}}
\toprule\noalign{}
\begin{minipage}[b]{\linewidth}\raggedright
Residual columns at \(p\)
\end{minipage} & \begin{minipage}[b]{\linewidth}\raggedright
\(L_{p^c}(n_1,\allowbreak{}n_2;k)\)
\end{minipage} \\
\midrule\noalign{}
\endhead
\bottomrule\noalign{}
\endlastfoot
\(p\nmid n_1n_2,\allowbreak{} p\mid k\) & \(Q^{c-1}\chi_p(n_1/n_2)^c(Q-1)\) \\
\(p\nmid n_1n_2k,\allowbreak{} 6\nmid c\) & \(-Q^{c-1}\chi_p(n_1/n_2)^c\) \\
\(p\nmid n_1n_2k,\allowbreak{} 6\mid c\) & \(Q^{c-1}(Q-2)\) \\
\(p\) divides exactly one of \(n_1,\allowbreak{}n_2\) &
\(Q^{c-1}(Q-1)\mathbf1_{6\mid c}\mathbf1_{p\nmid k}\) \\
\end{longtable}

For the first three rows, each solution modulo \(p\) has \(Q^{c-1}\)
lifts. If \(p\mid k\), then \(x/y=n_1/n_2\). Otherwise put
\(y=(k/n_1)t\), \(x=(k/n_2)(1+t)\); the ratio \((1+t)/t\) runs through
\(\mathbf F_Q^\times\setminus\{1\}\). Its character sum is \(-1\) for a
nonprincipal power and \(Q-2\) for a principal power. If exactly one
residual column is divisible by \(p\), the other variable is a fixed
unit and the free unit-character sum gives the last row. This proves the
table and \(|L_C|\le q_C\).

For separation, extend this table to all residual pairs by setting the
local value to zero whenever \(p\mid(n_1,\allowbreak{}n_2)\). This produces a
periodic coefficient modulo \(\operatorname{rad}C\) in each column.
\textbf{It equals the original correlation only on the coprime locus.}
Accordingly insert \[
\mathbf1_{(n_1,n_2)=1}=\sum_{d\mid n_1,\ d\mid n_2}\mu(d)
\tag{B1.18}
\] before separating columns; no common prime can be omitted from this
projector.

A second form is useful when a moment induction has already extracted
the complete common prime support. If \[
u=Da,\quad v=Eb,\qquad (a,b)=(ab,DE)=1,
\] then, without any squarefreeness condition on \(D,\allowbreak{}E\), \[
F(Da,Eb;j)=F(D,E;j)R(a,E)R(b,D)R(a,b)\chi_a(j)\overline{\chi_b(-j)}.
\tag{B1.19}
\] Indeed the congruence fixes the residues modulo \(a,\allowbreak{}b\). In the
remaining congruence set \(x'=bx\pmod D\), \(y'=ay\pmod E\); the unit
factors cancel by (B1.5), leaving precisely the displayed signs. One may
extend the right side to all \((a,\allowbreak{}b)\) with \((ab,\allowbreak{}DE)=1\), retaining all
character zeros, but must again use the full projector (B1.18). This is
the common-support identity used in both moment inductions.

\hypertarget{appendix-b2.-completed-reflection-and-marked-row-energy}{%
\section{Appendix B2. Completed reflection and marked row
energy}\label{appendix-b2.-completed-reflection-and-marked-row-energy}}

This appendix proves the reflected-energy estimate used in the moment
argument. A whole-index prime mark changes the reflection conductor when
its local frequency vanishes. The proof therefore keeps the compatible
conductor branches until the row and marked-prime phases have separated;
a quadratic--cubic large sieve then bounds the remaining norm.

\hypertarget{b2.1.-the-completed-family-and-fixed-conventions}{%
\subsection{B2.1. The completed family and fixed
conventions}\label{b2.1.-the-completed-family-and-fixed-conventions}}

Write \(\mathcal O=\mathbb Z[\omega]\), \(\lambda=1+2\omega\),
\(q_a=|a|^2\), and \(\alpha(a)=a/|a|\). Throughout \(Z\ge2\). Ideals
prime to \(\lambda\) are represented by their unique primary generators
\(a\equiv1\pmod3\). Let \(S\) be a fixed finite set containing the
primes above \(6\); ``good'' means prime to \(S\). We use the sextic
characters and normalized Gauss coefficients of the arithmetic section.
Every integer power of a residue character is zero on a nonunit,
including exponent zero: \[
\chi_p(A)^0=1_{p\nmid A}.
\tag{B2.1}
\] In particular, exponent reduction modulo six never removes a
coprimality condition.

Fix a finite family of finite-ray characters \(\nu\), with conductors
supported on \(S\), extended by zero outside the primary ideals coprime
to \(S\). A squarefree good label \(f\) and a squarefree puncture
radical \(r_\rho\) may vary between invocations. Within a row norm they
are fixed, and \[
\rho(A)=1_{(A,r_\rho)=1},\qquad
\Psi_k(A)=\nu(A)\rho(A)\chi_A(k)\chi_A(f)^4.
\tag{B2.2}
\] The puncture may depend on previously frozen outer labels; it is
independent of the current row and, when \(f\) is averaged later, of
that averaging variable.

For a fixed number of pairwise disjoint prime lists \(\mathcal P_i\),
with \(q_p\asymp Z^{z_i}\) on the \(i\)-th list and \(|a_i(p)|\le1\),
put \[
d(A)=\sum_{(p_i)\in\prod_i\mathcal P_i}
       \prod_i a_i(p_i)\,1_{\prod_i p_i\mid A},
\qquad z=\sum_i z_i.
\tag{B2.3}
\] The completed sum at length \(X\) is \[
\mathcal T_d(X;k;V)=
\sum_{\substack{n\ {\rm sf},\,b\\ (nb,S)=1}}
\frac{\overline{\alpha(n)}\gamma_2(n)\Psi_k(n)\,
      \overline{\alpha(b)}^{\,3}\Psi_k(b)^3}
     {\sqrt{q_n}\,q_b}\,
d(nb^3)V(q_nq_b^3/X).
\tag{B2.4}
\] The two primal ideals may share primes. An annular smooth test is
allowed, as is any test rapidly decreasing with all Euler derivatives at
both endpoints. The factor \(y^{1/2}\) converting an annular uncompleted
sum to (B2.4) is included in \(V\).

We also allow one common smooth profile in the completed norm and the
individual marked-prime norms, supported in fixed logarithmic boxes in
the latter variables: expand (B2.3) in (B2.4), then put this profile
inside the tuple sum. Logarithmic Fourier inversion in those variables
reduces it to (B2.4), with an absolutely integrable common family of
tests \(V\). Bounds below depend on finitely many weighted Euler
seminorms of that profile. They are uniform in the prime lists, their
bounded coefficients, \(k,\allowbreak{}f,\allowbreak{}r_\rho\), and all norm labels whose
logarithms to base \(Z\) lie in fixed bounded ranges. Imaginary norm
powers incur only a fixed polynomial in their heights. Product form in
(B2.3) is required when slots are assigned to extracted factors; after
that assignment, the remaining tuple coefficient may be any bounded
function independent of the row and dual indices.

The only imported analytic results are the unconditional theta
expansions in Dunn--Radziwiłł {[}DR, §5 and Appendix A{]},
Goldmakher--Louvel's quadratic Hecke large sieve {[}GL, Definition 1 and
Theorem 1.1{]}, and Heath-Brown's cubic large sieve {[}HB, Theorem 2{]}.
All moving-conductor and marked-index reductions are proved below.

\hypertarget{b2.2.-exact-reflection-with-the-local-zeros-retained}{%
\subsection{B2.2. Exact reflection, with the local zeros
retained}\label{b2.2.-exact-reflection-with-the-local-zeros-retained}}

\hypertarget{b2.2.1.-the-finite-theta-source}{%
\subsubsection{B2.2.1. The finite theta
source}\label{b2.2.1.-the-finite-theta-source}}

Put \(\check e(z)=\exp(2\pi i(z+\bar z))\) and
\(e(z)=\check e(z/\lambda)\). The theta input {[}DR, (5.4), (5.7),
(5.9), (5.12)--(5.17), Appendix A{]} gives a smooth cubic theta function
on \(\mathbb H^3\), automorphic on \[
\Gamma_2=\langle {\rm SL}_2(\mathbb Z),\Gamma(3)\rangle,\qquad
\theta(gw)=\kappa(g)\theta(w),\qquad
\kappa\begin{pmatrix}a&b\\c&d\end{pmatrix}
=\left(\frac ca\right)_3\quad(g\in\Gamma(3),\ c\ne0).
\tag{B2.5}
\] The multiplier is one on \({\rm SL}_2(\mathbb Z)\) and on the \(c=0\)
elements of \(\Gamma(3)\). Let \[
L(t)=\begin{pmatrix}1&0\\t&1\end{pmatrix},\qquad
\Theta_j(w)=\theta(L(t_j)w),\qquad
(t_0,t_1,t_2)=(0,\omega^2,\omega).
\] Translations by \(\mathbb Z+3\mathcal O\) have multiplier one. Thus
\(L(\lambda)\) and \(L(-\lambda)\) give \(\Theta_1,\allowbreak{}\Theta_2\). Also, for
\(u_0=s+t\omega\), \(0\le s,\allowbreak{}t<3\),
\(\theta\bigl(\begin{smallmatrix}u_0&-1\\1&0\end{smallmatrix}w\bigr)=\Theta_t(w)\).

For each of these three sources define \(d_H\) by \[
\overline{\theta(H(z,v))}
=C_H(v)+\sum_{\mu\ne0}d_H(\mu)\,vK_{1/3}(4\pi|\mu|v)\check e(\mu z).
\tag{B2.6}
\] Their support and coefficient bound are \[
\mu=u\lambda^h nb^3,\quad h\ge-4,\quad n,b\equiv1\pmod3,\quad n\ {\rm sf},
\qquad |d_H(\mu)|\le27\,3^{h/6}|b|.
\tag{B2.7}
\] Here \(u\) is an Eisenstein unit. The dual \(n,\allowbreak{}b\) may share primes
and may contain primes of \(S\) other than \(\lambda\). This is the full
source support; the primal \(S\)-mask is not imposed on it.

For clarity, the bound follows directly from the cited coefficient
formulas. Their ordinary cubic Gauss sums \[
\widetilde g_3(n)=q_n^{-1/2}\sum_{v\bmod n}\left(\frac vn\right)_3\check e(v/n)
\] have modulus one for every primary squarefree \(n\), by the
finite-field Gauss identity and CRT. This includes the prime \(-2\) of
norm \(4\), where only the cubic symbol is used. The possible magnitudes
in the three sources are \(3^{h/6+8/3}|b|\), \(3^{h/6+3}|b|\), and
\(9|b|\) at \(h=-4\), all bounded by (B2.7). Valuations make the
displayed \(n,\allowbreak{}b\) unique when the representation exists. At infinity the
condition \(\lambda^3\mu\equiv1\pmod3\) selects exactly \[
d_I(\lambda^{-3}nb^3)=3^{5/2}|b|\widetilde g_3(n),\qquad
\widetilde g_3(n)=\chi_n(\lambda)^{-2}\gamma_2(n)\quad((n,S)=1).
\tag{B2.8}
\] No coprimality between \(n\) and \(b\) is required.

\hypertarget{b2.2.2.-all-nonzero-element-rows-use-one-fixed-modulus}{%
\subsubsection{B2.2.2. All nonzero element rows use one fixed
modulus}\label{b2.2.2.-all-nonzero-element-rows-use-one-fixed-modulus}}

Write a nonzero row uniquely as \(k=u k_S k_{\rm good}\), using fixed
generators at \(S\) and primary generators elsewhere. The established
reciprocity identity gives, on primary \(S\)-units, \[
\chi_A(k)=\chi_A(uk_S)R(A,k_{\rm good})
          \prod_{p\mid k_{\rm good}}\chi_p(A)^{v_p(k)}.
\tag{B2.9}
\] It holds at collisions too: both sides are zero, including when
\(6\mid v_p(k)\). Fix \(u\), the \(S\)-valuations modulo six, and the
class of \(k_{\rm good}\) in the fixed reciprocity ray group. The first
two factors, multiplied by \(\nu\), then belong to one fixed finite
family of characters \(\Psi_0\) supported on \(S\). The supplementary
characters \(A\mapsto\chi_A(uk_S)\) have conductor supported there; only
the finitely many exponents modulo six occur. The fourth-power label
introduces no reciprocity sign, since \(R(A,\allowbreak{}f)^4=1\).

Choose once a modulus \(E\) supported exactly on \(S\), divisible by the
conductors of this finite family. Choose a sufficiently divisible \(L\),
with the same prime support and \(18\mid L\), \(E\mid L\), such that \[
\varphi(x)=
\begin{cases}
\chi_x(\lambda)^2\Psi_0(x),&x\equiv1\pmod3,\ (x,E)=1,\\
0,&\text{otherwise}
\end{cases}
\] is \(L\)-periodic. Such an \(L\) exists because the cubic
supplementary law makes \((\lambda/x)_3\) periodic modulo \(9\), and the
other characters and zero mask have fixed moduli. The zero branch is
evaluated before either character. Set \[
\widehat\varphi(h_0)=q_L^{-1}\sum_{x\bmod L}\varphi(x)e(-h_0x/L),
\qquad \mathfrak M=\lambda^{12}L^4.
\tag{B2.10}
\] Every good prime of \(k,\allowbreak{}f,\allowbreak{}r_\rho\) stays in a moving local set
\(\mathcal P\), with its exponents combined modulo six. No moving good
prime enters \(L\). Thus it is enough to reflect
\(\Psi(A)=\Psi_0(A)\prod_{p\in\mathcal P}\chi_p(A)^{j_p}\),
\(0\le j_p<6\).

\hypertarget{b2.2.3.-the-reflection-identity}{%
\subsubsection{B2.2.3. The reflection
identity}\label{b2.2.3.-the-reflection-identity}}

For \(m\not\equiv0\pmod6\), let \[
\tau_{p,m}^{\pm}=q_p^{-1/2}\sum_{y\bmod p}\chi_p(y)^m e(\pm y/p),
\qquad |\tau_{p,m}^{\pm}|=1.
\] An active branch consists of \(h_0\bmod L\) and a subset of
\(\mathcal P\) containing every prime with \(j_p\ne0\). Write
\(r=\prod_{p\ {\rm active}}p\). Reduce \(\lambda^2h_0/L=a_F/c_F\),
normalizing its numerator to be primary if \(\lambda\mid c_F\), and its
denominator otherwise. Put \(c=c_Fr\). For active \(p\) define, modulo
\(p\), \[
\sigma_p=\lambda^2c/p,\qquad
\epsilon_p=-\bigl((\lambda^3c/p)\sigma_p\bigr)^{-1}
          =-\lambda^{-5}(c/p)^{-2},
\] \[
B_{p,j}(x)=
\begin{cases}
\chi_p(x)^{-j-2},&j\ne0,4,\\
q_p^{-1/2}(-1+q_p1_{p\mid x}),&j=4,\\
q_p^{-1/2}\chi_p(x)^{-2},&j=0,
\end{cases}
\quad
\omega_{p,j}=
\begin{cases}
\tau_{p,j}^-\tau_{p,j+2}^+\chi_p(\epsilon_p)^{-j-2},&j\ne0,4,\\
\tau_{p,4}^-,&j=4,\\
-\tau_{p,2}^+\chi_p(\epsilon_p)^{-2},&j=0.
\end{cases}
\tag{B2.11}
\] Indices on \(\tau\) are reduced modulo six.

Define \[
K_*=\frac{(2\pi)^4}{27},\qquad
\mathcal R(t)=\prod_\pm\frac{\Gamma(1+t\pm1/6)}{\Gamma(1-t\pm1/6)},\qquad
V^\sharp(x)=\frac1{2\pi i}\int_{(\sigma)}
\widehat V(-t)\mathcal R(t)(K_*x)^{-t}\,dt,\quad \sigma\ge0.
\tag{B2.12}
\] The integral is independent of the nonnegative line. Write
\(\mathcal T_1(X;\Psi;V)\) for (B2.4) with no mark and with \(\Psi_k\)
replaced by \(\Psi\). For each branch, there are a source \(d\) from
(B2.6), a fixed-modulus additive character \(\vartheta\) on
\(\mathcal O\), and a unit scalar \(\kappa_F\), all determined by
\(h_0\) and \(r\bmod\mathfrak M^2\), such that \[
\begin{aligned}
\mathcal T_1(X;\Psi;V)
&=\sum_{\rm branches}\zeta
\sum_{0\ne\mu\in\lambda^{-4}\mathcal O}
\frac{d(\mu)\alpha(\mu)\vartheta(\lambda^4\mu)}{\sqrt{q_\mu}}
\prod_{p\ {\rm active}}B_{p,j_p}(\lambda^4\mu)
V^\sharp(q_\mu X/q_c^2),\\
\zeta&=-\frac{i}{81}\overline{\alpha(c)}^{\,2}
\widehat\varphi(h_0)\overline{\kappa_F}
\prod_{p\ {\rm inactive}}(1-q_p^{-1})
\prod_{p\ {\rm active}}\chi_p(\sigma_p)^{-2}\omega_{p,j_p}.
\end{aligned}
\tag{B2.13}
\] All dual sums converge absolutely and \(|\zeta|\le1/81\). There are
\(O_L(2^{|\mathcal P|})\) branches, or \(O_L(4^{|\mathcal P|})\) after
splitting the two terms of every \(j=4\) factor. The branch choices are
simultaneous across local sets: adding an inactive \(j=0\) prime changes
only the scalar, by \(1-q_p^{-1}\). It leaves the conductor, source,
additive character, other local phases, and kernel unchanged.

\textbf{Proof of the conductor and phase assertions.} Finite Fourier
inversion gives \[
C_{p,j}(h)=q_p^{-1}\sum_{x\bmod p}\chi_p(x)^j e(-hx/p)
=
\begin{cases}
q_p^{-1/2}\tau_{p,j}^-\chi_p(h)^{-j},&j\ne0,\ h\ne0,\\
0,&j\ne0,\ h=0,\\
-q_p^{-1},&j=0,\ h\ne0,\\
1-q_p^{-1},&j=0,\ h=0.
\end{cases}
\tag{B2.14}
\] Hence the masked theta function is \[
F(z,v)=\sum_{h_0,(h_p)}
\widehat\varphi(h_0)\prod_pC_{p,j_p}(h_p)\,
\overline{\theta(z+z_h,v)},\qquad
z_h=\lambda^2\left(h_0/L+\sum_ph_p/p\right).
\tag{B2.15}
\] At an infinity index \(\lambda^3\mu=nb^3\), the whole-index mask
first removes \((nb,\allowbreak{}E)>1\). On the remaining domain its coefficient
multiplier is \(\chi_n(\lambda)^2\Psi(n)\Psi(b)^3\), also at a shared
moving prime. Equations (B2.8) therefore identify its direct series with
the completed sum.

Fix \(h_0\) and the active frequencies. Then \[
z_h=a/c,\qquad c=c_Fr,\qquad
a=a_Fr+\lambda^2c_F\sum_{p\mid r}(r/p)h_p.
\tag{B2.16}
\] At a prime of \(c_F\), \(a\equiv a_Fr\) is a unit; at \(p\mid r\),
\(a\equiv\sigma_ph_p\) is a unit. Thus the fraction is reduced even for
nonunit \(h_0\). Lift every \(h_p\) to be \(0\bmod\mathfrak M^2\) while
retaining its class modulo \(p\), and fix \(r\bmod\mathfrak M^2\). Then
\(a\bmod\mathfrak Mc_F\) is fixed. Choose \(\delta'\) by CRT: \[
\begin{gathered}
\delta'\equiv a^{-1}\pmod{\mathfrak Mc_F}\text{ at primes of }c_F,\\
\delta'\equiv0\pmod{\mathfrak M}\text{ at other primes of }\mathfrak M,\\
\delta'\equiv a^{-1}\pmod p\ (p\mid r).
\end{gathered}
\] Set \(b=(a\delta'-1)/c\) and
\(g=\bigl(\begin{smallmatrix}a&b\\c&\delta'\end{smallmatrix}\bigr)\).
Then \(b\equiv0\bmod\mathfrak M\) at primes of \(c_F\), and
\(b\equiv-c^{-1}\bmod\mathfrak M\) elsewhere. All required fixed-prime
residues of the matrix are therefore independent of the active
frequencies. Periodic changes of lifts, followed by adding a multiple of
the CRT modulus to \(\delta'\), avoid zero matrix entries if a symbol
calculation requires this.

There are three cases. If \(v_\lambda(c)\ge2\), take \(H=I\); then
\(g\in\Gamma(3)\) and \(\kappa(g)=(c_F/a)_3(a/r)_3\). If
\(v_\lambda(c)=1\), choose \(u_0\in\{\lambda,\allowbreak{}-\lambda\}\) with
\(u_0\equiv c\pmod3\), set \(H=L(u_0)\), and write \(A=a-u_0b\). Then
\(G=gH^{-1}\in\Gamma(3)\), \(a,\allowbreak{}A\) are coprime primary elements, and \[
a(c-u_0\delta')=cA-u_0,\qquad bc\equiv-1\pmod a.
\] Cubic reciprocity gives \[
\kappa(G)=
\left(\frac{-u_0}{A}\right)_3\left(\frac aA\right)_3^{-1}
=\left(\frac{-u_0}{A}\right)_3
 \left(\frac{c_F/u_0}{a}\right)_3\left(\frac ar\right)_3.
\] If \(\lambda\nmid c\), choose \(u_0=s+t\omega\equiv a\pmod3\),
\(0\le s,\allowbreak{}t<3\), and
\(H=\bigl(\begin{smallmatrix}u_0&-1\\1&0\end{smallmatrix}\bigr)\). Then
\[
G=gH^{-1}=
\begin{pmatrix}-b&a+bu_0\\-\delta'&c+\delta'u_0\end{pmatrix}\in\Gamma(3).
\] Both \(-b,\allowbreak{}c\) are primary and \(-bc=1-a\delta'\equiv1\pmod9\).
Factoring \(\delta'=u\lambda^h\delta_0\), the supplementary laws and
reciprocity give \((\delta'/(-bc))_3=1\): the unit and
\(\lambda\)-factors use the congruence modulo \(9\), and after
reciprocity the remaining numerator is \(1\bmod\delta_0\). Consequently
\(\kappa(G)=(a/c_F)_3(a/r)_3\). In every case, \[
\kappa(G)=\kappa_F\prod_{p\mid r}\chi_p(a)^2,
\tag{B2.17}
\] where \(\kappa_F\) and the choice of \(d_H\) are fixed by the
displayed sector. This proves the needed conductor uniformity; it does
not assume a unit-frequency support condition on \(\widehat\varphi\).

For \(x=\lambda^4\mu\) and \(D_F=\lambda^3c_F\), additive CRT gives \[
\check e(-\delta'\mu/c)
=\vartheta(x)\prod_{p\mid r}e(\epsilon_ph_p^{-1}x/p),\qquad
\vartheta(x)=e(-\delta_F'r^{-1}x/D_F).
\tag{B2.18}
\] The fixed congruences determine \(\delta_F'\bmod D_F\), and the
sector determines \(r^{-1}\). Conjugating (B2.17) supplies
\(\overline{\kappa_F}\prod_p\chi_p(\sigma_p)^{-2}\chi_p(h_p)^{-2}\).
Multiply by (B2.14) and replace \(h_p^{-1}\) by \(y\). For \(j\ne0,\allowbreak{}4\)
the resulting sum is a Gauss sum of exponent \(j+2\); for \(j=4\) it is
\(\sum_{y\ne0}e(\epsilon_pxy/p)=-1+q_p1_{p\mid x}\); an active \(j=0\)
gives the last line of (B2.11). These are exactly
\(\omega_{p,\allowbreak{}j}B_{p,\allowbreak{}j}(x)\), including their zeros. An absent prime and
an inactive zero-exponent prime have identical active radical, so their
only difference is the coefficient \(1-q_p^{-1}\). This also proves
simultaneous branch compatibility.

\textbf{Proof of the Mellin normalization and continuation.} Put \[
D(s,\Psi)=\sum_{n\ {\rm sf}}\overline{\alpha(n)}\gamma_2(n)\Psi(n)q_n^{-s},
\quad L_0(s,\Psi)=\sum_b\overline{\alpha(b)}^{\,3}\Psi(b)^3q_b^{-3s+1/2},
\quad T=L_0D.
\] The sums are over primary \(S\)-units and converge absolutely for
\(\Re s>1\). For
\(J(s)=\int_0^\infty\partial_{\bar z}F(0,\allowbreak{}v)v^{2s-1}\,\allowbreak{}dv\), use \[
\int_0^\infty K_{1/3}(y)y^{w-1}\,dy
=2^{w-2}\Gamma((w-1/3)/2)\Gamma((w+1/3)/2).
\] Since
\(\partial_{\bar z}\check e(\mu z)=2\pi i\bar\mu\check e(\mu z)\) and
\(\overline{\alpha(\lambda^{-3})}=-i\), (B2.8) yields \[
J(s)=\frac{3^{5/2}}4\left(\frac{27}{(2\pi)^2}\right)^s
\Gamma(s+1/3)\Gamma(s+2/3)T(s,\Psi).
\tag{B2.19}
\] For one translated cusp \(g=GH\), its inverse action is \[
g^{-1}(z+a/c,v)=
\left(-\delta'/c-\frac{\bar z}{c^2(v^2+|z|^2)},
      \frac{v}{q_c(v^2+|z|^2)}\right).
\] At \(z=0\), the chain rule changes \(\partial_{\bar z}\) into
\(-\overline{\kappa(G)}(cv)^{-2}\partial_z\) at
\((-\delta'/c,\allowbreak{}1/(q_cv))\). The derivative removes the constant Fourier
term at every cusp. Substituting \(v\mapsto1/(q_cv)\) and the Bessel
integral gives, for \(\Re s<0\), \[
-\frac i4(2\pi)^{2s-2}\overline{\alpha(c)}^{\,2}q_c^{1-2s}
\overline{\kappa(G)}
\prod_\pm\Gamma(3/2-s\pm1/6)
\sum_{\mu\ne0}d_H(\mu)\alpha(\mu)\check e(-\delta'\mu/c)q_\mu^{-(1-s)}.
\tag{B2.20}
\] The series converges absolutely there by (B2.7). The differentiated
theta function decays exponentially at both height endpoints, using the
direct expansion at infinity and this inverse-cusp formula at zero.
Repeated integration by parts in \(\log v\) makes \(J\) entire and
rapidly decreasing on fixed vertical strips. Equation (B2.19) continues
\(T\) to an entire function. The direct series bounds it on a line right
of \(1\), and (B2.20), divided by (B2.19), gives polynomial growth left
of \(0\). Within the intervening strip, the preliminary bound from \(J\)
is \(O((1+|s|)^C e^{\pi|\Im s|})\). Apply the maximum principle to
\(T(s)(B+s)^{-C'}e^{\epsilon s^2}\), choosing \(-B\) outside the strip;
then let \(\epsilon\downarrow0\). This gives polynomial strip growth.

Mellin inversion expresses (B2.4), with no marks, as
\((2\pi i)^{-1}\int_{(a)}\widehat V(s-1/2)X^{s-1/2}T(s,\allowbreak{}\Psi)\,\allowbreak{}ds\),
\(a>1\). Shift to \(\Re s<0\), insert (B2.20)/(B2.19), and put
\(t=1/2-s\). The gamma quotient is \(\mathcal R(t)\); the remaining
factor is \[
-\frac i{81}\overline{\alpha(c)}^{\,2}q_\mu^{-1/2}
(K_*q_\mu X/q_c^2)^{-t}.
\] Absolute convergence and rapid Mellin decay justify both interchanges
and the contour shift. Combining the finite local factors proves
(B2.13).

\hypertarget{b2.3.-common-profiles-and-whole-series-tails}{%
\subsection{B2.3. Common profiles and whole-series
tails}\label{b2.3.-common-profiles-and-whole-series-tails}}

For \(A,\allowbreak{}B\ge0\), with \(B\) an integer, set \[
\mathcal M_{A,B}(V)=
\sup_{-A\le\eta\le1/4}\int_{\mathbb R}
(1+|t|)^B|\widehat V(\eta+it)|\,dt.
\] The first negative numerator pole of \(\mathcal R\) is \(-5/6\), and
on \(-1/4\le\sigma\le A\),
\(|\mathcal R(\sigma+it)|\ll_A(1+|t|)^{4\sigma}\). Shifting (B2.12) to
\(-1/4,\allowbreak{}0,\allowbreak{}A\) therefore proves \[
|(x\partial_x)^jV^\sharp(x)|
\ll_{A,j}\min\{x^{1/4},(1+x)^{-A}\}
\mathcal M_{A,\lceil4A\rceil+j+2}(V).
\tag{B2.21}
\] These are finite seminorms: \(B+2\) integrations by parts bound
\(\mathcal M_{A,\allowbreak{}B}\) by \[
\ll_{A,B}\int_0^\infty
(1+x^{-A}+x^{1/4})\bigl(|V(x)|+|(x\partial_x)^{B+2}V(x)|\bigr)\frac{dx}{x}.
\] Replacing \(V\) by \(x^{it_0}V\) costs at most \((1+|t_0|)^B\).

On a fixed compact logarithmic box, let
\(F_Y(\mathbf y)=W(\mathbf y)V^\sharp(Y\prod_i y_i^{a_i})\), where \(Y\)
is constant within the block. Euler differentiation and (B2.21) bound
every fixed derivative of \(F_Y\) by \[
m_A(Y)=\min\{Y^{1/4},(1+Y)^{-A}\}
\] times fixed seminorms. Fourier integration by parts in all
logarithmic variables shows that the common separating density has
weighted \(L^1\)-norm bounded by the same scalar. For a row vector
linear in this whole profile, write \(F_Y=m_A(Y)\widetilde F_Y\), apply
its single common Fourier inversion, and then Minkowski: \[
\|U(F_Y)\|_{\ell^2(I)}
\le m_A(Y)\int |\widehat{\widetilde F_Y}(\mathbf t)|
\|B(\mathbf t)\|_{\ell^2(I)}\,\frac{d\mathbf t}{(2\pi)^d}.
\tag{B2.22}
\] Only now may the positive row set be enlarged. The separated formula
defines \(B\) on the added rows; no kernel estimate is asserted there.
Squaring gives \(m_A(Y)^2\). In particular the small-argument gain for a
linear reflected sum is \(Y^{1/2}\) in its energy. Choosing separate
Fourier densities for different rows would not justify (B2.22).

For either lattice \(\Lambda=\mathcal O,\allowbreak{}\lambda^{-4}\mathcal O\), if
\(|K(x)|\le C_Ax^{-A}\) for \(x\ge1\), \(A>1\), then \[
\sum_{\substack{0\ne\mu\in\Lambda\\a q_\mu>U}}|K(aq_\mu)|
\ll_A C_A(1+a^{-1})U^{1-A}\qquad(a>0,\ U\ge1).
\tag{B2.23}
\] Indeed the shell \(2^jU<aq_\mu\le2^{j+1}U\) has \(O(1+2^jU/a)\)
lattice points; sum the two resulting geometric series. This estimate
controls the entire discarded dual series, without presupposing a bound
on its retained indices.

\hypertarget{b2.4.-the-quadraticcubic-norm-including-both-collision-masks}{%
\subsection{B2.4. The quadratic--cubic norm, including both collision
masks}\label{b2.4.-the-quadraticcubic-norm-including-both-collision-masks}}

\hypertarget{b2.4.1.-imported-sieves-and-their-domains}{%
\subsubsection{B2.4.1. Imported sieves and their
domains}\label{b2.4.1.-imported-sieves-and-their-domains}}

For squarefree primary good ideals, the quadratic sieve used here is \[
\sum_{\substack{a\ {\rm sf}\\q_a\le U}}
\left|\sum_{\substack{b\ {\rm sf}\\q_b\le V}}c_b\chi_a(b)^3\right|^2
\ll_\epsilon(UV)^\epsilon(U+V)\sum_b|c_b|^2.
\tag{B2.24}
\] Here is the family check needed to apply {[}GL{]}. The quadratic
symbol at the primary generator of an ideal \((x)\), initially on
\(\lambda\)-units, is \[
\vartheta_a((x))=\epsilon_\lambda(x)^{e(a)}(x/a)_2,\qquad
e(a)=(q_a-1)/2\bmod2,
\] where \(\epsilon_\lambda\) is the nontrivial character modulo
\(\lambda\). The powers of \(\omega\) used to make a generator primary
are squares; its sign contributes exactly the displayed
\(\lambda\)-character. CRT gives primitive conductor \(a\lambda^{e(a)}\)
and trivial infinite type. When \(e(a)=0\), take the primitive inducing
character with the \(\lambda\)-factor absent. In a fixed primary square
class modulo \(4\), two indices have the same \(e(a)\); for coprime such
indices the primitive product character has conductor exactly their
product. The established reciprocity factor is a bicharacter of a fixed
ray group. These are the three conditions of {[}GL, Definition 1{]}; its
theorem and transpose duality give (B2.24).

Heath-Brown's cubic sieve gives \[
\sum_{\substack{a\ {\rm sf},\,a\equiv1(3)\\q_a\le U}}
\left|\sum_{\substack{b\ {\rm sf},\,b\equiv1(3)\\q_b\le V}}
c_b(b/a)_3\right|^2
\ll_\epsilon(UV)^\epsilon
\{U+V+(UV)^{2/3}\}\sum_b|c_b|^2.
\tag{B2.25}
\] Squarefreeness is in \(\mathcal O\); squarefree rational norms and an
exclusion of rational prime factors are not hypotheses. Conjugation and
cubic reciprocity permit the opposite orientation. In either sieve,
fixed row restrictions decrease the positive outer sum, and fixed column
restrictions are implemented by zero coefficients. A row--column
collision mask must first be resolved.

\hypertarget{b2.4.2.-quadratic-reduction-for-a-squarefree-index-times-a-cube}{%
\subsubsection{B2.4.2. Quadratic reduction for a squarefree index times
a
cube}\label{b2.4.2.-quadratic-reduction-for-a-squarefree-index-times-a-cube}}

Let \(k\) be squarefree good with \(q_k\ll K\), \(n\) squarefree primary
with \(q_n\asymp N\), and \(b\) primary with \(q_b\asymp B\). The dual
\(n,\allowbreak{}b\) may contain primes of \(S\) other than \(\lambda\). Let
\(\beta(n,\allowbreak{}b)\) be arbitrary, independent of \(k\); fixed support
restrictions are allowed.

Write uniquely \[
b=gt^2,\qquad c=(n,g),\qquad n=cm,\qquad g=ch,
\] where \(g\) is squarefree and \(c,\allowbreak{}m,\allowbreak{}h\) are pairwise coprime
squarefree ideals. There is no restriction \((g,\allowbreak{}t)=1\). On dyads
\(q_c\asymp C,\allowbreak{}\ q_g\asymp G,\allowbreak{}\ q_t\asymp T\), we have \(B\asymp GT^2\)
and the exact identity \[
\chi_k(nb)^3=\chi_k(mh)^3\,1_{(k,ct)=1}.
\tag{B2.26}
\] Cauchy in \(t\) costs \(O(T)\). Fix \(t\), retain \((k,\allowbreak{}t)=1\) as a
row restriction, and expand the entire remaining mask:
\(1_{(k,\allowbreak{}c)=1}=\sum_{r\mid(k,\allowbreak{}c)}\mu(r)\). Rowwise divisor Cauchy costs
\((KNB)^\epsilon\). For fixed \(r\), write \(k=rk'\), \(c=rc'\).
Squarefreeness gives fixed restrictions \((k',\allowbreak{}r)=(c',\allowbreak{}r)=1\), and terms
with \((r,\allowbreak{}t)>1\) vanish. There is no remaining \((k',\allowbreak{}c')\)-mask in an
individual \(r\)-summand.

The factor \(\chi_r(mh)^3\) is column-only. Split the finitely many
squarefree \(S\)-parts of \(m,\allowbreak{}h\), whose characters are row scalars.
Their good product \(j=m_{\rm good}h_{\rm good}\) is squarefree with
norm \(O(NG/C^2)\). Group by \(j\), use (B2.24) with row length
\(K/q_r\), then Cauchy over the divisor-bounded factorizations of \(j\)
and the \(O(C/q_r)\) choices of \(c'\). If \(\mathcal D_r(t)\) retains
all original conditions on \((c',\allowbreak{}m,\allowbreak{}h)\), the resulting bound is \[
\begin{aligned}
Q_{C,G,T}(\beta)
\ll_\epsilon (KNB)^\epsilon T
\sum_{q_t\asymp T}
\sum_{\substack{r\ {\rm sf,good}\\q_r\ll\min(K,C)}}
\left(\frac K{q_r}+\frac{NG}{C^2}\right)\frac C{q_r}
\sum_{(c',m,h)\in\mathcal D_r(t)}
|\beta(rc'm,rc'ht^2)|^2.
\end{aligned}
\tag{B2.27}
\] This more precise positive bound will preserve the tuple independence
in the next step. In particular, for \(|\beta|\le1\), the triple count
is \(O(NG/(q_rC))\). In a dyad \(q_r\asymp r_0\), counting its
\(O(r_0)\) divisors and the \(O(T)\) values of \(t\) gives \[
\ll (KNB)^\epsilon\frac{NB}{r_0}
\left(\frac K{r_0}+\frac{NG}{C^2}\right)
\ll(KNB)^\epsilon NB(K+NB).
\] Triangle inequality over the logarithmically many \(C,\allowbreak{}G,\allowbreak{}T\)-dyads
proves \[
\sum_k\left|\sum_{n,b}\beta(n,b)\chi_k(nb)^3\right|^2
\ll_\epsilon(KNB)^\epsilon(K+NB)NB.
\tag{B2.28}
\] Subunit quotient ranges are empty; bounded ranges, including the unit
ideal, are covered by fixed comparison constants.

\hypertarget{b2.4.3.-hybrid-norm}{%
\subsubsection{B2.4.3. Hybrid norm}\label{b2.4.3.-hybrid-norm}}

Let \(P\) be squarefree good with \(q_P\asymp L\). Suppose
\(|a(P)|,\allowbreak{}|\beta(n,\allowbreak{}b)|\le1\), with \(a\) independent of \(k,\allowbreak{}n,\allowbreak{}b\) and
\(\beta\) independent of \(k,\allowbreak{}P\). Then \[
\begin{aligned}
\mathcal H&=
\sum_k\left|
\sum_P\frac{a(P)}{\sqrt{q_P}}1_{(P,k)=1}
\sum_{n,b}\beta(n,b)\chi_k(nb)^3\chi_P(n)^{-2}1_{(P,b)=1}
\right|^2\\
&\ll_\epsilon(KNBL)^\epsilon
(K+NB)B\{N+L+(NL)^{2/3}\}.
\end{aligned}
\tag{B2.29}
\] Fixed support restrictions and finite ray sectors are allowed if they
preserve those two independence conditions.

To prove this, expand the complete mask \(1_{(P,\allowbreak{}k)=1}\) over
\(D\mid(P,\allowbreak{}k)\), and apply rowwise divisor Cauchy. Put \(P=DP'\),
\(k=Dk'\), and \[
a_D(P')=a(DP')1_{(D,P')=1},\qquad
\beta_D(n,b)=\beta(n,b)\chi_D(nb)^3\chi_D(n)^{-2}1_{(D,b)=1}.
\] The energy obtains a factor \(q_D^{-1}\); the new row and tuple
lengths are \(K/q_D,\allowbreak{}L/q_D\). No \((P',\allowbreak{}k')\)-mask is reinstated. Apply
(B2.27) to \(k'\), treating \[
\beta_D(n,b)\sum_{P'}a_D(P')q_{P'}^{-1/2}
\chi_{P'}(n)^{-2}1_{(P',b)=1}
\] as its row-independent coefficient. In its notation the positive
output for one \(C,\allowbreak{}G,\allowbreak{}T\)-block is bounded by \[
\begin{aligned}
(KNBL)^\epsilon
\sum_D\frac1{q_D}T\sum_{q_t\asymp T}\sum_r
\left(\frac K{q_Dq_r}+\frac{NG}{C^2}\right)\frac C{q_r}
\sum_{c',m,h}|\beta_D(rc'm,rc'ht^2)|^2\\
\hspace{15mm}\cdot
\left|\sum_{P'}\frac{a_{D,r}(P')}{\sqrt{q_{P'}}}
\chi_{P'}(c'm)^{-2}1_{(P',ht)=1}\right|^2,
\qquad a_{D,r}(P')=a_D(P')\chi_{P'}(r)^{-2}.
\end{aligned}
\tag{B2.30}
\] The identity used to obtain it, with zeros included, is \[
\chi_{P'}(rc'm)^{-2}1_{(P',rc'ht^2)=1}
=\chi_{P'}(r)^{-2}\chi_{P'}(c'm)^{-2}1_{(P',ht)=1}.
\] The character supplies the zeros at \(r,\allowbreak{}c',\allowbreak{}m\).

Only in this positive sum discard \(|\beta_D|^2\le1\). For fixed
\(h,\allowbreak{}t,\allowbreak{}D,\allowbreak{}r\), the \(P'\)-coefficients are independent of \(j=c'm=n/r\),
and their squared mass is
\(\sum_{P'}q_{P'}^{-1}|a_{D,\allowbreak{}r}(P')|^2 1_{(P',\allowbreak{}ht)=1}\ll(KNBL)^\epsilon\).
Group \(c',\allowbreak{}m\) by the squarefree primary \(j\), with norm \(O(N/q_r)\).
Its allowed \(S\)-parts remain. The cubic sieve (B2.25), then the
\(O(G/C)\) choices of \(h\), give, on \(q_D\asymp D_0,\allowbreak{}\ q_r\asymp r_0\),
\[
\ll(KNBL)^\epsilon\frac GC
\left\{\frac N{r_0}+\frac L{D_0}+
\left(\frac{NL}{r_0D_0}\right)^{2/3}\right\}.
\] The \(t\)-sum and the factor \(C/q_r\) contribute
\(T^2(C/r_0)(G/C)\asymp B/r_0\). There are \(O(r_0)\) choices of \(r\),
and \(\sum_{q_D\asymp D_0}q_D^{-1}\ll1\). Replacing \(D_0,\allowbreak{}r_0\) by \(1\)
only increases the sieve factors; \(NG/C^2\ll NB\). Summing the dyads
proves (B2.29), without introducing either \((g,\allowbreak{}t)=1\) or a new
pair-dependent mask.

\hypertarget{b2.5.-whole-index-marks-and-separation-of-their-phases}{%
\subsection{B2.5. Whole-index marks and separation of their
phases}\label{b2.5.-whole-index-marks-and-separation-of-their-phases}}

For a set \(\mathcal Q\) of marked primes disjoint from the base local
set, expand \(\prod_{q\in\mathcal Q}1_{q\mid A}=\prod_q(1-\chi_q(A)^0)\)
and apply (B2.13). Its simultaneous branch compatibility makes the
inactive contribution at \(q\) \[
1-(1-q_q^{-1})=q_q^{-1}.
\] The active contribution occurs only in the subtracted zero-mask
transform. It contributes \(q_q^{-1/2}\chi_q(\lambda^4\mu)^{-2}\), with
scalar \(-\omega_{q,\allowbreak{}0}\). Thus the marked reflection is \[
\begin{aligned}
\mathcal T_{\mathcal Q}(X;\Psi;V)
&=\sum_{\mathcal B}\zeta_{\mathcal Q}
\sum_{\mu\ne0}\frac{d(\mu)\alpha(\mu)\vartheta(\lambda^4\mu)}{\sqrt{q_\mu}}
\prod_{p\in\mathcal A_0}B_{p,j_p}(\lambda^4\mu)
\prod_{q\in A_{\mathcal Q}}q_q^{-1/2}\chi_q(\lambda^4\mu)^{-2}
V^\sharp(q_\mu X/q_c^2),\\
\zeta_{\mathcal Q}
&=-\frac i{81}\overline{\alpha(c)}^{\,2}
\widehat\varphi(h_0)\overline{\kappa_F}
\prod_{p\in\mathcal I_0}(1-q_p^{-1})\prod_{q\in I_{\mathcal Q}}q_q^{-1}
\prod_{p\in\mathcal A_0}\chi_p(\sigma_p)^{-2}\omega_{p,j_p}
\prod_{q\in A_{\mathcal Q}}\chi_q(\sigma_q)^{-2}(-\omega_{q,0}).
\end{aligned}
\tag{B2.31}
\] Here \(\mathcal A_0,\allowbreak{}\mathcal I_0\) are the base active and inactive
sets, \(\mathcal Q=A_{\mathcal Q}\sqcup I_{\mathcal Q}\), and
\(c=c_F\prod_{p\in\mathcal A_0\cup A_{\mathcal Q}}p\). An inactive
marked prime is absent from that branch's conductor and kernel. If a
marked prime belongs to the base local set, its primal term is
identically zero and is removed first: its whole-index mark and the base
zero extension cannot both be nonzero.

The phase at an active prime has the exact form \[
\chi_p(\sigma_p)^{-2}\omega_{p,j}
=\xi_{p,j}\chi_p(c/p)^{2j+2},
\qquad |\xi_{p,j}|=1,
\tag{B2.32}
\] where \(\xi_{p,\allowbreak{}j}\) depends only on \(p,\allowbreak{}j\). This follows by
inserting \(\epsilon_p=-\lambda^{-5}(c/p)^{-2}\) into (B2.11); for
\(j=4\), use \(2j+2\equiv-2\pmod6\). Consequently, two distinct active
primes contribute \[
\chi_p(q)^{2j_p+2}\chi_q(p)^{2j_q+2}
=\left(\frac qp\right)_3^{j_p+j_q+2}.
\tag{B2.33}
\] A residual row prime has \(j_p=1\), and an active mark has \(j_q=0\).
Their phase is therefore one. In contrast, a \(j=4\) prime would leave a
nontrivial cubic phase with the residual row, so every nonresidual base
prime is fixed before the row norm.

Let \(R\) be the residual squarefree row, \(P\) the product of active
marked primes, and \(F_{\rm act}\) the frozen active base product. Fix
separate classes of \(R\) and \(P\) modulo the full \(\mathfrak M^2\).
This fixes \(d,\allowbreak{}\vartheta,\allowbreak{}\kappa_F\) without imposing a pair-dependent
sector restriction. Equation (B2.33) assigns row--row phases to a row
scalar, mark--mark phases to a tuple scalar, and frozen interactions to
the relevant one of these; \(\overline{\alpha(c)}^2\) separates
multiplicatively. The removed primal collisions remain as
\(1_{(P,\allowbreak{}R)=1}\).

After fixing the dual unit, ramified power, and extracted frozen
factors, the remaining dual part is \(nb^3\). Its moving columns are
exactly \[
\chi_R(nb^3)^3=\chi_R(nb)^3,\qquad
\chi_P(nb^3)^{-2}=\chi_P(n)^{-2}1_{(P,b)=1}.
\tag{B2.34}
\] The source coefficient, fixed additive character, and frozen local
columns depend only on the full extracted dual index. A collision with
an extracted frozen prime gives a fixed row, tuple, or dual restriction.
Hence after common smooth separation the tuple coefficient is
independent of \(R,\allowbreak{}n,\allowbreak{}b\), the dual coefficient is independent of
\(R,\allowbreak{}P\), and the remaining mixed masks are exactly those in (B2.29).

\hypertarget{b2.6.-reflected-energy-and-its-two-consequences}{%
\subsection{B2.6. Reflected energy and its two
consequences}\label{b2.6.-reflected-energy-and-its-two-consequences}}

\hypertarget{b2.6.1.-block-parameters-and-estimate}{%
\subsubsection{B2.6.1. Block parameters and
estimate}\label{b2.6.1.-block-parameters-and-estimate}}

Fix \(q_f\asymp Z^V\), \(Q=\log_Zq_{r_\rho}\), \(X=Z^{N_*}\), and rows
\(0<q_k\ll Z^M\). Split the ideal of \(k\) as \[
(k)=k_{\rm pow}k_{\rm sup}R.
\] The maximal powerful part contains every prime of valuation at least
two; \(k_{\rm sup}\) consists of the valuation-one primes in
\(S\cup\operatorname{supp}(fr_\rho)\); and \(R\) is the remaining
squarefree part. Their supports are disjoint. Freeze the actual
\(k_{\rm pow},\allowbreak{}k_{\rm sup}\), the unit, and the fixed sectors above. On
dyads \(q_{k_{\rm pow}}\asymp Z^O\), \(q_R\asymp Z^H\), retain the
actual row cutoff. Necessarily \[
0\le H\le M-O+O(1/\log Z).
\tag{B2.35}
\] It need not equal \(M-O\).

Choose the active zero-exponent base primes and split every \(j=4\)
factor. For a large \(j=4\) term use the disjoint assignment \[
1_{p\mid nb^3}=1_{p\mid n}+1_{p\nmid n}1_{p\mid b}.
\] Let \(A_0\) be the logarithmic norm of all active non-slot primes
outside \(R\); let \(S_0\) count the small \(j=4\) terms and active base
zero masks; and let \(N_0,\allowbreak{}B_0\) count the large terms assigned
respectively to the squarefree and cube factors. Let
\(z_a=\sum_{i\ {\rm active}}z_i\le z\). The active marked product
satisfies \(q_P\asymp Z^{z_a}\). All these parameters are nonnegative.

For actual frozen products \(F_{\rm act},\allowbreak{}N_F,\allowbreak{}B_F\) of logarithmic norms
\(A_0,\allowbreak{}N_0,\allowbreak{}B_0\), write a dual index as \[
\mu=u\lambda^hN_Fn(B_Fb)^3,\quad h\ge-4,\quad
q_n\asymp Z^v,\quad q_b\asymp Z^{\ell_b},\quad
e_\lambda=h\log_Z3,
\] with \(v,\allowbreak{}\ell_b\ge0\). All original source restrictions and permitted
\(S\)-parts remain. Define \[
T_d=2H+2A_0+2z_a-N_*-N_0-3B_0,\qquad
u_h=\min\{v,z_a,(v+z_a)/3\}.
\tag{B2.36}
\] For any fixed \(\tau_{\rm ref}>0\), all dyads outside \[
v+3\ell_b+e_\lambda\le T_d+\tau_{\rm ref}
\tag{B2.37}
\] have arbitrarily small total size. A retained dyad, including the sum
over active tuples, satisfies \[
\begin{aligned}
E_{\rm ref}
={}&\frac O2+\max(H,v+\ell_b)-S_0-B_0+z_a-u_h-\ell_b-\frac23e_\lambda\\
&-\frac12(T_d-v-3\ell_b-e_\lambda)_+.
\end{aligned}
\tag{B2.38}
\] Precisely, for the fixed \(k_{\rm pow},\allowbreak{}k_{\rm sup}\) its energy over
\(R\) is \(\ll Z^{E_{\rm ref}-O/2+\epsilon}\); summing these fixed parts
in their \(O\)-dyad and all local choices gives
\(\ll Z^{E_{\rm ref}+\epsilon}\). Constants are uniform under the scope
in §B2.1. Current row restrictions must be independent of the active
tuple and dual indices apart from \(1_{(P,\allowbreak{}R)=1}\); the character,
puncture, and original product mark are also independent of a
subsequently averaged \(f\). The thresholds depend only on the bounded
log-length ranges, fixed profiles and moduli,
\(\epsilon,\allowbreak{}\tau_{\rm ref}\). Negative \(e_\lambda\) has size
\(O(1/\log Z)\).

\textbf{Proof.} An inactive slot has absolute mass
\(\sum_{p\in\mathcal P_i}q_p^{-1}\ll Z^\epsilon\), including bounded
lists. Triangle inequality in the row Hilbert space therefore freezes
those primes at this cost; their collision exclusions become fixed row
restrictions. Keep all active tuples inside the norm.

Equation (B2.7), before the local factors, gives \[
\frac{|d(\mu)|}{\sqrt{q_\mu}}
\le27\,q_{\lambda^h}^{-1/3}q_{N_Fn}^{-1/2}q_{B_Fb}^{-1}.
\] A small \(j=4\) or active base zero term contributes \(q_p^{-1/2}\).
A large \(j=4\) term contributes \(q_p^{1/2}\): assigning \(p\) to
\(N_F\) cancels that square-root denominator; assigning it to \(B_F\)
leaves \(q_p^{-1/2}\). If \(N_F\) meets \(b\), its entire cube
occurrence remains in \(b\). Thus the common coefficient, per active
tuple, has exponent \[
-v/2-\ell_b-e_\lambda/3-(S_0+B_0+z_a)/2.
\tag{B2.39}
\]

The denominator is \(c=c_FF_{\rm act}RP\). With actual annular
coordinates \[
y_R=q_R/Z^H,\quad y_n=q_n/Z^v,\quad
y_b=q_b/Z^{\ell_b},\quad y_i=q_{p_i}/Z^{z_i},
\] the kernel argument is exactly \[
\frac{q_\mu X}{q_c^2}
=Y\,y_ny_b^3y_R^{-2}\prod_{i\ {\rm active}}y_i^{-2},
\qquad
Y=q_{c_F}^{-2}Z^{v+3\ell_b+e_\lambda-T_d}.
\tag{B2.40}
\] This identity remains true at shared prime factors. On the original
annuli the monomial is bounded above and below by fixed constants.
Equation (B2.22) therefore permits the scalar \[
m_{\rm ker}=Z^{-(T_d-v-3\ell_b-e_\lambda)_+/4}
\] to be factored from the whole linear row vector, with bounded
normalized seminorms.

Remove tails before this Fourier separation. For sufficiently large
\(Z\), a dyad violating (B2.37) has actual argument
\(>Z^{\tau_{\rm ref}/2}\). The full source satisfies
\(|d(\mu)|/\sqrt{q_\mu}\le27\,\allowbreak{}3^{4/3}\), and every local column factor
is at most \(q_p^{1/2}\). Row and tuple counts, branch counts, and
\(1+q_c^2/X\) cost a fixed power \(Z^{B_{\rm tail}}\), chosen
independently of the kernel order. Apply (B2.23) with \(a=X/q_c^2\),
\(U=Z^{\tau_{\rm ref}/2}\) and \[
A>1+2(B_{\rm tail}+D)/\tau_{\rm ref}.
\] The entire discarded sum is \(O(Z^{-D})\) times a finite test
seminorm, with its corresponding fixed polynomial height cost. This
counts all discarded \(\mu\).

For a retained block, apply the single common separation to (B2.40),
then enlarge only the separated positive row norm to \(q_R\ll Z^H\). By
§B2.5 the separated arithmetic expression is the one in (B2.29), with \[
K=Z^H,\qquad N=Z^v,\qquad B=Z^{\ell_b},\qquad L=Z^{z_a}.
\] Leave \(q_P^{-1/2}\) inside that norm. The outside squared
coefficient is \(Z^{-v-2\ell_b-2e_\lambda/3-S_0-B_0+\epsilon}\). Since
\[
\max\{v,z_a,2(v+z_a)/3\}=v+z_a-u_h,
\] (B2.29) gives the first line of (B2.38), without \(O/2\).
Multiplication by \(m_{\rm ker}^2\) supplies its last line.

Every powerful ideal is uniquely \(x^2y^3\) with \(y\) squarefree. Ideal
counting and \(\sum_yq_y^{-3/2}<\infty\) give \(O(Z^{O/2})\) choices in
its dyad. The supported part divides the radical of \(Sfr_\rho\) and has
\(O(Z^\epsilon)\) choices; the local branches have the same divisor
bound. Rowwise divisor Cauchy and the positive sum over frozen parts add
\(O/2+\epsilon\) once. Active tuples were already bounded by the hybrid
norm and are not counted again. This proves (B2.38).

\hypertarget{b2.6.2.-general-label-and-puncture-width}{%
\subsubsection{B2.6.2. General label and puncture
width}\label{b2.6.2.-general-label-and-puncture-width}}

At an active non-slot prime, the coefficient of its log-norm in
\(2A_0-N_0-S_0-4B_0\) is

\begin{longtable}[]{@{}lr@{}}
\toprule\noalign{}
Local choice & Coefficient \\
\midrule\noalign{}
\endhead
\bottomrule\noalign{}
\endlastfoot
\(j\) odd or \(j=2\) & \(2\) \\
\(j=0\), active & \(1\) \\
\(j=4\), small or assigned to \(n\) & \(1\) \\
\(j=4\), assigned to \(b\) & \(-2\) \\
\(j=0\), inactive & \(0\) \\
\end{longtable}

A residual row prime contributes two. These coefficients are bounded
prime by prime by the valuations in
\((q_k^2/q_{k_{\rm pow}})q_fq_{r_\rho}\): a powerful row valuation
\(e\ge2\) supplies \(e\), a simple row prime supplies two, and a prime
occurring only in \(f\) or \(r_\rho\) supplies one. Intersections only
increase the available valuation. Hence \[
q_R^2 Z^{2A_0-N_0-S_0-4B_0}
\le C_S\,\frac{q_k^2}{q_{k_{\rm pow}}}\,q_fq_{r_\rho},
\] and therefore \[
T_d-S_0-B_0
\le2M-O+Q+V-N_*+2z_a+O(1/\log Z).
\tag{B2.41}
\] This is the general width bound used when completion shortens the
column length to \(N_*=N-3\ell_h\). It is uniform in the actual \(f\)
and puncture, without replacing \(H\) by \(M-O\).

With no slots, \(z_a=u_h=0\). On a retained dyad, (B2.38) gives \[
E_{\rm ref}\le O/2+\max\{H,T_d-S_0-B_0\}
+O(1/\log Z)+\tau_{\rm ref}.
\] For the row branch drop the nonpositive terms; for the column branch
use \(v\le T_d-3\ell_b-e_\lambda+\tau_{\rm ref}\). Apply
(B2.35),(B2.41), take \(\tau_{\rm ref}\) and the small-power losses
sufficiently small, and sum the logarithmically many blocks: \[
\sum_{\substack{0<q_k\ll Z^M\\q_{k_{\rm pow}}\asymp Z^O}}
|\mathcal T_1(Z^{N_*};k;V)|^2
\ll_\epsilon
Z^{\,O/2+\max\{M-O,\;2M-O+Q+V-N_*\}+\epsilon}.
\tag{B2.42}
\] In particular, with \(f=r_\rho=1\), summing \(O\) gives the unmarked
completed bound \[
\sum_{0<q_k\ll Z^M}|\mathcal T_1(Z^{N_*};k;V)|^2
\ll_\epsilon Z^{\max(M,\,2M-N_*)+\epsilon}.
\tag{B2.43}
\] This includes the balanced energy \(Z^{1+\epsilon}\) at \(M=N_*=1\).

\hypertarget{b2.6.3.-short-marked-energy-without-label-or-puncture}{%
\subsubsection{B2.6.3. Short marked energy without label or
puncture}\label{b2.6.3.-short-marked-energy-without-label-or-puncture}}

When \(f=r_\rho=1\), every active non-slot prime outside \(R\) divides
the powerful part. Thus \[
2A_0\le O+O(1/\log Z),\qquad 0\le N_0\le A_0.
\] If \(M+2z\le N_*\), then for every \(\epsilon>0\), \[
\sum_{0<q_k\ll Z^M}|\mathcal T_d(Z^{N_*};k;V)|^2
\ll_\epsilon
Z^{\,M+(N_*+2z-2M)_+/4+\epsilon}.
\tag{B2.44}
\] To see this directly from (B2.38), first ignore fixed \(O(1/\log Z)\)
errors and choose \(\tau_{\rm ref}\) arbitrarily small. Since \[
T_d-H\le M+2z-N_*\le0,
\] retention forces \(\max(H,\allowbreak{}v+\ell_b)=H+O(\tau_{\rm ref}+1/\log Z)\).
If \(u_h=z_a\), (B2.38) is at most \(M\), up to those errors. Otherwise
\(u_h\ge v/2\). Put \(y_d=v+3\ell_b+e_\lambda\); because \(\ell_b\ge0\)
and \(e_\lambda\ge-O(1/\log Z)\), \[
u_h+\ell_b+\frac23e_\lambda
+\frac12(T_d-y_d)_+
\ge \frac14y_d+\frac12(T_d-y_d)_+-O(1/\log Z)
\ge\frac14T_d-O(1/\log Z).
\] With \(\Delta_H=M-O-H\ge-O(1/\log Z)\), substitution into (B2.38)
gives \[
\begin{aligned}
4(E_{\rm ref}-M)
&\le N_*+2z_a-2M-2\Delta_H
-(2A_0-N_0+B_0+4S_0)+O(\tau_{\rm ref}+1/\log Z)\\
&\le N_*+2z-2M+O(\tau_{\rm ref}+1/\log Z).
\end{aligned}
\] Restoring the row branch, tails, and small-power losses proves
(B2.44). No moment estimate or zero-free conclusion enters this proof.

\hypertarget{appendix-b3.-the-inverse-and-plain-moment-estimates}{%
\section{Appendix B3. The inverse and plain moment
estimates}\label{appendix-b3.-the-inverse-and-plain-moment-estimates}}

This appendix proves UM and MI, the unconditional zero-slot moment, and
the canonical common-profile positive-slot moment in the range stated in
B3.4. Its arithmetic inputs are the finite Gauss identities
(B1.1)--(B1.19) and the completed reflected-energy estimate
(B2.35)--(B2.41), proved in the preceding appendices. In particular, the
inverse-moment induction terminates at RE; the plain-moment induction is
independent of the inverse moment.

\hypertarget{b3.1.-common-conventions-and-analytic-bookkeeping}{%
\subsection{B3.1. Common conventions and analytic
bookkeeping}\label{b3.1.-common-conventions-and-analytic-bookkeeping}}

Write \(q_a=Na\), sum ideals through their primary generators, and
retain every zero of every displayed residue symbol, including powers
divisible by six. Element rows include units and arbitrary prime
valuations. All column ideals avoid the fixed set \(S\). A fixed ray
character has its entire defining modulus supported in \(S\); a moving
coprimality condition is displayed separately. Put \[
a_0(n)=\overline{\alpha(n)}\gamma_2(n)\qquad(n\text{ squarefree}).
\] For disjoint prime lists \(\mathcal P_i\), fixed throughout the
current row and label averages, define a \textbf{product mark} \[
d_I(A)=\sum_{(p_i)\in\prod_{i\in I}\mathcal P_i}
       \prod_{i\in I}a_i(p_i)\mathbf1_{p_i\mid A},
\qquad q_{p_i}\asymp Z^{z_i},\quad |a_i(p_i)|\le1.
\tag{B3.1}
\] The number of slots is fixed; their coefficients need not be
multiplicative. Smooth slot weights are included in the individual
\(a_i\). The useful property is the product form, not just the
consequent divisor bound. In particular, \[
\mathbf1_{p\mid D_1\cdots D_h n}
=\sum_{j=1}^h\mathbf1_{p\mid D_j}\prod_{a<j}\mathbf1_{p\nmid D_a}
 +\mathbf1_{p\mid n}\prod_{a\le h}\mathbf1_{p\nmid D_a}
\tag{B3.2}
\] assigns each slot first to one extracted factor or leaves it on the
remaining column. Assigned coefficients are divisor bounded; surviving
coefficients retain their original product form. Each copy of a squared
mark is allocated independently.

\textbf{Common-profile lemma.} Suppose finitely many operations use
smooth functions on fixed compact intervals of positive normalized
norms, radial Schwartz Fourier kernels, and finite Fourier expansions on
fixed ray groups. All arithmetic restrictions that are not smooth remain
in the outer weights. For a coupled profile \(H(y_1,\allowbreak{}\ldots,\allowbreak{}y_d)\),
multiply by larger individual cutoffs equal to one on the full current
supports and use \[
H(y)=(2\pi)^{-d}\int_{\mathbf R^d}\widehat H(t)\prod_jy_j^{it_j}\,dt,
\qquad
\int|\widehat H(t)|\langle t\rangle^Jdt
 \ll_{J,d}\max_{|\beta|\le2r}\|D^\beta H\|_\infty,
\quad 2r>J+d,
\tag{B3.3}
\] where \(D_j=y_j\partial_{y_j}\) and \(\langle t\rangle=1+|t|\). The
coefficient measure is defined on the ambient box before any current row
or live label is chosen. Relations among those norms restrict evaluation
points, not the measure. Real inverse roots belong to this coupled
profile once; the final column tests receive only their individual
cutoffs and Fourier powers.

The estimate follows by integration by parts with \((1-\Delta)^r\) in
logarithmic coordinates. If a kernel is \(K(A\prod_jy_j^{e_j})\), its
Euler derivatives are fixed multiples of \((x\partial_x)^rK(x)\),
uniformly for every \(A>0\). Consequently arbitrarily high fixed
derivative orders do not introduce a power of \(Z\). Norm twists
introduce only fixed polynomials in their heights. For any \(Y\ge1\),
radial decay and lattice counting also give \[
\sum_{h\ne0:\,a q_h>Y}|K(aq_h)|
 \ll_B(1+a^{-1})Y^{1-B}\qquad(a>0,\ B>1).
\tag{B3.4}
\] Indeed the shell \(2^jY<aq_h\le2^{j+1}Y\) has \(O(1+2^jY/a)\) points.
This proves the displayed bound by two geometric series.

Here and below a tail is discarded on the genuine arithmetic expression
before extending a coprime sum or taking an absolute Fourier integral.
To do so, enclose the retained frequencies in an outer ball, fixed for
the current norm sector, whose complement forces the actual kernel
argument above \(Z^\tau\). Keep the full kernel inside that ball. If the
remaining raw factors and \(1+a^{-1}\) cost \(Z^{B_0}\), choose
\(B>1+(B_0+T)/\tau\) in (B3.4) to obtain \(O(Z^{-T})\). All raw scales
below are polynomially bounded, so \(B_0\) is fixed before \(B\) is
chosen.

For completeness, the finite-order uniformity used in both inductions
follows from this rule. At a terminal vertex choose its required profile
and height orders. Along an edge, (B3.3) integrates the polynomial
demanded by its child after increasing finitely many parent derivative
orders. Work backwards through a bounded-depth acyclic graph. For a
discrete family with weighted count \(\ll Z^b\), extract \(Z^b\) once
and normalize its counting measure before this argument. Cauchy--Schwarz
averages two child exponents; it does not add them. Minkowski takes
their maximum after the coefficient mass is accounted for. Thus the path
length bounds the total numerical exponent loss, while all derivative
orders remain finite. Fresh quotient cutoffs, larger separation cutoffs
and bounded scale-clipping families are fixed through this finite depth
before the final threshold in \(Z\). A nonempty window at a nominal
negative length \(n\) has \(1\le Z^{-n}\le B\), where \(B\) is its fixed
upper endpoint; replacing it by length zero is therefore a bounded
dilation, not a growing support enlargement.

\textbf{Masked primitive Poisson.} If \(\psi\) is primitive modulo
\(\mathfrak m\), zero on nonunits, and \(\Phi\) is radial Schwartz, then
\[
\sum_{k\in\mathcal O}\psi(k)\mathbf1_{(k,R)=1}\Phi(q_k/K)
=\frac{K\gamma(\psi;\mathfrak m)}{\sqrt{q_{\mathfrak m}}}
 \sum_{d\mid\operatorname{rad}R}\frac{\mu(d)\psi(d)}{q_d}
 \sum_h\overline{\psi(h)}\,\mathcal F\Phi
       \left(\frac{Kq_h}{q_dq_{\mathfrak m}}\right).
\tag{B3.5}
\] Insert the full Möbius expansion of the mask, write \(k=dk'\), and
apply lattice Poisson in residue classes modulo \(\mathfrak m\). The
finite transform is
\(\sqrt{q_{\mathfrak m}}\gamma(\psi;\mathfrak m)\overline\psi(h)\); it
vanishes at nonunits by primitivity. Finite Parseval gives
\(|\gamma|=1\). For the primitive principal character use
\(\mathfrak m=1\), \(\gamma=1\), including \(h=0\). Its nonzero
frequencies have absolute sum \(O_\Phi(d(\operatorname{rad}R))\), since
\(A\sum_{h\ne0}|\mathcal F\Phi(Aq_h)|\ll1\). Thus principal nonzero
terms may be restored with the same retained outer mask at their
original diagonal cost. Fixed ray restrictions are first expanded into
characters; deleted conductor zeros are kept in \(R\).

\hypertarget{b3.2.-a-canonical-marked-energy}{%
\subsection{B3.2. A canonical marked
energy}\label{b3.2.-a-canonical-marked-energy}}

Fix a squarefree puncture \(\rho(n)=\mathbf1_{(n,\allowbreak{}r_\rho)=1}\),
independent of the current row \(k\) and squarefree label \(f\), and put
\(Q=\log_Zq_{r_\rho}\). Let \(q_f\asymp Z^V\),
\(0\le w(f)\le C d(f)^C\), with \(w\) depending only on \(f\). For
\(F=N+V\), define \[
\mathcal E(N,V,M)=Z^{-V}\sum_f w(f)\sum_{0<q_k\ll Z^M}
 \left|Z^{-N/2}\sum_{n\ {\rm sf}}a_0(n)\nu(n)\rho(n)
       \chi_n(k)\chi_n(f)^4d_I(n)W(q_n/Z^N)\right|^2.
\tag{B3.6}
\] Here \(\nu\) is a fixed finite ray character, all real lengths range
over fixed bounded sets, the slot sum is at most a fixed cap \(z_0\),
and no further coefficient on \(n\) is permitted.

\textbf{Canonical estimate.} For any fixed \(c_*>0\), \[
F-M-Q-z_0\ge c_*,\qquad4F-3M-6z_0\ge c_*
\quad\Longrightarrow\quad
\mathcal E(N,V,M)\ll Z^{F+\epsilon}.
\tag{B3.7}
\] The bound is uniform over the moving punctures and all label-only
weights in the displayed class. It holds simultaneously for every
surviving slot subcollection with the original cap \(z_0\).

\textbf{Proof: terminal completion.} Choose \(0<d\le c_*/200\). Put
\(\Psi_k(n)=\nu(n)\rho(n)\chi_n(k)\chi_n(f)^4\), multiplicatively on all
ideals. For \(V_*(y)=y^{1/2}W(y)\) and a mark function \(d_*\), write
the completed marked sum as \[
T_{d_*}(X;k)=\sum_{n\ {\rm sf},b}
 \frac{a_0(n)\Psi_k(n)\overline{\alpha(b)}^3\Psi_k(b)^3}
      {\sqrt{q_n}\,q_b}\,d_*(nb^3)V_*(q_nq_b^3/X).
\] The exact cube inversion is \[
Z^{-N/2}\sum_{n\ {\rm sf}}a_0(n)\Psi_k(n)d_I(n)W(q_n/Z^N)
=\sum_h\frac{\mu(h)\overline{\alpha(h)}^3\Psi_k(h)^3}{q_h}
 T_{d_h}(Z^N/q_h^3;k),\quad d_h(A)=d_I(h^3A).
\tag{B3.8}
\] After expansion put \(c=hb\). The mark is \(d_I(nc^3)\), independent
of the divisor \(h\), and its coefficient contains
\(\sum_{h\mid c}\mu(h)\); only \(c=1\) remains. This also verifies the
central normalization and every mask.

Use fixed dyadic windows, with the unit window centered at zero. If
\(V<d\) and the center \(\ell_1\) of \(h\) is less than \(d\), freeze
\(h\), assign its slots by (B3.2), and apply RE at length
\(N_*=N-3\log_Zq_h\), with surviving cap \(z_a\le z_0\). The coefficient
\(\Psi_k(h)^3\) is a contraction in the joint \((f,\allowbreak{}k)\) Hilbert space,
and \(\sum_{q_h\asymp Z^{\ell_1}}q_h^{-1}\ll1\). To record the terminal
exponent check, let \(c\) be the current margin in (B3.7), and let
\(\eta\) bound fixed-window logarithmic errors. The RE support
inequalities give \[
H\le M-O+\eta,\qquad
T_d-S_0-B_0\le M-O+2z_a-z_0-c+5d+4\eta.
\tag{B3.9}
\] The second inequality follows from
\(T_d-S_0-B_0\le2M-O+Q+V-N_*+2z_a+\eta\): use the first margin and
\(2V+3\log_Zq_h<5d+3\eta\).

On retained reflected blocks \(v+3b+e_\lambda\le T_d+\tau_{\rm ref}\)
and \(e_\lambda\ge-\eta\). In the row branch of RE, drop its nonpositive
terms to obtain \(E\le M+z_a+5\eta/3\le F-c+5\eta/3\). In the column
branch, \(u=z_a\) gives \(E\le O/2+T_d-S_0-B_0+5\eta/3+\tau_{\rm ref}\);
if \(u\ne z_a\), its definition gives \(u\ge v/2\) and
\(E\le O/2+(T_d-S_0-B_0)/2+z_a+7\eta/6+\tau_{\rm ref}/2\). Substitution
of (B3.9), followed respectively by the first and second margins, yields
\[
E\le F-2c+5d+17\eta/3+\tau_{\rm ref},\qquad
E\le F-M/4-3c/4+5d/2+19\eta/6+\tau_{\rm ref}/2.
\tag{B3.10}
\] All are strictly below \(F\) when \(c\ge c_*/2\) and
\(\eta,\allowbreak{}\tau_{\rm ref}\), together with the reflected divisor loss, are
at most \(c_*/1000\). Summing the fixed-label RE bound costs
\(Z^{V+o(1)}\), canceled by \(Z^{-V}\).

\textbf{Nonterminal block.} If \(V<d\) and \(\ell_1\ge d\), reopen the
completion in (B3.8), combine its cube variable with \(h\), and call the
result \(b\). If its center is \(\ell\), put \(r=N-3\ell\). If
\(V\ge d\), instead take \(b=1,\allowbreak{}\ell=0,\allowbreak{}r=N\). Triangle inequality in the
joint Hilbert space and the common-profile lemma reduce every remaining
term to \[
\begin{gathered}
Z^{-r-2\ell-V}\sum_f w(f)\sum_k\Phi(q_k/Z^M)\\
\times\left|\begin{gathered}\sum_{q_b\asymp Z^\ell}\beta(b,f)\chi_b(k)^3\\
\times\sum_{n\ {\rm sf}}a_0(n)\nu(n)\rho(n)\chi_n(k)\chi_n(f)^4
        d_I(nb^3)W(q_n/Z^r)\end{gathered}\right|^2,
\end{gathered}
\tag{B3.11}
\] where \(V+\ell\ge d\), \(\beta\) is divisor bounded and independent
of \(k,\allowbreak{}n\), and includes the original \((b,\allowbreak{}f)=1\) mask. The variables
\(b,\allowbreak{}n\) need not be coprime. We now reduce (B3.11) to (B3.6) with
strictly shorter rows.

\textbf{First Poisson transformation.} In the square write \[
\mathfrak B=\operatorname{rad}(b_1b_2),\qquad
n_i=A_iCu_i,\quad A_i=(n_i,\mathfrak B),\quad
C=(n_1/A_1,n_2/A_2).
\] Then \(u_1,\allowbreak{}u_2\) are individually squarefree, coprime to one another
and to \(C\mathfrak B\). At \(p\mid\mathfrak B\), set
\(a_i=\mathbf1_{p\mid A_i}\), \(\pi=v_p(b_1b_2)\bmod2\), and
\(t_p=a_1-a_2+3\pi\bmod6\). Put \(R_1=\prod_{t_p\ne0}p\). The primitive
row character and its residual zero mask are \[
\chi_{u_1}\overline{\chi_{u_2}}\prod_{t_p\ne0}\chi_p^{t_p},
\qquad\mathfrak m_1=u_1u_2R_1,
\qquad R_{\rm mask}=C\mathfrak B/R_1.
\tag{B3.12}
\] Apply (B3.5), introducing \(d_k\mid C\mathfrak B/R_1\) and frequency
\(h\). Write \(A_i,\allowbreak{}B,\allowbreak{}R,\allowbreak{}\delta\) also for the centers of
\(A_i,\allowbreak{}C,\allowbreak{}R_1,\allowbreak{}d_k\). The actual conductor length is
\(\log_Z(q_{n_1}q_{n_2}/(q_{A_1}q_{A_2}q_C^2))+\log_Zq_{R_1}\). The
principal case forces \(u_1=u_2=1\), \(\pi=0\), and \(a_1=a_2\) at every
prime, hence \(n_1=n_2\) and \(b_1b_2\) a square. Even counting all cube
pairs, its normalized contribution is \(O(Z^{M+o(1)})\). The same bound
covers restoring principal nonzero frequencies with the mask below.

CRT, \(\gamma_1(u)\gamma_2(u)=\mu(u)\alpha(u)G(u)\), and reciprocity
change the two residual columns to inverse coefficients. Their remaining
joint ray factor is
\(G([u_1u_2^{-1}])R([u_1u_2^{-1}],\allowbreak{}[\prod p^{t_p}])\); expand it on the
fixed ray group before extending coprimality. Put
\(E=\prod_{\pi=1}p^{a_1+a_2}\), \(y=hf^2E\). The residual exponent at
each prime of \(\mathfrak B\) is
\(4a_i+\mathbf1_{t_p\ne0}(1\pm t_p)+\pi(a_1+a_2)\bmod6\). The following
exhaustive table gives the same exponent on both sides:

\begin{longtable}[]{@{}rrrrr@{}}
\toprule\noalign{}
\(\pi\) & \(a_1\) & \(a_2\) & \(t_p\) & Residual exponent \\
\midrule\noalign{}
\endhead
\bottomrule\noalign{}
\endlastfoot
0 & 0 & 0 & 0 & 0 \\
0 & 1 & 0 & 1 & 0 \\
0 & 0 & 1 & 5 & 0 \\
0 & 1 & 1 & 0 & 4 \\
1 & 0 & 0 & 3 & 4 \\
1 & 1 & 0 & 4 & 4 \\
1 & 0 & 1 & 2 & 4 \\
1 & 1 & 1 & 3 & 4 \\
\end{longtable}

An exponent zero still means a puncture. Write \(b_1b_2=q^2s\), \(s\)
squarefree, and put \[
J_2=\prod_{\pi=0,\ a_1=a_2=1}p,
\quad J=sJ_2,\quad q=J_2q_0.
\] Then \(J\) is squarefree and the residual factor is exactly
\(\chi_n(J)^4\mathbf1_{(n,\allowbreak{}\operatorname{rad}q_0)=1}\). If \(j\) centers
\(J\), the actual identities \[
\log_Zq_{q_0}=\tfrac12\log_Z(q_{b_1}q_{b_2})
 -\tfrac12\log_Zq_s-\log_Zq_{J_2}\ge0,
\qquad
\log_Z\frac{q_{R_1}q_E}{q_{A_1}q_{A_2}}
=\log_Zq_s-2\log_Zq_{J_2}
\tag{B3.13}
\] give the retained row enclosure \[
H_c=2r-2B-M+4\ell+2V+\delta-j,
\qquad H_+=H_c+12\eta+\tau,
\qquad 0<q_y\le Z^{H_+}.
\tag{B3.14}
\] Indeed the actual first kernel argument is
\(Z^Mq_h/(q_{d_k}q_{\mathfrak m_1})\). If it is at most \(Z^\tau\),
(B3.13) bounds \(\log_Zq_y\) by \(H_c+12\eta+\tau\). Discard only the
complement of this outer ball by (B3.4). A negative \(H_+\) leaves no
nonzero row.

Now define the quotient-free expression on all individually allowed
\(u_1,\allowbreak{}u_2\), keeping its full kernel, formal product conductor and outer
ball, and insert the complete identity
\(\mathbf1_{(u_1,\allowbreak{}u_2)=1}=\sum_{t'\mid(u_1,\allowbreak{}u_2)}\mu(t')\). Write
\(u_i=t'x_i\) and let \(t\) center \(t'\). No Poisson formula is
asserted on noncoprime conductors: the full Möbius sum annihilates
precisely those artificial terms. Squarefreeness retains \((x_i,\allowbreak{}t')=1\).
With \[
s_i=r-A_i-B-t,\qquad
\varkappa_i=M-2r-2\ell-V-\delta+A_i+B-R/2,
\tag{B3.15}
\] the separated inverse polynomial on side \(i\) is \[
P_i(y)=\sum_{x\ {\rm sf}}\mu(x)\nu_i(x)\rho(x)
 \overline{\chi_x(y)}\chi_x(CJ)^4\chi_x(d_k)
 \mathbf1_{(x,\operatorname{rad}(q_0t'))=1}d_{I_i}(x)W_i(q_x/Z^{s_i}).
\tag{B3.16}
\] Here \(\nu_i\) is fixed in its ray summand. The fourth-power factor
at \(f\) became \(\overline{\chi_x(f^2)}\), which explains both the row
\(y\) and its conjugation. Priority allocation to \((b_i,\allowbreak{}A_i,\allowbreak{}C,\allowbreak{}t';x_i)\)
leaves an original slot subcollection: all earlier survival exclusions
are already supplied by \(CJ\), \(q_0\), and \(t'\).

For precision, the first coupled profile uses the independent normalized
norms of \((A_1,\allowbreak{}A_2,\allowbreak{}C,\allowbreak{}d_k,\allowbreak{}R_1,\allowbreak{}t',\allowbreak{}h,\allowbreak{}x_1,\allowbreak{}x_2)\). Apart from its larger
cutoffs, its real root and kernel are \[
\begin{gathered}
y_d^{-1}y_R^{-1/2}y_t^{-1}(y_1y_2)^{-1/2}
 W(y_{A_1}y_Cy_ty_1)\overline{W(y_{A_2}y_Cy_ty_2)}\\
\times K_1\!\left(A_1^{\rm ker}\frac{y_h}{y_dy_Ry_t^2y_1y_2}\right),
\\ A_1^{\rm ker}=Z^{M+h_1-\delta-R-2t-s_1-s_2},
\end{gathered}
\tag{B3.17}
\] where \(h_1\) centers the bare frequency. Multiply this profile by
\(Z^{-9\eta/2}\); its outside side factors are
\(Z^{(\varkappa_i+9\eta/2)/2}\). This is an exact separation of the full
prefactor and kernel, with full fresh \(x_i\)-supports satisfying
\(|\log_Zq_{x_i}-s_i|\le4\eta\). The common-profile lemma applies
uniformly to \(A_1^{\rm ker}>0\).

Apply weighted Cauchy--Schwarz over the \textbf{whole} outer tuple
before bounding its weights. At fixed old reconstruction and \(y\ne0\),
\[
\sum_{f^2E\mid(y)}w(f)\ll d((y))^{C+1}\ll Z^{o(1)}.
\tag{B3.18}
\] Thus \(f,\allowbreak{}h\) disappear together, since \(h=y/(f^2E)\), while the
original \(-V\) remains in \(\varkappa_i\). Only at this positive-norm
stage replace the retained ball by a fixed smooth majorant at scale
\(Z^{H_+}\). Keep \(J\) averaged and \(q_0\) fixed. From the table, if
\(r_{\rm diff}\) is the actual radical length of even-parity primes with
\(a_1\ne a_2\), \(\log_Zq_{R_1}=\log_Zq_s+r_{\rm diff}\). Hence \[
\log_Zq_{q_0}\le\ell+R/2-j+5\eta/2.
\tag{B3.19}
\] This bounds the count of \(q_0\); counting \(J\) as well costs
\(Z^{\ell+R/2+7\eta/2+o(1)}\). Its pieces are not counted independently
later.

\textbf{Second Poisson transformation.} Fix a first Cauchy side. Expand
\(|P_i(y)|^2\), write \(x_a=g'z_a\), \(g'=(x_1,\allowbreak{}x_2)\), and center \(g'\)
at \(g\). Apply (B3.5) with primitive character
\(\overline{\chi_{z_1}}\chi_{z_2}\), modulus \(z_1z_2\), residual mask
\(g'\), divisor \(d_2\mid g'\) of center \(\theta\), and frequency
\(k''\). Put \[
\begin{gathered}
\lambda_c=H_c-\theta-(s_i-g),\\
M_c=2(s_i-g)-H_c+\delta+2\theta\\
=M-4\ell-2A_i-2t-2g+2\theta-2V+j,
\\ M'=M_c+\eta.
\end{gathered}
\tag{B3.20}
\] Retain \(0<q_{d_kd_2k''}\le Z^{M'}\). The actual conductor differs
from \(2(s_i-g)\) by at most \(10\eta\). The two appearances of \(d_2\)
and one of \(d_k\) add \(3\eta\); subtracting \(H_+=H_c+12\eta+\tau\)
cancels the same tail tolerance \(\tau\). Thus its complement is an
actual-ratio tail, exactly as required by (B3.4).

The principal case \(z_1=z_2=1\) has nominal total exponent \[
\varkappa_i+H_c+s_i+B+t+\ell+R/2=F-B-j.
\tag{B3.21}
\] The first prefactor, row enclosure, diagonal column window and outer
counts add at most \(26\eta+\tau+o(1)\). It is therefore within the
target bound, including any principal nonzero terms restored with the
identical second mask.

The Gauss signal is
\(\mu(z)\gamma_{-1}(z)=a_0(z)\chi_z(-1)\overline{G(z)}\). After the
second side is placed inside its conjugate, the relative ray factor is
\(G([z_2z_1^{-1}])\). It is unchanged on removing a common factor.
Extend the separate columns and full kernel off coprimality and insert
the full Möbius sum \(z_a=v'n_a\), with \(v'\) squarefree of center
\(v\). The common extracted factor and the old gcd restriction retain
\((v',\allowbreak{}g'CJ)=1\), \((g',\allowbreak{}CJ)=1\), and the frequency zero \((v',\allowbreak{}k'')=1\).
Set \[
\begin{gathered}
\gamma=(q_0,t',r_g),\\ r_g=g'/d_2,\\
k_{\rm new}=d_kd_2k'',\\ f_{\rm new}=JCd_2v',
\\\rho_\gamma(n)=\rho(n)\mathbf1_{(n,\operatorname{rad}(q_0t'r_g))=1}.
\end{gathered}
\tag{B3.22}
\] The retained common zeros make \(J,\allowbreak{}C,\allowbreak{}d_2,\allowbreak{}v'\) pairwise coprime and
squarefree. The new label is therefore squarefree, and \[
\chi_n(d_k)\chi_n(k'')\overline{\chi_n(d_2)}
 \chi_n(CJ)^4\chi_n(v')^4
=\chi_n(k_{\rm new})\chi_n(f_{\rm new})^4.
\tag{B3.23}
\] This uses \(\overline{\chi_n(d_2)}=\chi_n(d_2)\chi_n(d_2)^4\), valid
also at nonunits, and \(a_0(v'n)=a_0(v')a_0(n)\chi_n(v')^4\). Expanding
\(G([n_2n_1^{-1}])\) yields the same fixed character \(\nu'\) on both
child sides. Priority allocation leaves original slot subcollections;
their old survival zeros are supplied by \(f_{\rm new}\) and
\(\rho_\gamma\). Freeze only \(\gamma\); the four factors of
\(f_{\rm new}\) remain inside its average.

The formal child lengths and their actual enclosing windows are \[
\begin{gathered}
N_c=r-A_i-B-t-g-v,\\ V_c=B+\theta+v+j,\\ F_c=N_c+V_c,
\\ N'=\max(0,N_c),\\ F'=N'+V_c,
\\0\le\delta_N=N'-N_c\le6\eta.
\end{gathered}
\tag{B3.24}
\] The full fresh column window has logarithmic error \(6\eta\), the
full label window \(4\eta\). An empty window is omitted; otherwise the
clipping is a bounded dilation as in the common-profile lemma. With
\(D_c=\ell+A_i+t+g-\theta+V\), direct substitution gives \[
\begin{aligned}
F_c&=F-D_c-2\ell+j,&M_c&=M-2D_c-2\ell+j,\\
F_c-M_c&=F-M+D_c,&
4F_c-3M_c&=4F-3M+2(A_i+t+g-\theta+V)+j.
\end{aligned}\tag{B3.25}
\] The fixed triples range over the projection of the genuine outer
reconstruction set, before positive row or label enlargement. Every such
triple has a source witness, giving \(g-\theta\ge-2\eta\),
\(j\le2\ell+3\eta\), and \(\log_Zq_{q_0}\le\ell+\eta\). Consequently its
actual puncture satisfies
\(Q'\le Q+\ell+t+g-\theta+4\eta\le Q+D_c+4\eta\). Therefore \[
F'-M'-Q'-z_0\ge c-5\eta,\quad
4F'-3M'-6z_0\ge c-7\eta,\quad
M-M'\ge2(V+\ell)-8\eta,\quad F'\le F+11\eta.
\tag{B3.26}
\] These inequalities hold for the structural triples themselves, not
for every point in a containing count ball.

To finish the normalization, the second common profile uses normalized
norms of \((g',\allowbreak{}d_2,\allowbreak{}v',\allowbreak{}k'',\allowbreak{}n_1,\allowbreak{}n_2)\), with real root
\(y_d^{-1}y_v^{-1}(y_1y_2)^{-1/2}\), old windows
\(W_i(y_gy_vy_1)\overline{W_i(y_gy_vy_2)}\), and kernel
\(K_2(A_2^{\rm ker}y_k/(y_dy_v^2y_1y_2))\),
\(A_2^{\rm ker}=Z^{H_++h_2-\theta-2v-2N_c}\). Multiplying this profile
by \(Z^{-6\eta}\) leaves the exact outside factor
\(Z^{\lambda_c+18\eta+\tau}\). Fresh tests are
\(\omega_a(Z^{\delta_N}y)y^{\pm i t_a}\); they and the density are
independent of the current \(k_{\rm new},\allowbreak{}f_{\rm new},\allowbreak{}\gamma\). The
signed outer weight retains both common extracted squared magnitudes,
every priority mask, one \(\mu(d_2)\), one \(\mu(v')\), and the second
frequency mask. Expanding the separated columns and inverting (B3.3)
restores the full expression term by term. Only then apply weighted
Cauchy over all outer tuples and bound their common zeros by one.

Here is an explicit sufficient reconstruction bound, which prevents a
hidden label cost. Let \(K_0\) bound the original slot count. At fixed
\(\gamma,\allowbreak{}f_{\rm new},\allowbreak{}k_{\rm new}\), splitting \(f_{\rm new}\) into
\(J,\allowbreak{}C,\allowbreak{}d_2,\allowbreak{}v'\), splitting \(J=sJ_2\), reconstructing
\(b_1b_2=q_0^2J_2^2s\), and choosing \(A_i,\allowbreak{}d_k\) gives at most \[
D(f_{\rm new})C(\gamma),\qquad
D(f)=d(f)^{9+4K_0},\quad
C(\gamma)=d(q_0)^{5+2K_0}d(t')^{2K_0}d(r_g)^{2K_0}.
\tag{B3.27}
\] Indeed \(d_k\mid f_{\rm new}\operatorname{rad}q_0\); all extracted
slot primes divide these fixed reconstructed factors and contribute only
the indicated additional divisor powers. Once these factors are chosen,
the row frequency is determined. No old \(f,\allowbreak{}h\) remains. By (B3.19), the
distinct triples have weighted count \[
\ll Z^{C_c+11\eta/2+o(1)},\qquad
C_c=\ell+R/2-j+t+g-\theta.
\tag{B3.28}
\] Extract this factor once and normalize the triple measure. The
unnormalized squared child norm equals \(Z^{F'}\mathcal E(N',\allowbreak{}V_c,\allowbreak{}M')\).
If the child estimate is \(\mathcal E\ll Z^{F'+\epsilon_{\rm child}}\),
its contribution is \(Z^{2F'+\epsilon_{\rm child}}\). The exact balance
\[
\varkappa_i+\lambda_c+C_c+2F_c=F
\tag{B3.29}
\] then gives \(Z^{F+40\eta+\tau+\pi+\epsilon_{\rm child}}\): the two
profile allowances and triple count contribute
\(9\eta/2+18\eta+11\eta/2=28\eta\), and \(2\delta_N\le12\eta\). Here
\(\pi\) is the aggregate freely chosen divisor, ray-expansion, dyadic
and coefficient loss. Both Cauchy inequalities take geometric means, so
neither this exponent nor the triple mass is doubled.

\textbf{Closure.} Set \(D=\lceil(M_{\max}+2)/d\rceil+1\). At depth
\(h\), use margins \(c_h=c_*-7h\eta\) and row cap \(M\le M_{\max}-hd\).
Choose \[
\eta\le\min\{d/16,c_*/(14D),\epsilon/(160D),c_*/1000\},
\qquad\tau,\pi\le\epsilon/(4D),\qquad
\tau\le c_*/1000,
\tag{B3.30}
\] and the reflected losses as above. Then (B3.26) gives a decrease of
at least \(d\), both margins remain at least \(c_*/2\), and \(F\) stays
in a bounded range through depth \(D\). Nonnegative rows preclude a
strict call at depth \(D\). Choose all local small-power shares within
\(\pi\), then the raw tail orders by (B3.4), then the required finite
smooth and height orders backwards by (B3.3), and finally \(Z\) large
enough for all fixed-window errors. The accumulated path loss is less
than \(\epsilon\). The terminal completion, principal counts and this
strict descent prove (B3.7). ∎

\hypertarget{b3.3.-the-marked-inverse-estimate-and-inverse-amplification}{%
\subsection{B3.3. The marked inverse estimate and inverse
amplification}\label{b3.3.-the-marked-inverse-estimate-and-inverse-amplification}}

Let \(\varphi_u(n)=\nu(n)\chi_n(u)^{\varepsilon_\chi}\), with one common
sign \(\varepsilon_\chi\in\{\pm1\}\), and define the centrally
normalized inverse and prime polynomials as in the main text. We prove
the full-width version of MI: \[
r+2z\le m-c_1,\qquad2r+8z\le3m-c_2
\quad\Longrightarrow\quad
\sum_{0<q_u\ll Z^m}|M_u(Z^r)Q_u|^2\ll Z^{m+\epsilon}.
\tag{B3.31}
\] The prime coefficients are arbitrary bounded row-independent
coefficients, and no mesh is required.

Let \(P=\prod_i p_i\), \(j=(n,\allowbreak{}P)\), and fix the subset of slot primes
dividing \(j\). Write \(n=jn_0\), \(P=jP_0\); then \((n_0,\allowbreak{}P_0)=1\) and
\(c=n_0P_0\) is squarefree and prime to \(j\). The exact identities
\(\mu(n_0)=\mu(c)\mu(P_0)\) and
\(\varphi_u(n)\varphi_u(P)=\varphi_u(j)^2\varphi_u(c)\) retain all
zeros. If \(G\) is the sum of assigned slot lengths, put \[
z_0=z-G,\qquad D_0=r+z-2G,
\qquad Z^{-(r+z)/2}=Z^{-G}Z^{-D_0/2}.
\tag{B3.32}
\] There are \(O(Z^{G+o(1)})\) assigned tuples. Their Hilbert-space
triangle mass is canceled by \(Z^{-G}\), while \(\varphi_u(j)^2\) is a
contraction. For a fixed \(j\), the old inverse profile has argument
\(q_n/Z^r=(q_j/Z^G)(q_c/Z^{D_0})/\prod_{i\ {\rm surviving}}(q_{p_i}/Z^{z_i})\).
Separate this one common profile by (B3.3), before choosing surviving
primes. Absorb \(\mu(p_i)\) in their coefficients. Conjugate the whole
residual polynomial when necessary; it remains to bound \[
R_0(u)=Z^{-D_0/2}\sum_{c\ {\rm sf}}\mu(c)\nu_0(c)
 \mathbf1_{(c,j)=1}\overline{\chi_c(u)}d_0(c)W_0(q_c/Z^{D_0}).
\tag{B3.33}
\] The prior Fourier measure is uniform over the fixed \(q_j/Z^G\)
interval.

Expand its square in a smooth row ball, write \(c_i=Cz_i\),
\((z_1,\allowbreak{}z_2)=1\), and apply (B3.5) to
\(\overline{\chi_{z_1}}\chi_{z_2}\), with divisor \(d'\mid C\) and
frequency \(h'\). Let \(B,\allowbreak{}\theta\) center \(C,\allowbreak{}d'\), put
\(P_1=B-\theta\ge-2\eta\), and retain \[
M_c=2D_0-m-2P_1,\qquad M_{\rm init}=M_c+6\eta+\tau_{\rm init},
\qquad0<q_{d'h'}\le Z^{M_{\rm init}}.
\tag{B3.34}
\] The actual conductor differs from \(2(D_0-B)\) by at most \(4\eta\),
and the two appearances of \(d'\) add \(2\eta\). Thus the omitted
complement is an actual-ratio tail. The principal diagonal is
\(O(Z^{m+o(1)})\), also for a bounded nonempty negative \(D_0\)-window.

The same Gauss calculation as the second canonical transform produces
\(a_0(z)\nu_*(z)\overline{\chi_z(d')}\chi_z(h')\), with the same
\(\nu_*\) on both sides after expanding \(G([z_2z_1^{-1}])\). Insert the
full mutual-gcd Möbius expansion \(z_a=sn_a\), and retain the common
zeros \((C,\allowbreak{}j)=1\), \((s,\allowbreak{}Cj)=1\), \((s,\allowbreak{}h')=1\). Set \[
t=C/d',\qquad k=d'h',\qquad f=d's,\qquad
\rho_{t,j}(n)=\mathbf1_{(n,\operatorname{rad}(tj))=1}.
\tag{B3.35}
\] The zero-preserving identity \[
a_0(sn)\nu_*(sn)\overline{\chi_{sn}(d')}\chi_{sn}(h')
=a_0(s)\nu_*(s)\overline{\chi_s(d')}\chi_s(h')
 a_0(n)\nu_*(n)\chi_n(k)\chi_n(f)^4
\tag{B3.36}
\] holds on \((s,\allowbreak{}n)=1\), using
\(\overline{\chi_n(d')}=\chi_n(d')\chi_n(d')^4\). The new label \(f\) is
squarefree. Slot priority allocation to \((t,\allowbreak{}d',\allowbreak{}s;n)\) leaves an
original subcollection; its survival zeros are exactly supplied by
\(\rho_{t,\allowbreak{}j}\) and \(\chi_n(f)^4\).

If \(v\) centers \(s\), the formal child parameters are \[
\begin{gathered}
N_c=D_0-B-v,\\ V_c=\theta+v,\\ F_c=D_0-P_1,
\\ N'=\max(0,N_c),\\0\le\delta_N=N'-N_c\le3\eta,
\\ F'=F_c+\delta_N.
\end{gathered}
\tag{B3.37}
\] For the six normalized norms \((C,\allowbreak{}d',\allowbreak{}s,\allowbreak{}h',\allowbreak{}n_1,\allowbreak{}n_2)\), the exact root
is \(y_d^{-1}y_s^{-1}(y_1y_2)^{-1/2}\), the old windows are
\(W_0(y_Cy_sy_1)\overline{W_0(y_Cy_sy_2)}\), and the kernel is
\(K_{\rm init}(A^{\rm ker}y_h/(y_dy_s^2y_1y_2))\), with
\(A^{\rm ker}=Z^{m+H-\theta-2(D_0-B)}\) and \(H\) the bare-frequency
center. Multiply the profile by \(Z^{-3\eta}\); the exact outside factor
is \(Z^{\varkappa_c+3\eta}\), where \(\varkappa_c=m-2D_0+P_1\). Fresh
tests are bounded clipped cutoffs times their Fourier powers. Inversion
restores the full kernel before Cauchy, with the mask (B3.34) and all
three common zeros outside it.

At fixed \(t,\allowbreak{}f,\allowbreak{}k\), a reconstruction is determined by \(d'\mid f\),
since \(s=f/d'\), \(C=td'\), \(h'=k/d'\). With at most \(K_0\) original
slots, its multiplicity is bounded by \(d(f)^{1+2K_0}d(t)^{2K_0}\). The
distinct \(t\)'s are projected from the genuine outer set once; each has
\(\log_Zq_t\le P_1+2\eta\), so their weighted count is
\(Z^{P_1+2\eta+o(1)}\). Normalize this measure after extracting that
factor. The new actual puncture obeys \(Q'\le P_1+G+3\eta\). Direct
substitution gives \[
\begin{aligned}
F'-M_{\rm init}-Q'-z_0&\ge c_1-9\eta-\tau_{\rm init},\\
4F'-3M_{\rm init}-6z_0&\ge c_2-22\eta-3\tau_{\rm init},\\
\varkappa_c+P_1+2F_c&=m.
\end{aligned}\tag{B3.38}
\] The first two follow respectively from the exact formal differences
\(m-r-2z+2G\) and \(3m-2r-8z+10G+2P_1\). Apply the canonical estimate
with margin \(\min(c_1,\allowbreak{}c_2)/2\), after choosing \(\eta,\allowbreak{}\tau_{\rm init}\)
sufficiently small. The last identity and \(2\delta_N\le6\eta\) give
total exponent at most
\(m+11\eta+\pi_{\rm init}+\epsilon_{\rm canonical}\). Choose these
shares below the requested \(\epsilon\), then fix raw tail and backwards
seminorm orders as in the common-profile lemma. This proves (B3.31),
hence MI at \(m=1\). ∎

\textbf{Unmarked consequence (UM).} Take no slots in (B3.31), \(Z=H\),
\(m=1\), and \(r=\log_H D\). If \(H\ge D^{1+\vartheta}\), then
\(1-r\ge\vartheta/(1+\vartheta)\) and \(3-2r\ge1\). Bounded scales are
direct; otherwise the lengths lie in a fixed bounded range. Thus for
every fixed \(\vartheta,\allowbreak{}\epsilon>0\), \[
\sum_{0<q_u\le H}\left|D^{-1/2}\sum_n\mu(n)\nu(n)
 \chi_n(u)^{\varepsilon_\chi}W(q_n/D)\right|^2
\ll H(HD)^\epsilon,\qquad H\ge D^{1+\vartheta}.
\tag{B3.39}
\] All element rows and all original zero extensions have survived. The
same finite-seminorm and polynomial-height bounds follow from the proof
above.

\hypertarget{b3.4.-the-enlarged-plain-moment-family}{%
\subsection{B3.4. The enlarged plain-moment
family}\label{b3.4.-the-enlarged-plain-moment-family}}

The induction requires a larger family than the physical application.
Let \(\psi_k(n)=\tau(n)\chi_n(k)\), where \(0<q_k\ll Z^m\). The
displayed twist \(\tau\) is a product of a fixed finite ray character
and moving residue-symbol factors. The union of the primes in those
moving factors, \textbf{before} canceling phases or reducing exponents
modulo six, has norm at most \(Z^q\). Put \(M=m+q\). An additional mask
\(\mathbf1_{(n,\allowbreak{}R)=1}\), with \(R\) squarefree and \(q_R\le Z^B\), may be
applied to every plain and prime factor. It is fixed throughout the
current row sum, although it may depend on previously frozen labels. A
cancellation between \(\tau\) and \(\chi_n(k)\) retains its natural row
zero and is never recast as an independently chosen row mask.

For a positive-slot family fix a nonzero integral ray modulus
\(\mathfrak M\), a ray subgroup containing the global units, and one
common smooth profile \(W\) supported in a fixed compact annulus. Fix an
anchor character \(\eta_0\) and an ideal \(Q\le\mathfrak M\). Write
\(Q_0=\operatorname{internalQ}(Q,\allowbreak{}\eta_0)\), and require \(Q_0\ne0\),
\(Q_0\ne\mathcal O\), and \(Q_0\subseteq(72)\). Each slot coefficient is
the individual relative ray character
\(\nu_i=\operatorname{relativeCharacter}(\mathfrak M,\allowbreak{}H,\allowbreak{}\eta_0,\allowbreak{}\theta_i)\),
where \(\theta_i\) is a character of the fixed ray quotient. With
bounded real powers \(\sigma_i\) and height parameters \(v_i\), write \[
Q_{k,i}=Z^{-z_i/2}\sum_{p\in\mathcal P_i^{\rm can}}
 \psi_k(p)\nu_i(p)W(q_p/Z^{z_i})
 (q_p/Z^{z_i})^{-\sigma_i+iv_i}.
\] Here \(\mathcal P_i^{\rm can}\) is the genuine canonical good-prime
pool at the scale \(Z^{z_i}\). Every natural zero and the same common
puncture remain in the summand. Arbitrary prime subsets are allowed in
MI, but are not asserted for MP. The physical application has separated
annular scales. There is no assertion here for arbitrary slot-specific
\(W_i\) or for a finite-character-expansion extension.

Write \(z=\sum_i z_i\), \(A=n_1+n_2+z\), and let \(S_k(Z^{n_i})\) be the
centrally normalized plain sums with two independent smooth profiles.
Their norm-twist heights may be independent and have polynomial costs.

\textbf{Plain-moment theorem.} The zero-slot assertion is unconditional.
For positive slots assume \[
7/10\le\kappa\le3/4,\qquad \beta_*\ge51/100,\qquad
2\beta_*-1\le\kappa.
\] Then \[
\begin{array}{ll}
z=0:&\displaystyle
 \sum_{\psi_k^*\ne1}|S_k(Z^{n_1})S_k(Z^{n_2})|^2\ll Z^{M+\epsilon}
 \quad\text{for all bounded }n_1,n_2\ge0;\\[2mm]
z>0:&\displaystyle
 n_1+n_2+6\kappa z\le M
 \quad\Longrightarrow\quad
 \sum_{k\ {\rm eligible}}|S_k(Z^{n_1})S_k(Z^{n_2})\prod_iQ_{k,i}|^2
 \ll Z^{M+\epsilon}.
\end{array}\tag{B3.40}
\] The positive assertion is the original-family
\sourcecode{Cycle25.\allowbreak{}Energy.\allowbreak{}terminal\_\allowbreak{}positive\_\allowbreak{}at}, with \sourcecode{Positive\allowbreak{}At}
defined in \sourcecode{OAI.\allowbreak{}Number\allowbreak{}Theory.\allowbreak{}Dirichlet\allowbreak{}L.\allowbreak{}Energy.\allowbreak{}Bands}. Its state
is \sourcecode{Centered\allowbreak{}Moment\allowbreak{}Energy\allowbreak{}State.\allowbreak{}Natural\allowbreak{}State}: its width is the row
width plus the character width; the presented character modulus has norm
at most \(Z^q\); and the common nonzero puncture has norm at most
\(Z^B\). Its radial weight and every original coefficient zero are
retained. The kept rows exclude \textbf{all} \sourcecode{Fixed\allowbreak{}Inducing\allowbreak{}Row}
presentations supported by the fixed modulus and common mask. This is an
explicit eligibility condition, not an identification with
\(\psi_k^*\notin\Theta\) for an arbitrary finite group.

For each requested loss and bounded width, puncture and length ranges,
the common mesh is chosen before the arithmetic family, slot count,
profiles and target height, uniformly for \(\kappa\in[7/10,\allowbreak{}3/4]\). After
fixing a finite slot index type and the common \(W\), one finite
seminorm set and one height order precede the fixed character and fixed
ideal. Their constant and eventual threshold precede the live slot
subset, slot characters, lengths, real powers, heights, radial state and
the two plain profiles. For the positive theorem a common whole-product
norm power is allowed explicitly; independent twists of the two plain
profiles are controlled by their seminorms. The zero-slot theorem
separately retains independent plain heights, unrestricted bounded
lengths, and all nonprincipal inducing characters, including
fixed-conductor rows.

Fixed conductor constants are absorbed by an arbitrarily small positive
width padding before the eventual threshold, and that padding is charged
to the requested loss. Moving support remains part of the character
width; it is not absorbed into a profile seminorm.

The physical translation is the proved adapter
\sourcecode{OAI.\allowbreak{}Seven\allowbreak{}Eighths.\allowbreak{}Cycle25Weighted\allowbreak{}Numerator\allowbreak{}Physical.\allowbreak{}finite\_\allowbreak{}physical\_\allowbreak{}pair\_\allowbreak{}positive\allowbreak{}At}
in \sourcecode{Cycle25.\allowbreak{}Energy.\allowbreak{}Physical.\allowbreak{}Physical\allowbreak{}Moment\allowbreak{}Input}. It constructs
the original \sourcecode{initial\allowbreak{}State}, identifies the exact coefficients,
and uses \sourcecode{weighted\_\allowbreak{}source\_\allowbreak{}keep\_\allowbreak{}iff} and
\sourcecode{weighted\_\allowbreak{}source\_\allowbreak{}exceptional\_\allowbreak{}norm} from
\sourcecode{Cycle25.\allowbreak{}Energy.\allowbreak{}Physical.\allowbreak{}Source\allowbreak{}State}. The exceptional
sixth-power-free physical rows have a fixed norm bound and disappear
from the growing lower row range. In that adapter \(W=\overline{F.W}\),
each ray-quotient character \(\theta_i\) is one, so the relative slot
coefficient retains \(\eta^{-1}\), the real power is \(33/50\), and the
slot height is \(-\Im z_i\). Thus the physical application uses
precisely the proved common-profile endpoint.

The zero-slot endpoint is
\sourcecode{Zeta\allowbreak{}Zero\allowbreak{}Free.\allowbreak{}Analytic.\allowbreak{}Moments.\allowbreak{}Unconditional.\allowbreak{}all\_\allowbreak{}nonprincipal\_\allowbreak{}at}.
Fixed nonprincipal rows are restored by masked lattice cancellation. The
common moving-support norm is counted before cancellation, the puncture
is excluded from the width, and independent profile and height controls
remain explicit. For the child analysis, fix a finite conjugation-stable
group \(\Theta\) of characters supported on \(S\), containing the fixed
phases, their inverses and \(n\mapsto\chi_n(-1)\). This group is
bookkeeping for the exceptional child terms; it is not substituted for
the original positive-slot eligibility condition. Taking
\(m=1,\allowbreak{}q=0,\allowbreak{}n_1=n_2\) gives MP in that positive family and MP0 in the full
nonprincipal zero-slot family.

\textbf{Removing fixed extra masks.} Call a presentation natural if its
only zeros are those from its primitive conductor, redundant displayed
moving factors and row radical, and \(S\). Let \(\psi_k^0\) be obtained
by deleting the additional fixed \(R\)-mask. Multiplicativity gives the
exact expansions \[
\begin{gathered}
S_{k,R}(X)=\sum_{d\mid R}\frac{\mu(d)\psi_k^0(d)}{\sqrt{q_d}}
 S_k^0(X/q_d),\\
Q_{k,i,R}=Q_{k,i}^0-Z^{-z_i/2}\sum_{p\mid R}\psi_k^0(p)\nu_i(p)\\[-1pt] \qquad{}\times W(q_p/Z^{z_i})(q_p/Z^{z_i})^{-\sigma_i+iv_i}.
\end{gathered}
\tag{B3.41}
\] The total absolute coefficient mass is \(Z^{o(1)}\), because
\(\prod_{p\mid R}(1-q_p^{-1/2})^{-1}\ll_\epsilon Z^\epsilon\); on a
frozen prime's annulus \(Z^{-z_i/2}\ll q_p^{-1/2}\). Apply both
identities simultaneously and use Minkowski. The remaining affine length
\(A+(6\kappa-1)z\) only decreases. If all slots disappear, the zero-slot
rows include the original eligible rows. Nonempty subunit plain scales
are clipped to one by bounded dilation and still do not exceed their old
nonnegative lengths. Thus it suffices to prove the natural assertion
uniformly, keeping this deletion available at each child.

\textbf{Reflection of a whole plain factor.} Let \(\psi_k^*\) induce
\(\psi_k^0\), let \(R_{0,\allowbreak{}k}\) be its redundant natural radical, and let
\(C_k=3Q_k/(2\pi)^2\), where \(Q_k\) is its primitive conductor norm.
Each good prime occurs once in either that conductor or the redundant
radical, so \[
C_kq_{R_{0,k}}\ll Z^M.
\tag{B3.42}
\] Unit and \(S\)-supported row factors, with valuations reduced modulo
six, give only finitely many fixed ray sectors. For a nonprincipal
inducing character, the ordinary Hecke functional equation {[}GZ,
equation (1.1){]} and Mellin inversion give the centrally normalized
primitive reflection at scale \(C_k/X\), with root number of modulus one
and profile \[
\widehat{W^\sharp}(s)=\widehat W(1-s)\frac{\Gamma(s)}{\Gamma(1-s)}.
\tag{B3.43}
\] This identity follows by substituting the functional equation in the
Mellin integral and replacing \(s\) by \(1-s\). There is no \(L\)-pole.
Moving the inverse Mellin contour left shows that \(W^\sharp\) and its
Euler derivatives are bounded at zero; moving it right gives arbitrary
decay at infinity. Stirling's formula and
\(\widehat{Wy^{it}}(s)=\widehat W(s+it)\) give finite-seminorm and
polynomial-height bounds.

Delete the redundant Euler factors before reflection and restore them
geometrically afterward. This gives the exact natural-character identity
\[
S_k^0(X;W)=\varepsilon_k
 \sum_{d_0\mid R_{0,k}}\sum_{\operatorname{rad}(h_0)\mid R_{0,k}}
 \frac{\mu(d_0)\psi_k^*(d_0)\overline{\psi_k^*(h_0)}}{\sqrt{q_{d_0}q_{h_0}}}
 S_{\overline{\psi_k^0}}\!\left(\frac{C_kq_{d_0}}{Xq_{h_0}};W^\sharp\right).
\tag{B3.44}
\] Its coefficient mass is \(Z^{o(1)}\). Decompose \(W^\sharp\) into
annuli. Lower annuli carry summable central coefficients
\(O(2^{-j/2})\); upper annuli have arbitrary summable decay. Discard the
upper tail after relative enlargement \(Z^{\xi/2}\). A main scale
\(Y<Z^{-\xi}\) is also negligible: absolute counting and rapid decay
give \(O_B(Y^{B-1/2})\) for its normalized sum, for arbitrarily large
fixed \(B\). Row-dependent coefficients and retained scales are handled
first by their absolute mass and then by a rowwise scale supremum. The
latter is bounded by integrating the function and its logarithmic scale
derivative over \(O(\log Z)\) unit boxes; for two scales use the
two-variable version. The derivative replaces a profile by \(-W/2-yW'\),
so the profile lemma applies. Conjugating a whole reflected factor
inside a product's absolute value restores the original row character.
By (B3.42), every retained nonempty reflected scale has declared length
\[
n_{\rm ref}\le M-n+\xi.
\tag{B3.45}
\] All fixed endpoint and scale-box factors are included by increasing
the final \(Z\)-threshold. A negative right side means the retained
range is empty. Reflecting each plain longer than \(M/2\) once reduces
the zero-slot problem to the padded core \[
n_i\le M/2+\xi,\qquad A\le M+2\xi\le M+\delta.
\tag{B3.46}
\]

\textbf{Prime cost.} On the original eligible rows, each individual slot
character is nonprincipal. The hypothesis \(2\beta_*-1\le\kappa\)
permits moving the smooth von Mangoldt Mellin contour to \(s_\kappa+e\),
\(s_\kappa=(1+\kappa)/2\). Here
\(|L'/L(s_\kappa+e+it)|\ll_e\log(Q(3+|t|)^2)\): take the holomorphic
Euler branch of \(g=\log L\) on the disk centered at \(2+it\) with
radius \(2-s_\kappa-e/4\). It is zero free. Polynomial strip growth from
the ordinary functional equation, Stirling and the strip principle
bounds \(\Re g\) by \(O(\log(Q(3+|t|)^2))\), while \(g(2+it)=O(1)\) by
the Euler product. Borel--Carathéodory on a slightly smaller disk and
Cauchy's derivative estimate on a further disk containing
\(s_\kappa+e+it\) give the claim. All gaps are fixed multiples of \(e\).
Prime powers of degree at least two cost \(O(P^{1/2+\epsilon})\).
Dividing the annular weight by \(\log(Py)\) removes the von Mangoldt
factor. Redundant natural zeros delete only \(O(\log Z)\) primes. After
central normalization and multiplication over the fixed slot list, \[
\left|\prod_iQ_{k,i}\right|^2
 \ll Z^{\kappa z+\epsilon_1}\,p_J(W)^b\langle\mathbf t\rangle^h.
\tag{B3.47}
\] The contour displacement costs \(2ez\), depending on total length
rather than slot count; height twists merely translate rapidly
decreasing Mellin transforms. Thus the exponent is independent of the
required derivative and height orders.

\hypertarget{b3.5.-induction-order-and-the-centered-comparison}{%
\subsection{B3.5. Induction order and the centered
comparison}\label{b3.5.-induction-order-and-the-centered-comparison}}

Choose a width floor \(\rho>0\), step \(\sigma>0\), and padding
\(\delta\ll\min(\rho,\allowbreak{}\sigma)\). Divide the bounded \(M\)-range into
bands of length \(\sigma/4\). Prove all zero-slot bands in increasing
width, then all positive-slot bands. Within each band prove the
uncentered range \(A\le5M/6\) first, then the remaining allowed range
using a centered difference. Reflection then supplies unrestricted
zero-slot lengths. At \(M\le\rho\), absolute counting gives the padded
core with excess at most \(\rho+\delta\). The positive-slot case follows
from the completed zero-slot estimate and (B3.47), since \[
z\le M/(6\kappa),\qquad\kappa z\le M/6.
\tag{B3.48}
\] Its additional terminal excess is at most \(\rho/6\).

For \(A>5M/6\), set \(L=M/4\). If an original plain has length \(b<L\),
reflect the other. Otherwise compare with plain lengths \(L,\allowbreak{}A-z-L\) and
the same profiles, reflecting the longer comparison factor only when
estimating that comparison separately. In the positive-slot range, with
\(q_\kappa=6\kappa-1>2\), \(z<(M/6)/q_\kappa<M/12\); hence both
comparison lengths are at least \(L\). Whole-factor reflection gives \[
\begin{gathered}
A_{\rm comp}\le3M/2-A+2z+\xi,
\\
\min\{5M/6-A_{\rm comp},\,M-A_{\rm comp}-q_\kappa z\}
\ge \mathfrak m(7/10)M-\xi,
\\
\mathfrak m(\kappa)=\frac{2\kappa-1}{2(6\kappa-1)}.
\end{gathered}
\tag{B3.49}
\] For example, write \(t=M-A<M/6\) and use \(z\le t/q_\kappa\): the
first margin is at least \(M/3-t(1+2/q_\kappa)-\xi\), and the second at
least \(M/2-t(2+2/q_\kappa)-\xi\). Both have the displayed lower bound.
For zero slots each comparison has margin at least \(M/6-\xi\) to the
uncentered boundary. Choose \(\xi<\rho\mathfrak m(7/10)/2\). Every
comparison then calls the already proved uncentered stage in its own
band.

It remains to treat originals with both plain lengths at least \(L\).
Put \(X_i=Z^{n_i}\), \(Y_1=Z^L\), \(Y_2=X_1X_2/Y_1\). Subtract the
unreflected comparison, keeping the two equal-product rectangles
together. After plain powers \(b_i\) are extracted, its coefficient is
\[
D_b(l_1,l_2)=W_1(q_{b_1l_1}/X_1)W_2(q_{b_2l_2}/X_2)
 -W_1(q_{b_1l_1}/Y_1)W_2(q_{b_2l_2}/Y_2),
\quad X_1X_2=Y_1Y_2.
\tag{B3.50}
\] The same \(b_i\) is used in both terms even if an allocation makes
one term zero. The comparison has already been bounded on the
permissible original rows. Enlarge the nonnegative norm of the
\textbf{whole difference} to all rows of a smooth ball. In the
uncentered range enlarge the original squared norm instead. No Poisson
formula is applied to a character-eligibility indicator.

The invariant coefficient carried through the transformations is \[
\prod_{i\in I'}\nu_i(p_i)W(q_{p_i}/Z^{z_i})(q_{p_i}/Z^{z_i})^{-\sigma_i+iv_i}\,
D_b(l_1,l_2)\,\tau_1(u)q_u^{it}\,\mathbf1_{(u,R)=1}\mathbf1_{s\mid u},
\qquad u=\Bigl(\prod_{i\in I'}p_i\Bigr)l_1l_2,
\tag{B3.51}
\] where \(s\) is squarefree and fixed. For an uncentered input omit the
second rectangle. The character, mask and norm power act on the whole
product; this is valid even when its factors share primes, because all
zero powers retain their zeros.

Use one aggregate support tolerance. For the eventual fixed slot count
\(N\), choose \(H_N\) containing the sum of absolute logarithmic
endpoints of all its slot windows, the two plain windows and the
finitely many later support enlargements. Put \(\theta_N=H_N/\log Z\).
Every slot subset has total logarithmic ratio error at most
\(\theta_N\). Extracted plain powers always use exact norms. Thus one
stage's boundary errors are \(O(\theta_N)\), with a numerical
coefficient independent of \(N\); there is no tolerance charged once per
slot.

\hypertarget{b3.6.-first-poisson-from-plain-products-to-gauss-polynomials}{%
\subsection{B3.6. First Poisson: from plain products to Gauss
polynomials}\label{b3.6.-first-poisson-from-plain-products-to-gauss-polynomials}}

Set \(H=Z^m\), \(X=Z^A\). At frequency zero the mean of the row
character can be nonzero only when all net prime exponents are zero
modulo six. No prime then has total multiplicity one in the product of
the two full column products, so that product is powerful. A powerful
ideal is a square times a cube of a squarefree ideal; summing over the
latter counts \(O(Y^{1/2+\epsilon})\) powerful ideals up to norm \(Y\).
With \(Y\asymp X^2\), divisor-bounded allocations, central factor
\(X^{-1}\), and row mean at most \(H\), the zero term is
\(O(Z^{m+o(1)})\).

For nonzero frequencies write the two full products as \(Ca,\allowbreak{}Db\), where
\(C,\allowbreak{}D\) contain their \textbf{complete} common prime support. Thus
\((a,\allowbreak{}b)=(ab,\allowbreak{}CD)=1\). At every common prime fix its full valuation in
each plain variable and its occurrence in a slot. Extract all those
powers, freeze any slot supplying the prime, and puncture every
remaining factor there. The same allocation acts on both rectangles. A
scalar zero forced by an old character or mask remains zero before any
enlargement.

Let \(c,\allowbreak{}d,\allowbreak{}p\) be the logarithmic norms of
\(C,\allowbreak{}D,\allowbreak{}\operatorname{rad}(CD)\). Let \(r\) contain common primes of
nonzero net exponent modulo six, let \(R=\log_Zq_r\), and let \(\xi_r\)
be their primitive character. The remaining common row mask introduces a
divisor \(e\), \(E=\log_Zq_e\le p-R\). If \(A_C,\allowbreak{}A_D\) are the allocated
convolution coefficients, omitting the whole-product twist, put
\(\tau_C(a)=\tau(a)\chi_a(er)\overline{\xi_r(a)}\) and similarly for
\(\tau_D\). Residue-class Poisson and CRT give, up to a bounded frozen
scalar, \[
\frac{H}{Xq_e\sqrt{q_r}}\sum_{h\ne0}\mathcal G_{\xi_r}(r,h)
\sum_{(a,b)=1}\frac{A_C(a)\overline{A_D(b)}}{\sqrt{q_aq_b}}
 \tau_C(a)\overline{\tau_D(b)}R(a,b)
 \mathcal G(a,h)\overline{\mathcal G(b,-h)}
 K_1\!\left(\frac{Hq_h}{q_eq_rq_aq_b}\right).
\tag{B3.52}
\] Here \(\mathcal G\) is the normalized two-variable Gauss sum of
(B1.12); \(|\mathcal G_{\xi_r}|\le1\). To check (B3.52), Poisson modulo
\(rab\) after \(k=ek'\) contributes \(H/(q_eq_rq_aq_b)\); the three
unnormalized Gauss sums restore \(\sqrt{q_rq_aq_b}\). CRT contributes
\(\chi_a(br)\overline{\chi_b(ar)}\overline{\xi_r(ab)}\), and reciprocity
replaces the \(a,\allowbreak{}b\) cross ratio by \(R(a,\allowbreak{}b)\). These are exactly the
displayed factors. Every new moving prime occurs in the whole-product
character and punctures all residual factors, including when phases
cancel. Its full support length is bounded by \[
\widetilde q=q+R+E,\qquad
K_0=2A-c-d+R+E-m.
\tag{B3.53}
\] The first kernel's nominal frequency scale is \(Z^{K_0}\).

To localize both Poisson formulas used here, choose one nonnegative
dyadic partition of frequency norms with fixed relative supports. In a
sector of aggregate outer norms let \(T_{\rm sec}\) be the maximum
nominal frequency scale. Retain whole partition weights meeting
\(q_h\le T_{\rm sec}Z^{\xi/2}\). Their union has upper length
\(K\le K_0+\delta_1\), \(0\le\delta_1<\xi\), after \(Z\) exceeds a fixed
endpoint threshold. Each discarded weight has actual kernel argument at
least \(Z^{\xi/4}\), because the full column-product support costs at
most \(e^{2H_N}\). Formula (B3.4), \(|\mathcal G(a,\allowbreak{}h)|\le\sqrt{q_a}\),
and the finite bound \(|F(u,\allowbreak{}v;j)|\le q_uq_v\) justify arbitrary tail
savings in the two formulas respectively. The original columns have
bounded polynomial norms and divisor-bounded coefficients, so these raw
bounds cost a fixed power before the decay order is chosen. Retained
weights and the full kernel are common to both rectangles and
independent of live columns.

On each retained dyad define the displayed bridge off \((a,\allowbreak{}b)=1\) by its
separate Gauss sums, fixed bicharacter, masks and formal norms. It is an
artificial coefficient there, not a formula for the original Gauss
correlation. Insert the \textbf{complete} projector
\(\sum_{s\mid(a,\allowbreak{}b)}\mu(s)\) before taking absolute values. Let
\(s_0=\log_Zq_s\). Separate the kernel and inverse roots in the
normalized row and two whole-product norms by (B3.3), adding one
aggregate outer-ratio coordinate if necessary. This introduces one norm
power on each entire column, never different powers on its two plains or
rectangles. Only now use Cauchy--Schwarz and
\(|\mathcal G_{\xi_r}|\le1\). The positive C-side norm is \[
\mathcal N_C=\sum_{h\ne0}\omega_K(q_h/Z^K)
 \left|Z^{-(A-c)/2}\sum_{s\mid a}A_C(a)\tau'_C(a)q_a^{it}\mathcal G(a,h)\right|^2.
\tag{B3.54}
\] It has coefficient (B3.51), and similarly for D. The outside exponent
is \(m-A-R/2-E\). The common-support and Möbius labels cost
\(Z^{p+s_0+o(1)}\): count the radical first, then its bounded-range
powers by
\(\#\{v:\operatorname{rad}v\mid c,\allowbreak{}\ q_v\le Z^T\}\le Z^{aT}\prod_{p\mid c}(1-q_p^{-a})^{-1}\),
choosing \(a>0\) small; all allocations are divisor bounded.

It suffices to prove \[
\mathcal N_C\ll Z^{A-c+\widetilde q+B_c-s_0+\epsilon_G},
\qquad B_c=(3c-5d-R)_+/6,
\tag{B3.55}
\] and its exchanged counterpart. Indeed their geometric mean, with the
outside factors, has exponent \(M+\{B_c+B_d-(c+d-2p-R)\}/2\). Primewise,
\[
B_c+B_d\le c+d-2p-R.
\tag{B3.56}
\] For multiplicities \(i\ge j\ge1\) and \(r_p=\mathbf1_{6\nmid i-j}\),
only the larger side can have a positive local numerator. Six times the
difference between the right side and \((3i-5j-r_p)_+/6\), when the
latter is positive, is \(3i+11j-12-5r_p\ge0\); otherwise the right side
is already nonnegative. Positive parts of sums are bounded by sums of
positive parts. This proves (B3.56).

\hypertarget{b3.7.-enlarging-a-gauss-row-norm}{%
\subsection{B3.7. Enlarging a Gauss-row
norm}\label{b3.7.-enlarging-a-gauss-row-norm}}

Let \(w\) be a column length extracted from (B3.54), and \(w_o\le w\)
the additional moving-radical length, whose squared scalar coefficient
is \(Z^{-w_o}\). Put \(a_0=A-c-w\). The allowance (B3.55) for the
unweighted remaining norm is \[
\Lambda_c=a_0+\widetilde q+w_o+w+B_c-s_0+\epsilon_G.
\tag{B3.57}
\] Define \(J=d-c+(K_0-K)-2w+w_o\), \(J_+=\max(J,\allowbreak{}0)\). With zero slots
take \(w=w_o=0\), \(\ell=0\), and enlarge the row norm to a smooth ball
of length \(K+g\), \(g=J_++\sigma\).

With positive slots, put \(\ell_* =\sigma/3\), choose the mesh below
\(\sigma/6\), and use the fixed pool
\(\mathcal P=\{p:Z^{\ell_*}/2<q_p\le Z^{\ell_*},\allowbreak{}\ p\notin S\}\). The
fixed-field prime ideal theorem {[}TZ, Theorem 1.1, specialized to the
identity extension of \(\mathbb K\){]} gives
\(|\mathcal P|=Z^{\ell_*+o(1)}\). The pool is disjoint from every live
slot window once \(Z\) is large. For each \(h\ne0\), omit primes
dividing \(h,\allowbreak{}s\) or frozen supports; only \(O(\log Z)\) are omitted.

For \(p\nmid uh\), the prime-power Gauss identities (B1.13) and CRT give
\[
\mathcal G(p^iu,h)=R(p,u)^i\chi_p(u)^{2i}\mathcal G(p^i,h)\mathcal G(u,h),
\qquad\mathcal G(u,hp^6)=\mathcal G(u,h).
\tag{B3.58}
\] Replacing \(h\) by \(hp^6\) changes only column valuations
\(i=1,\allowbreak{}6,\allowbreak{}7\). Allocate their powers to the two plain variables, since the
pool is disjoint from the slots, and puncture the residual product at
\(p\). The same allocation updates both rectangles. Jensen in the
eligible primes and this finite allocation list bounds the original
squared value by the average of the amplified squared value and these
extracted squared errors. A fixed output \(hp^6\) has boundedly many
representations, since all pool primes have norm \(\asymp Z^{\ell_*}\)
and the output norm has bounded logarithmic length. Thus the amplified
norm gets one density factor \(Z^{-\ell_*+o(1)}\). For it take
\(\ell=\ell_*\), \(g=J_++2\sigma\), which contains the new rows because
\(6\ell_*=2\sigma\).

The extracted alternatives are exactly:

\begin{longtable}[]{@{}
  >{\raggedleft\arraybackslash}p{(\columnwidth - 6\tabcolsep) * \real{0.2500}}
  >{\raggedleft\arraybackslash}p{(\columnwidth - 6\tabcolsep) * \real{0.2500}}
  >{\raggedleft\arraybackslash}p{(\columnwidth - 6\tabcolsep) * \real{0.2500}}
  >{\raggedleft\arraybackslash}p{(\columnwidth - 6\tabcolsep) * \real{0.2500}}@{}}
\toprule\noalign{}
\begin{minipage}[b]{\linewidth}\raggedleft
Column valuation \(i\)
\end{minipage} & \begin{minipage}[b]{\linewidth}\raggedleft
Squared extraction coefficient
\end{minipage} & \begin{minipage}[b]{\linewidth}\raggedleft
\(w\)
\end{minipage} & \begin{minipage}[b]{\linewidth}\raggedleft
\(w_o\)
\end{minipage} \\
\midrule\noalign{}
\endhead
\bottomrule\noalign{}
\endlastfoot
1 & \(q_p^{-1}\) & \(\ell_p\) & \(\ell_p\) \\
6 & \((1-q_p^{-1})^2\) & \(6\ell_p\) & 0 \\
7 & \(q_p^{-1}\) & \(7\ell_p\) & \(\ell_p\) \\
\end{longtable}

Here \(\ell_p=\log_Zq_p\in[\sigma/6,\allowbreak{}\sigma/3]\). At valuation six the
exact zero power \(\chi_p(u)^{12}\) is its already common puncture;
retain the fixed phase \(R(p,\allowbreak{}u)^6\). At valuations one and seven the
local character remains displayed and its prime is counted in \(w_o\).
Freeze \(p\) for each error, set \(\ell=0,\allowbreak{}g=J_++\sigma\), and apply the
second transform directly, with no further amplification. The normalized
average \(|\mathcal P|^{-1}\sum_p\) costs no second prime count. Always
\[
0\le w_o\le w\le7\sigma/3.
\tag{B3.59}
\] The coefficient remains (B3.51) times \(\mathcal G(u,\allowbreak{}h)\), including
its whole difference, common puncture and \(s\)-divisibility.

Add row zero only when passing to the final smooth ball. Formula (B1.13)
gives \(\mathcal G(u,\allowbreak{}0)=0\) unless \(u\) is a sixth power and
\(|\mathcal G(u,\allowbreak{}0)|\le\sqrt{q_u}\). If \(Y=e^{H_N}Z^{a_0}\) is a
nonempty column upper bound and \(s\mid u\), then \(q_s^6\le Y\).
Counting \(O(Y^{1/6})\) sixth powers bounds the normalized squared zero
row by \(Z^{-a_0}Y^{4/3+o(1)}\le Z^{a_0-s_0+2\theta_N+o(1)}\), within
(B3.57). Zero is never amplified.

\hypertarget{b3.8.-second-transform-and-the-actual-children}{%
\subsection{B3.8. Second transform and the actual
children}\label{b3.8.-second-transform-and-the-actual-children}}

Restore live slot sums before squaring. If \(B(u)\) is the complete
convolution coefficient, with \(s\mid u\), (B1.14) gives \[
\begin{gathered}
Z^{-\ell}\sum_h\Phi_2(q_h/Z^{K+g})
 \left|Z^{-a_0/2}\sum_u B(u)\tau_1(u)q_u^{it}\mathcal G(u,h)\right|^2\\
=Z^{K+g-\ell-a_0}\sum_{u,v}
 \frac{B(u)\overline{B(v)}\tau_1(u)\overline{\tau_1(v)}q_u^{it}q_v^{-it}}
      {\sqrt{q_uq_v}}\\
\qquad{}\times\sum_jF(u,v;j)K_2\!\left(\frac{Z^{K+g}q_j}{q_uq_v}\right).
\end{gathered}
\tag{B3.60}
\] At \(j=0\), (B1.16) leaves \(u=v\), \(F(u,\allowbreak{}u;0)=\phi(u)\le q_u\).
There are \(O(Z^{a_0-s_0+\theta_N})\) supported multiples of \(s\), so
its exponent is \(K+g-\ell-s_0+\theta_N+o(1)\). Its nominal excess above
(B3.57) is \[
(A-M)-d-(K_0-K)-w_o-B_c+g-\ell
\le A-M+5\sigma/3+\delta_1.
\tag{B3.61}
\] For the inequality use \(J_+\): when \(J\ge0\), the part
\(-d-(K_0-K)-w_o+J_+\) equals \(-c-2w\); otherwise it is at most
\(\delta_1\). The largest remaining increment is
\(2\sigma-\ell_*=5\sigma/3\). This is a terminal loss, with an
additional \(\delta\) in the padded zero-slot core.

Localize the genuine nonzero \(j\)-sum at its nominal scale
\(Z^{2a_0-K-g}\) by the whole-dyad rule of §B3.6, before extracting
common support. Now write \(u=D_2a,\allowbreak{}v=E_2b\), extracting their complete
common support. Put \[
\begin{gathered}
c_2=\log_Zq_{D_2},\\ d_2=\log_Zq_{E_2},\\ b_2=(c_2+d_2)/2,
\\ p_2=\log_Zq_{\operatorname{rad}(D_2E_2)},\\ G_c=(D_2,E_2),\\ g_2=\log_Zq_{G_c}.
\end{gathered}
\] Use the exact complete-support factorization (B1.19) \[
F(D_2a,E_2b;j)=F(D_2,E_2;j)R(a,E_2)R(b,D_2)R(a,b)
 \chi_a(j)\overline{\chi_b(-j)}
\tag{B3.62}
\] only on \((a,\allowbreak{}b)=(ab,\allowbreak{}D_2E_2)=1\). It permits arbitrary powers and all
slot/plain coincidences. All extracted primes puncture all remaining
factors, and any slot supplying them is frozen. In particular a plain
\(p^5\) and slot \(p\) on one side leave the zero-extended factor
\(\chi_p(j)^6\), not the constant one.

The local correlation table (B1.17) implies \(G_c\mid j\). At a common
prime of equal multiplicity \(i\not\equiv0\pmod6\), split according as
\(j/G_c\) is a unit or nonunit. Let \(t_2\) be the radical length of the
unit primes and \(V\) that of the nonunit primes; let \(V_{\rm id}\) be
the latter radical. Write \(j=G_cV_{\rm id}h'\). The partitioned scalar
\(\mathbf1_{\rm part}F(D_2,\allowbreak{}E_2;G_cV_{\rm id}h')/Z^{g_2-t_2}\), extended
by zero off its partition, has modulus at most one. The necessary local
bounds are

\begin{longtable}[]{@{}lr@{}}
\toprule\noalign{}
Common multiplicities and divided frequency & Bound \\
\midrule\noalign{}
\endhead
\bottomrule\noalign{}
\endlastfoot
\(i=i,\allowbreak{}\ 6\nmid i\), unit & \(P^{i-1}\) \\
\(i=i,\allowbreak{}\ 6\nmid i\), nonunit & \(P^i\) \\
\(i=i,\allowbreak{}\ 6\mid i\), either & \(P^i\) \\
\(i>j_0,\allowbreak{}\ 6\mid j_0\), unit & \(P^{j_0}\) \\
\end{longtable}

Every other unequal case is zero. At the fixed factor
\(v=G_cV_{\rm id}\), write each valuation modulo six. A nonzero residue
retains its moving character; a zero residue is exactly the already
present common puncture. Active fixed primes are contained in those
counted by \(t_2+V\). This operation factors only \(\chi_n(v)\); it
keeps the natural varying-row zeros of \(\chi_n(h')\). Hence the nominal
child widths are \[
m'=2a_0-K-g-g_2-V,\quad q'=\widetilde q+w_o+t_2+V,
\qquad M'=m'+q'=M+J-g-g_2+t_2\le M-\sigma.
\tag{B3.63}
\] The actual retained row weight, pulled back under
\(j=G_cV_{\rm id}h'\), lies in \(q_{h'}\le Z^{m'+\delta_2}\),
\(0\le\delta_2<\xi\). A nonempty range has \(m'+\delta_2\ge0\); declare
\(m'_{\rm act}=m'+\delta_2\), \(M'_{\rm act}=M'+\delta_2\). Extend the
pulled-back weight and partitioned scalar by zero on the rest of this
ball. They are still part of the signed expression.

The condition \(s\mid u,\allowbreak{}v\) forces \(s\mid\operatorname{rad}(D_2E_2)\).
Thus common-support labels cost \(Z^{p_2-s_0+o(1)}\), not \(Z^{p_2}\).
The second-transform exponent ledger is

\begin{longtable}[]{@{}lr@{}}
\toprule\noalign{}
Factor & Exponent \\
\midrule\noalign{}
\endhead
\bottomrule\noalign{}
\endlastfoot
Original squared central normalization & \(-a_0\) \\
Poisson factor and inverse roots & \(K+g-a_0\) \\
Amplifier density & \(-\ell\) \\
Common-support count with fixed \(s\) & \(p_2-s_0\) \\
Common correlation & \(g_2-t_2\) \\
Conversion to normalized residual products & \(a_0-b_2\) \\
\end{longtable}

Subtracting this sum from (B3.57) leaves the inner allowance \[
M'+\Delta,\qquad \Delta=b_2-p_2+w+B_c+\ell\ge w+\ell.
\tag{B3.64}
\] The last inequality uses \(b_2\ge p_2\).

Before separating child factors, define the right side of (B3.62)
artificially on all individually allowed residual pairs, keeping the
original common-support masks, formal full-product norms, inverse roots,
kernel and pulled-back row weight. Insert the full residual coprimality
projector \(\sum_{t\mid(a,\allowbreak{}b)}\mu(t)\). Only its complete sum recovers
the genuine correlation; the extra common primes in an individual
\(t\)-term are not added to \(D_2,\allowbreak{}E_2\) or its frequency restrictions.
Separate the full kernel and roots in the row and two whole-product
norms as before. Then allocate each prime of \(t\) with \[
\mathbf1_{p\mid\prod_i n_i}
=\sum_{\varnothing\ne J}(-1)^{|J|+1}\prod_{i\in J}\mathbf1_{p\mid n_i},
\tag{B3.65}
\] where the factors are the two plains and live slots. A selected plain
is \(n_i=pl_i\) with no restriction on \(l_i\); divide both rectangle
scales for it by \(q_p\). A selected slot is frozen. Selecting two
different primes in one prime slot is zero. Quotients may retain \(p\),
and no new one-variable puncture is inserted.

This proves the needed coefficient invariance: each separated child side
has one whole-product character, common mask and norm power in both
plains and both rectangles; surviving slots retain their original
\(\nu_i\), without Gauss factors. The two sides may have different
whole-product norm powers. Their fixed-inducing predicates agree under
the paired reflection, with the same fixed modulus and retained natural
masks; this is \sourcecode{Family.\allowbreak{}paired\_\allowbreak{}inducing\_\allowbreak{}iff} in
\sourcecode{OAI.\allowbreak{}Number\allowbreak{}Theory.\allowbreak{}Dirichlet\allowbreak{}L.\allowbreak{}Moments.\allowbreak{}Second\allowbreak{}Exceptional\allowbreak{}Family},
under its fixed-modulus condition \(Q\le(72)\). At each fixed-modulus
stage enlarge \(\Theta\) to contain every primitive character whose
conductor divides that modulus and the finitely many fixed ray twists
used here. This is still a finite group. On an exceptional row with
primitive inducing character \(\vartheta\), fixing the live slots leaves
precisely \[
\vartheta(l_1)\vartheta(l_2)\mathbf1_{(l_1l_2,R_*)=1}
 q_{l_1l_2}^{it}D_b(l_1,l_2)
\tag{B3.66}
\] for one squarefree polynomial-size \(R_*\), common to the two
variables and rectangles. Every extracted row scalar, including zero
scalars, is retained until the following positive or absolute estimate.

\textbf{Exact extraction normalization.} Put \(\alpha_1=a_0-c_2\),
\(\alpha_2=a_0-d_2\). For side \(j\), let \(F_j\) be the entrance slots
frozen by the genuine extractions, and \(I_j\) the live slots before
(B3.65). If \(T_j\) is the common product of its two pre-allocation
plain scales, then \(T_j\prod_{i\in I_j}Z^{z_i}=Z^{\alpha_j+e_j}\),
where \(e_j=\sum_{i\in F_j}\log_Z(q_{p_i}/Z^{z_i})\). For the final
\(t\)-allocation, let \(d_{j,\allowbreak{}i}\) be the product selected in plain
\(i\), and \(J_j\subset I_j\) the newly frozen slots. Put \[
a_{j,i}=\log_Zq_{d_{j,i}},\quad
r_j=a_{j,1}+a_{j,2}+\sum_{i\in J_j}z_i,\quad
\omega_j=\sum_{i\in J_j}\log_Z(q_{p_i}/Z^{z_i}).
\] Use lower-endpoint dyads \(Z^{t_-}\le q_t<2Z^{t_-}\). Every selected
product is divisible by \(t\), so \[
r_j+\omega_j\ge t_-,\qquad |\omega_j|\le\theta_N,
\qquad |e_j+\omega_j|\le\theta_N.
\tag{B3.67}
\] The last bound uses disjoint subsets of the current entrance slots.
Ancestor slots are not counted again. The raw child normalizer is
\(Z^{-(\alpha_j-r_j)/2}\), and its exact extraction scalar relative to
\(Z^{-\alpha_j/2}\) is \(Z^{-r_j/2}\). The formal post product of scales
is \(Z^{\alpha_j+e_j-r_j}\). These signed raw rectangles need not have
nonnegative individual lengths.

\textbf{Exact child eligibility.} Split each actual child by the literal
\sourcecode{Fixed\allowbreak{}Inducing\allowbreak{}Row} predicate, retaining its fixed modulus and
source masks. For the retained family,
\sourcecode{Family.\allowbreak{}nonexceptional\_\allowbreak{}iff} in
\sourcecode{OAI.\allowbreak{}Number\allowbreak{}Theory.\allowbreak{}Dirichlet\allowbreak{}L.\allowbreak{}Moments.\allowbreak{}Second\allowbreak{}Exceptional\allowbreak{}Family}
identifies the nonexceptional part with \[
z\ne0\quad\text{and}\quad
\neg\operatorname{FixedInducingRow}(\tau_\chi,Q,\mathrm{fixedBadMask},1,z).
\] The actual radial norm is the canonical radial norm by
\sourcecode{source\allowbreak{}Radial\_\allowbreak{}eq\_\allowbreak{}canonical} in
\sourcecode{OAI.\allowbreak{}Number\allowbreak{}Theory.\allowbreak{}Dirichlet\allowbreak{}L.\allowbreak{}Moments.\allowbreak{}Second\allowbreak{}Nonexceptional\allowbreak{}Chosen\allowbreak{}Block}.
In \sourcecode{OAI.\allowbreak{}Number\allowbreak{}Theory.\allowbreak{}Dirichlet\allowbreak{}L.\allowbreak{}Energy.\allowbreak{}Child\allowbreak{}State},
\sourcecode{canonical\_\allowbreak{}mask\_\allowbreak{}nonexceptional} restores the common nonzero
puncture without changing the exclusion, and \sourcecode{canonical\allowbreak{}State}
supplies the required natural state. The smaller-width moment estimates
therefore apply to precisely these rows:
\sourcecode{actual\_\allowbreak{}canonical\_\allowbreak{}child\_\allowbreak{}bands} in
\sourcecode{OAI.\allowbreak{}Number\allowbreak{}Theory.\allowbreak{}Dirichlet\allowbreak{}L.\allowbreak{}Energy.\allowbreak{}Canonical\allowbreak{}Child\allowbreak{}Bound} is the
proved source adapter. The finite group \(\Theta\) records the
exceptional characters for lattice cancellation; membership in that
group is not used in place of the child's eligibility predicate.

\textbf{Nonexceptional children.} On this exact nonexceptional part, use
Cauchy to form positive child norms, then split a centered child into
its two rectangles. Only now bound the common row weights and extracted
scalars by one and enlarge to the declared ball. For a retained
rectangle with pre-allocation plain log scales \(x_i\), post scales
\(y_i=x_i-a_{j,\allowbreak{}i}\), set \(\pi_0=\sum_i(-x_i)_+\),
\(\pi=\sum_i(-y_i)_+\). A rectangle with a plain upper support endpoint
below the first ideal is identically zero and is omitted; otherwise
\(0\le\pi_0\le\pi\le\theta_N\). Its clipped reduction and total are \[
r_{{\rm clip},j}=r_j-(\pi-\pi_0)\ge0,\qquad
A_{{\rm clip},j}=\alpha_j+e_j-r_j+\pi,
\tag{B3.68}
\] and its exact scalar relative to the pre-normalizer is
\(Z^{(-r_{{\rm clip},\allowbreak{}j}+e_j+\pi_0)/2}\). For a paired divisor dyad, the
count and both scalars satisfy \[
t_--\frac{r_{{\rm clip},1}+r_{{\rm clip},2}}2
 +\frac{e_1+\pi_{0,1}+e_2+\pi_{0,2}}2\le2\theta_N.
\tag{B3.69}
\] Substitute (B3.68) and use (B3.67); there are only four rectangle
pairs. No common clipped reduction is assigned to a signed difference.

When no slot survives, (B3.69) and the already completed smaller-width
zero-slot assertion fit (B3.64), since \(\Delta\ge0\). For positive
slots, before the final allocation the child total on side one is
\(A'=A-c-w-c_2\). Direct substitution in (B3.63) gives \[
(A'-M')-(A-M)
=g+w-d-c_2-(K_0-K)-w_o+g_2-t_2
\le6(w+\ell)+\delta_1.
\tag{B3.70}
\] Use \(g_2-t_2\le c_2\). In the amplified main norm
\(g\le d+(K_0-K)+6\ell+\delta_1\); in an error
\(g\le d+(K_0-K)+\sigma+\delta_1\), and the three rows of the extraction
table give \(w-w_o+\sigma\le6w\). The other side is identical with
\(c_2,\allowbreak{}d_2\) exchanged.

Delete the extra mask by (B3.41). For each resulting actual natural
child with surviving slots, let
\(F_{\rm act}=(A_{\rm act}-M'_{\rm act}+(6\kappa-1)z_{\rm act})_+\).
Equations (B3.68)--(B3.70), the parent affine condition, and the
aggregate endpoint bound imply \[
F_{\rm act}\le6(w+\ell)+\delta_1+2\theta_N.
\tag{B3.71}
\] Remove whole slots using (B3.47) until the affine defect disappears
or no slot remains. One slot of length \(d\) reduces that defect by
\(6\kappa d\) and costs \(\kappa d\) in squared norm. A single greedy
prefix overshoots its required total length by at most one mesh, so \[
\kappa d_z\le F_{\rm act}/6+\kappa\eta_{\rm mesh}
\le\Delta+\delta_1/6+\theta_N/3+\eta_{\rm mesh}.
\tag{B3.72}
\] If slots remain, apply the smaller-width positive-slot assertion.
Otherwise apply the already completed unrestricted zero-slot assertion;
its plain lengths may now be arbitrary. Combining (B3.69), (B3.72), and
\(M'_{\rm act}-M'=\delta_2\), the strict edge costs at most
\(\delta_2+\delta_1/6+7\theta_N/3+\eta_{\rm mesh}\), besides its
allocated small-power share. Cauchy averages the two costs; it does not
double the mesh. This handles all nonexceptional children at smaller
width.

\hypertarget{b3.9.-fixed-character-rows-and-rectangle-cancellation}{%
\subsection{B3.9. Fixed-character rows and rectangle
cancellation}\label{b3.9.-fixed-character-rows-and-rectangle-cancellation}}

Call a child exceptional if it satisfies the complementary
fixed-inducing predicate above. Before taking absolute values, a nonzero
genuine second allocation is disjoint from every old moving support:
those primes are common column zeros. Thus no old character at such a
prime can cancel a factor supplied by \(G_cV_{\rm id}\). Let \(v_1\) be
the radical length of nonunit-partition primes whose common multiplicity
is one, and set \(f=2v_1\). At each such prime the fixed row factor
\(G_cV_{\rm id}\) has valuation two. Its primitive inducing character
has conductor dividing the fixed modulus and is therefore unramified at
these good primes. Exceptionality forces \(v_p(h')\equiv4\pmod6\). More
generally, at every good prime it specifies one residue class modulo
six, with residue zero outside the frozen moving supports. Unit and
\(S\)-valuations contribute finitely many sectors, while their sixth
powers remain in the free sixth-power factor. Therefore \((h')=h_0v^6\)
for a fixed sixth-power-free \(h_0\), and
\(q_{h_0}\ge Z^{4v_1}=Z^{2f}\). Ideal counting gives the stronger
exponent \((m'_{\rm act}-2f)/6\); we need only \[
\#\{h'\text{ exceptional}:q_{h'}\ll Z^{m'_{\rm act}}\}
\ll Z^{(m'-f)/6+\delta_2/6+o(1)}.
\tag{B3.73}
\] Below the forced initial norm the range is empty, apart from a fixed
endpoint factor. This counts the full fixed-inducing part, including the
principal character, and also rows admitted when a former partition
restriction is bounded by one. The corresponding source estimates are
\sourcecode{chosen\_\allowbreak{}exceptional\_\allowbreak{}block} in
\sourcecode{OAI.\allowbreak{}Number\allowbreak{}Theory.\allowbreak{}Dirichlet\allowbreak{}L.\allowbreak{}Moments.\allowbreak{}Second\allowbreak{}Exceptional\allowbreak{}Fixed\allowbreak{}QChosen\allowbreak{}Block}
and \sourcecode{original\_\allowbreak{}exceptional\_\allowbreak{}energy} in
\sourcecode{OAI.\allowbreak{}Number\allowbreak{}Theory.\allowbreak{}Dirichlet\allowbreak{}L.\allowbreak{}Moments.\allowbreak{}Second\allowbreak{}Exceptional\allowbreak{}Fixed\allowbreak{}QSource}.
No saving from the unit set \(t_2\) is used.

After the full Möbius allocation, take the pointwise absolute product on
these rows. Absolute volume on side \(j\), with its raw normalizer and
live slots, is \(O(Z^{(\alpha_j-r_j)/2+e_j})\). Together with both
extraction scalars and the divisor count, (B3.67) gives paired volume at
most \(Z^{a_0-b_2+2\theta_N+o(1)}\): indeed \(t_--r_1-r_2\le\omega_2\),
and \(r_1\ge0\). Its nominal excess over the inner allowance (B3.64),
including (B3.73), is exactly \[
\frac{m'-f}{6}+a_0-b_2-(M'+\Delta)
=A-\frac56M-F_1-F_2,
\tag{B3.74}
\] where \[
\begin{aligned}
F_1&=c/6+5\{d+(K_0-K)\}/6+w/3+\widetilde q/6+w_o+B_c-5g/6+\ell,\\
F_2&=2b_2-5g_2/6-p_2+t_2+V/6+f/6.
\end{aligned}\tag{B3.75}
\] These expressions obey \[
F_1\ge\frac23(c+w)-3\sigma-5\delta_1/6,
\qquad F_2\ge\frac23b_2.
\tag{B3.76}
\] For the first, initially take \(g=J_+,\allowbreak{}\ell=0\), and put
\(D_0=d+K_0-K\). If \(J\ge0\), substitution gives
\(F_1=c+2w+\widetilde q/6+w_o/6+B_c\). If \(J<0\), use
\(B_c+\widetilde q/6\ge(3c-5D_0)_+/6-5\delta_1/6\), since
\(\widetilde q\ge R\). Hence
\(F_1\ge[c+5D_0+(3c-5D_0)_+]/6+w/3-5\delta_1/6\ge2(c+w)/3-w/3-5\delta_1/6\).
Restoring \(g\le J_++2\sigma\), \(\ell\ge0\), and \(w\le7\sigma/3\)
loses less than \(3\sigma\). For the second, the local contributions in
units of \(\log_ZP\) are:

\begin{longtable}[]{@{}lrr@{}}
\toprule\noalign{}
Common case & \(F_2\) & \(b_2\) \\
\midrule\noalign{}
\endhead
\bottomrule\noalign{}
\endlastfoot
\(i=i,\allowbreak{}6\nmid i\), unit & \(7i/6\) & \(i\) \\
\(i=i,\allowbreak{}6\nmid i\), nonunit & \((7i-5)/6+\mathbf1_{i=1}/3\) & \(i\) \\
\(i=i,\allowbreak{}6\mid i\) & \(7i/6-1\) & \(i\) \\
\(i>j_0,\allowbreak{}6\mid j_0\), unit & \(i+j_0/6-1\) & \((i+j_0)/2\) \\
\end{longtable}

Each first entry is at least two thirds of the second: the only smallest
case is \(i=1\) in the second row, where the \(f/6\) contribution makes
equality. This proves (B3.76). Consequently (B3.74) is at most
\(3\sigma+O(\xi+\theta_N)\) when \(A\le5M/6\). This completes the
current uncentered stage and validates its comparisons from §B3.5.

\textbf{Masked lattice cancellation.} For a fixed
\(\vartheta\in\Theta\), a squarefree mask \(R_*\) of polynomial norm, a
fixed annular \(W_i\), and real \(t\), put
\(L_i(X,\allowbreak{}t)=\sum_l\vartheta(l)\mathbf1_{(l,\allowbreak{}R_*)=1}q_l^{it}W_i(q_l/X)\).
Then, uniformly for \(X>0\), \[
L_i(X,t)=c_{\vartheta,R_*}X^{1+it}I_i(t)
 +O(Z^{\epsilon_1}\langle t\rangle^J),\qquad |L_i(X,t)|\ll X,
\tag{B3.77}
\] where \(c_{\vartheta,\allowbreak{}R_*}\) is independent of \(X,\allowbreak{}t\), and \(I_i(t)\)
is the fixed lattice measure constant times
\(\int_0^\infty W_i(y)y^{it}dy\).

To prove it first omit the mask. The primary ideals outside \(S\), with
fixed weight \(\vartheta\), form a finite weighted set of classes in a
fixed lattice. Poisson for \(q_z^{it}W_i(q_z/X)=X^{it}f_t(z/\sqrt X)\)
gives zero-frequency main term \(c_\vartheta X^{1+it}I_i(t)\).
Integration by parts on the fixed annulus gives
\(|\widehat f_t(\xi)|\ll_J\langle t\rangle^J(1+|\xi|)^{-J}\). For
\(X\ge1\), summing nonzero frequencies bounds the error by
\(O(\langle t\rangle^J)\) for fixed sufficiently large \(J\); for
\(0<X<1\), both the finite lattice sum and main term are \(O(1)\).
Insert the mask with
\(L_i(X,\allowbreak{}t)=\sum_{d\mid R_*}\mu(d)\vartheta(d)q_d^{it}L_i^0(X/q_d,\allowbreak{}t)\),
omitting already excluded primes. Its main coefficient is
\(c_{\vartheta,\allowbreak{}R_*}=c_\vartheta\sum_{d\mid R_*}\mu(d)\vartheta(d)/q_d\);
the two norm powers of \(d\) cancel. The divisor bound gives the stated
error. Finally absolute counting on a nonempty annulus is \(O(X)\);
below its first possible ideal it is empty. This proves (B3.77).

For \(U_i=X_i/q_{b_i}\), \(V_i=Y_i/q_{b_i}\) and \(T=U_1U_2=V_1V_2\),
the sum of (B3.66) is \(L_1(U_1,\allowbreak{}t)L_2(U_2,\allowbreak{}t)-L_1(V_1,\allowbreak{}t)L_2(V_2,\allowbreak{}t)\). Its
two main terms agree. If all four scales are at least \(Z^r\),
\(r\ge0\), (B3.77) bounds the centrally normalized difference by \[
O\bigl(Z^{\epsilon_1}\langle t\rangle^{2J}T^{1/2}Z^{-r}\bigr).
\tag{B3.78}
\] Each one-error term uses one scale at most \(TZ^{-r}\); the two-error
term fits because \(1\le TZ^{-r}\). Without a nonnegative lower length
the all-scale volume fallback is \(O(T^{1/2})\). This is an
equality-of-main-terms argument for the same row, mask and norm power;
it would fail for independently modified rectangles.

Apply (B3.78) to the centered exceptional children. Initially all four
plain lengths are at least \(L\). Before the final \(t\)-allocation, the
two sides have lower lengths at least \(L-c-w-c_2\), \(L-c-w-d_2\). Let
\(r=(L-c-w-\min(c_2,\allowbreak{}d_2))_+\), and call the side with this saving side
one. Its final allocation removes at most \(r_1\) from either plain. If
\(r_1<r\), use (B3.78) with saving \(r-r_1\); otherwise use volume. Keep
the raw equal-product scales, with no clipping of a signed rectangle.
The divisor dyad and both reductions give exactly \[
t_--r_1-r_2-(r-r_1)_+
=(t_--r_2)-\max(r,r_1)\le\omega_2-r.
\tag{B3.79}
\] Thus the paired absolute volume improves to
\(Z^{a_0-b_2-r+2\theta_N+o(1)}\), uniformly in the live slots. Their
labels were absent from the Fourier density; sum them only now by
absolute volume, using the exact normalization (B3.67).

Put \(v=c+w+\min(c_2,\allowbreak{}d_2)\). From (B3.76),
\(F_1+F_2\ge2v/3-3\sigma-5\delta_1/6\). The remaining centered
exceptional deficit is bounded by \[
A-5M/6-2v/3-(L-v)_+\le(A-M)_+,
\qquad L=M/4.
\tag{B3.80}
\] The left side increases for \(v\le L\), decreases for \(v\ge L\), and
equals \(A-M\) at \(v=L\). In the positive-slot range \(A\le M\); in the
padded zero-slot core \(A\le M+\delta\). The exceptional terminal cost
is therefore at most \(3\sigma+\delta\), with its numerical support and
frequency errors already reserved. This finishes the centered stage.

\hypertarget{b3.10.-closure-and-uniform-common-mesh}{%
\subsection{B3.10. Closure and uniform common
mesh}\label{b3.10.-closure-and-uniform-common-mesh}}

Every strict child has nominal width at most \(M-\sigma\) by (B3.63).
The actual enclosing width and support corrections cost at most
\(C_*\xi\), for one numerical \(C_*\) independent of the eventual slot
count. This follows from the single aggregate \(\theta_N\), the two
frequency enclosures, the at most two plain clippings on a rectangle,
and the fixed list of aggregate lengths in (B3.53)--(B3.80). Every
ledger is affine or a positive part of an affine expression in those
lengths. In particular the positive-slot boundary changes by at most
\(|\Delta M|+|\Delta A|+5|\Delta z|\), because \(0<6\kappa-1\le5\). Thus
\(C_*\) does not grow with the slot count. The one greedy prefix
contributes one mesh, not one mesh per removed slot.

Choose \(\rho,\allowbreak{}\sigma,\allowbreak{}\delta\) from the bounded real ranges and final
\(\epsilon\), with \(\delta\ll\min(\rho,\allowbreak{}\sigma)\), so that \[
T_{\rm term}:=\rho+\delta+\rho/6+3\sigma+5\sigma/3<\epsilon/4.
\] Put \(D=2+\lceil2M_{\max}/\sigma\rceil\), and, still before
specifying the fixed slot system, choose \[
\begin{gathered}
\xi<\min\{\delta/2,\rho\mathfrak m(7/10)/2,\sigma/(4C_*),
                   \epsilon/(16C_*D)\},\\
\eta_{\rm mesh}<\min\{\sigma/6,\epsilon/(16C_*D)\},\\
\epsilon_0<\epsilon/(16C_*D).
\end{gathered}
\tag{B3.81}
\] These choices preserve the comparison margin, padded core and strict
width decrease of at least \(\sigma/2\). Every strict child leaves its
width band; there are at most \(D-2\) such calls on a branch. Same-band
comparisons use only the preceding uncentered stage. Zero-slot
reflection is done at most once per plain before centering, and
amplifier errors are not amplified again. Thus the induction is acyclic.

Now fix the arithmetic data, profiles and any finite slot count. Form
\(H_N\) and the finite support boxes through depth \(D\); choose \(Z\)
eventually so large that \(\theta_N\ll\xi\), the whole-dyad tail
inclusions hold and the pool avoids all slot windows. In each stage
allocate all divisor, ray, logarithmic and prime-density losses within
\(\epsilon_0\). For example, mask deletion can assign each of its at
most \(N+2\) masses a share \(\epsilon_0/(10(N+2))\); divisor
allocations use global bounds
\(C_N^{\omega(v)},\allowbreak{}d_{N+2}(v)^C\ll_{N,\allowbreak{}C,\allowbreak{}a}q_v^a\), with one sufficiently
small \(a\), rather than a loss per prime. The prime contour
displacement costs \(2ez\), so it also fits a share independently of the
number of slots. These subsidiary choices affect constants, not the
previously fixed mesh.

When at most \(d\) strict calls remain, permit exponent error
\(E_d=T_{\rm term}+(d+1)C_*(\eta_{\rm mesh}+\xi+\epsilon_0)\). The
width-floor, diagonal and exceptional calculations fit \(E_0\); one
strict edge uses its child envelope plus at most the displayed stage
allowance, by (B3.69)--(B3.72). Same-band comparisons and reflections
have bounded coefficient mass included in that allowance. Cauchy
averages child errors, and triangle inequalities take the larger one, so
the terminal loss occurs once along a branch. Formula (B3.81) gives
\(E_{D-2}<\epsilon\). This proves all zero-slot bands and then all
positive-slot bands simultaneously for every surviving subset of the
original slots.

Finally choose raw kernel-tail orders and all finite profile and height
orders backwards by the common-profile lemma; then take one final
threshold for their tails, endpoint ratios and logarithmic losses. Each
child extra mask is supported on frozen columns or one frozen amplifier
prime, so its norm remains polynomial through the finite depth.
Exceptional-row masks combine these with the row radical; (B3.77) is
uniform over them. Full moving supports remain explicit, including
canceled zero powers. This proves (B3.40) with its stated uniformity.

For positive slots the comparison margin is bounded below by
\(\mathfrak m(7/10)>0\), and \(6\kappa-1\le7/2\) throughout the proved
range. The cost-to-capacity ratio is \((\kappa d)/(6\kappa d)=1/6\). The
common mesh therefore works for \(7/10\le\kappa\le3/4\), with the
hypotheses and canonical common-profile family in (B3.40). The zero-slot
closure is unconditional and restores the fixed nonprincipal rows
separately. ∎

\hypertarget{appendix-b4.-the-probes-local-factors-and-conductor-allocation}{%
\section{Appendix B4. The probe's local factors and conductor
allocation}\label{appendix-b4.-the-probes-local-factors-and-conductor-allocation}}

The physical probe is defined in (B0.21), and its Euler representation
and local table are (B0.28)--(B0.30). We derive that table here,
including shared primes, then verify the coefficient estimates used to
average its numerator. The notation \(g_\chi(a,\allowbreak{}h)\) denotes the
unnormalized Gauss sum.

\hypertarget{b4.1.-full-modulus-poisson-summation}{%
\subsection{B4.1. Full-modulus Poisson
summation}\label{b4.1.-full-modulus-poisson-summation}}

For \(A=cn^3\), define the finite convolution \[
\mathcal F(a,A,H)=\sum_{\substack{d\bmod a\\H\equiv b_*Ad\pmod a}}
\chi_a(d)g_{\xi\chi_A}\!\left(b_*A,\frac{H-b_*Ad}{a}\right).
\tag{B4.1}
\] Expanding the \(a\)-Gauss sum and writing \(m=r+b_*A l\) gives \[
\sum_{m\bmod ab_*A}\xi(m)\chi_A(m)q_a^{-1/2}g_{\chi_a}(a,-m)e(Hm/(ab_*A))
=q_a^{1/2}\mathcal F(a,A,H).
\tag{B4.2}
\] The \(l\)-sum imposes the congruence in (B4.1); the \(r\)-sum is its
remaining Gauss sum. This calculation uses the full modulus and requires
no coprimality between \(a\) and \(A\).

The planar dilation contributes \(X/q_A\); with the physical square root
it leaves \(X^{1/2}/(\sqrt{q_{b_*}}q_A)\). Its Fourier argument is
\(Xq_H/q_A\). At each prime of \(S\), the local factor of \(\xi\) is
primitive and its Gauss frequency is a unit multiple of \(H\). Thus all
frequencies meeting \(S\), including zero, vanish. Write the remaining
frequency uniquely as \(H=u r^6\), with \(r\) primary and \(u\)
sixth-power-free retaining the unit of \(H\). No factor six is
introduced. Mellin inversion, followed by \(s=t+1-z\), gives (B0.28). On
the initial lines \((\Re s,\allowbreak{}\Re w,\allowbreak{}\Re z)=(3,\allowbreak{}3,\allowbreak{}2)\), the trivial
finite-sum bound gives the summable majorant \[
C_S\sum_{u,c,n,a,r}q_u^{-2}q_c^{-5/2}q_n^{-7}q_a^{-2}q_r^{-12}.
\tag{B4.3}
\] Together with the Gaussian and Mellin decay, this justifies the
initial interchanges.

We record the phase cancellation because it determines every local
numerator. CRT gives \[
g_{\xi\chi_A}(b_*A,K)=\xi(A)\chi_A(b_*)g_\xi(b_*,K)g_{\chi_A}(A,K).
\] At \(S\), \(K\equiv H/a\); its factor is
\(\sqrt{q_{b_*}}\tau\xi(A)\overline{\xi(u)}\xi(a)\). The normalizing
factors in the probe leave precisely \(\overline{\xi(u)}\). The
remaining \(b_*\)-symbols cancel with \(\chi_a(b_*)\Xi(A)^{-1}\).

For a good prime put \(t_p=v_p(A)\), \(k_p=v_p(a)\),
\(H_p=H/p^{v_p(H)}\), \(a_p'=a/p^{k_p}\), \(A_p'=A/p^{t_p}\). Before
cancellation, its CRT-unit factor is \[
\chi_p(H_p)^{k_p-t_p}\chi_p(a_p')^{t_p}\chi_p(A_p')^{t_p-k_p}\chi_p(b_*)^{t_p-k_p}.
\] Reciprocity (B1.5) gives the exact identity \[
R(A,a)\prod_p\chi_p(a_p')^{t_p}\chi_p(A_p')^{-k_p}
=\prod_pR(p,p)^{t_pk_p}.
\tag{B4.4}
\] Write \(t_p=e_p+3l_p\), \(e_p\in\{0,\allowbreak{}1\}\). After dividing by the
cross-prime factors in \(\gamma_1(c)\), the remaining pair factor for
\(p<r\) is \(R(p,\allowbreak{}r)^{e_pl_r+e_rl_p+l_pl_r}\). By the
quadratic-refinement identity (B1.4), their product is \[
R(c,n)G(n^3)\prod_p\overline{G(p^3)}^{\,l_p}
R(p,p)^{-e_pl_p-l_p(l_p-1)/2}.
\tag{B4.5}
\] The first two factors, together with \(G(c)\), cancel the inserted
\(\overline{G(cn^3)}\). Also
\(\alpha(c)\overline{\alpha(cn^3)}=\overline{\alpha(n)}^{\,\allowbreak{}3}\). These
identities allow \(c\) and \(n\) to share primes.

\hypertarget{b4.2.-recovering-the-six-row-euler-table}{%
\subsection{B4.2. Recovering the six-row Euler
table}\label{b4.2.-recovering-the-six-row-euler-table}}

Fix \(p\notin S\), and use \(Q,\allowbreak{}j,\allowbreak{}\rho,\allowbreak{}\mathfrak a,\allowbreak{}\mathfrak b,\allowbreak{}V,\allowbreak{}R\) from
(B0.29). Put \(\omega_p=\chi_p(-1)\), \(j'=j+6m\), and let \(t=e+3l\),
\(e\in\{0,\allowbreak{}1\}\), be the local completed exponent. The finite sum is
exactly \(C_p(t,\allowbreak{}k,\allowbreak{}j')\) in (B0.30). Prime-power lifting gives \[
g_{\chi_p^h}(p^t,K)=
\begin{cases}
Q^{t-1}g_{\chi_p^h}(p,K/p^{t-1}),&p^{t-1}\mid K,\\
0,&p^{t-1}\nmid K.
\end{cases}
\tag{B4.6}
\] For \(t=0\), \(C_p(0,\allowbreak{}0,\allowbreak{}j')=1\) and \(C_p(0,\allowbreak{}k,\allowbreak{}j')=\mathbf1_{j'=0}\)
when \(k>0\). For \(t>0,\allowbreak{}k=0\), it is the source Gauss sum. For
\(t>0,\allowbreak{}k=1,\allowbreak{}j'\ge1\), the change of variable in the prime Gauss sum gives
\[
C_p(t,1,j')=\omega_p^{-1}g_{\chi_p}(p,1)
g_{\chi_p^{t-1}}(p^t,p^{j'-1}).
\tag{B4.7}
\] All other cases vanish. For \(k\ge2\), the last lift
\(d=d_0+p^{k-1}h\) has a nontrivial additive character and sums to zero.
At \(k=1,\allowbreak{}j'=0\), the defining congruence has no solution.

The complete local summand, after (B4.4)--(B4.5), is \[
(-1)^e\gamma_1(p)^{-e}\eta(p)^e
(\mathfrak a\overline{\gamma_3(p)})^l
\rho^{k-t}\omega_p^{tk-el-l(l-1)/2}
C_p(t,k,j+6m)
Q^{-(s+1/2)e-(1+3s)l-wk-6zm}.
\tag{B4.8}
\] Here \(e=0,\allowbreak{}1\), \(l,\allowbreak{}k,\allowbreak{}m\ge0\); there is a single sum over these
indices. The bar on \(\gamma_3(p)\) follows from (B4.5), and cannot be
suppressed unless its square is one.

For clarity, (B4.6)--(B4.7) leave the following cases for \(t>0\):

\begin{longtable}[]{@{}lll@{}}
\toprule\noalign{}
\(t\) & \(k=0\) & \(k=1\) \\
\midrule\noalign{}
\endhead
\bottomrule\noalign{}
\endlastfoot
\(6r+6\) & \(j'\ge t-1\) & \(j'=t\) \\
\(6r+1\) & \(j'=t-1\) & \(j'\ge t\) \\
\(6r+3\) & \(j'=t-1\) & \(j'=t\) \\
\(6r+4\) & \(j'=t-1\) & \(j'=t\) \\
\end{longtable}

Increasing \(r,\allowbreak{}m\) together multiplies the summand by \(R\). The strict
tails in the first two rows therefore sum respectively to \[
\frac{R(1-Q^{-1})}{(1-R)(1-V)},\qquad
-\frac{\eta(p)(Q-1)Q^{-s-w}V^{\mathbf1_{j\le1}}}{(1-R)(1-V)}.
\tag{B4.9}
\] The equalities \(j'=t-1\) or \(t\), using (B1.7)--(B1.10), give the
six values \(J_j/(1-R)\) in (B0.29). The \(t=0\) cases give
\((1-V)^{-1}+\mathbf1_{j=0}W^\circ/(1-W^\circ)\). This derives the
entire table, including \(j=0,\allowbreak{}1\) and the shared-prime terms.

\hypertarget{b4.3.-polynomial-euler-corrections-and-their-weighted-mass}{%
\subsection{B4.3. Polynomial Euler corrections and their weighted
mass}\label{b4.3.-polynomial-euler-corrections-and-their-weighted-mass}}

The averaging estimate (B0.41) uses absolute coefficient mass only
before the marked moment is applied. We verify the precise weights. For
an unramified row prime, put \(v=\chi_p(u)\),
\(d_p=\eta(p)v^{-1}Q^{-s}\), and expand in \(t=Q^{-w}\). Equation
(B0.29) gives \[
\begin{aligned}
H_{p,0}&=\frac{1-R/Q-d_p(1-V)}{(1-R)(1-d_p)},\\
H_{p,1}&=\frac{\eta(p)Q^{-s}(1-QV)-vV+vR/Q}{(1-R)(1-d_p)},\\
H_{p,2}&=\frac{v\eta(p)(Q-1)Q^{-s}V-v^2R(1-V)}{(1-R)(1-d_p)}.
\end{aligned}
\tag{B4.10}
\] Thus on \(\Re s=a+16e,\allowbreak{}\Re z=17/50\), \[
|H_{p,0}-1|+\sum_{k=1}^2|H_{p,k}|Q^{k(a-1)}\ll Q^{-1-\epsilon_e}.
\tag{B4.11}
\] The leading linear exponent is \(-1-16e\); the other possible largest
exponents are \(a-1-6z_0\), \(2-4a-6z_0\), \(4-6a-6z_0\) and
\(-a-6z_0\), all strictly below \(-1\) for \(a\ge.51\).

To check the error bound without a local quotient, write
\(E_p=P_p^*+d_p=E_0+E_1t\), where \[
E_0=\frac R{1-R}\left(\frac{1-Q^{-1}}{1-V}-d_p\right),\quad
E_1=\frac{-\eta(p)(Q-1)Q^{-s}V/(1-V)+vR}{1-R}.
\tag{B4.12}
\] Insert this into (B0.33), using \(c_pd_p=v^{-1}\). Its constant
coefficient is bounded by \(Q^{-a}+Q^{-6z_0}+Q^{4-5a-6z_0}\ll Q^{-a}\).
Multiplication by \(t\) carries the weight \(Q^{a-1}<1\), and the
possible \(QV\) term has exponent \(1-6z_0<-1\). Hence \[
\sum_{k=0}^3|\Delta_{p,k}|Q^{k(a-1)}\ll Q^{-a}.
\tag{B4.13}
\]

At a row prime \(p^j\Vert u\), \(1\le j\le5\), both \(\psi_u(p)\) and
its conjugate vanish, so \(\Delta_p=G_p\). Write \(J_j=A_j+B_jt\) using
the six-row table, and abbreviate \(c=\overline{\eta(p)}Q^s\). For
\(j\ge2\), direct expansion gives \[
\begin{aligned}
H_0&=[1-R/Q+(1-V)A_j]/(1-R),\\
H_1&=[-\eta(p)(Q-1)Q^{-s}+(1-V)B_j]/(1-R),\\
G_0&=c[R(1-Q^{-1})+(1-V)A_j]/(1-R),\\
G_1&=[-Q+R/Q+(1-V)(cB_j-A_j)]/(1-R),\qquad G_2=-H_1.
\end{aligned}
\tag{B4.14}
\] For \(j=1\), \[
\begin{gathered}
H_0=\frac{1-R/Q}{1-R},\\ H_1=\frac{\eta(p)Q^{-s}(1-QV)}{1-R},\\
G_0=\frac{cR(1-Q^{-1})}{1-R},\\ G_1=\frac{R/Q-QV}{1-R},\\ G_2=-H_1.
\end{gathered}
\tag{B4.15}
\] The primitive conductor has exponent one at every such good prime.
Allocate \(\min(k,\allowbreak{}j-1)\) powers of a shift \(Q^k\) to
\(h_r\mid\Lambda_u=q_u/q_{\operatorname{rad}(u)}\); allocate the rest to
\(h_0\). With \(\delta=2a-1\), (B4.14)--(B4.15) prove \[
\sum_k|H_{p,k}|Q^{-k/2}Q^{\delta(k-j+1)_+/2}\le1+CQ^{-c_e},\qquad
\sum_k|\Delta_{p,k}|Q^{-(k+1)/2}Q^{\delta(k-j+1)_+/2}\ll1.
\tag{B4.16}
\] The largest error is \(G_1\asymp Q\) at \(j\ge2\): the factor
\(Q^{-1}\) cancels its size, and its shift belongs entirely to
\(\Lambda_u\). At \(j=2\), the quadratic error contributes
\(Q^{1-a}Q^{-3/2}Q^{\delta/2}=Q^{-1}\). The remaining non-tail constant
or linear terms have exponents \(-\delta\) at \(j=2\),
\(-3\delta/2,\allowbreak{}-\delta\) at \(j=3\), \(-3\delta/2\) at \(j=4\), and
\(-5\delta/2\) at \(j=5\); at \(j=1\), (B4.15) leaves only negative
powers. This checks every row valuation.

Products of (B4.16) cost \(U^\epsilon\), as do the divisor choices for
finitely many ramified error slots. Off the row, an error slot of norm
scale \(P\), after its \(P^{z_0-1/2}\) normalizer is removed, has
weighted mass at most \[
P^{-1/2}\sum_{q_p\asymp P}q_p^{-a}\ll P^{1/2-a}\le1.
\tag{B4.17}
\] The ordinary factors are summable by (B4.11). Multiplying these
estimates proves (B0.44), uniformly in the row, before any moment is
invoked. This is why the subsequent use of MP involves common plain
profiles and permitted prime coefficients only.

\hypertarget{t2.-primitive-reduction-and-the-precise-hecke-family}{%
\section{T2. Primitive reduction and the precise Hecke
family}\label{t2.-primitive-reduction-and-the-precise-hecke-family}}

The family in T0 is the finite-order, trivial-infinite-type family
represented upstream by \sourcecode{Hecke\allowbreak{}Family.\allowbreak{}Character}: a residue
character in a nonzero integral modulus, trivial on global units. No
restriction to norm lifts, real characters, or bounded conductor is made
in the B0--B4 bootstrap. Arbitrary target conductors enter the fixed set
\(S\) only after the common power saving has been chosen. The unit
factors in intermediate sextic rows are retained and need not themselves
be principal.

Let \(\eta\) in modulus \(\mathfrak m\) be induced by the primitive
character \(\eta^*\) of conductor \(\mathfrak f\mid\mathfrak m\). For
\(\Re s>1\), absolute Euler products give \[
L_{\mathfrak m}(s,\eta)=L(s,\eta^*)
\prod_{\substack{\mathfrak p\mid\mathfrak m\\\mathfrak p\nmid\mathfrak f}}
\left(1-\eta^*(\mathfrak p)(N\mathfrak p)^{-s}\right).
\tag{T2.1}
\] Indeed the imprimitive character deletes precisely the Euler factors
at primes dividing the presented modulus but not the primitive
conductor. Analytic continuation gives the same identity away from the
principal pole in the positive half-plane. Each displayed factor is
nonzero there:
\(|\eta^*(\mathfrak p)(N\mathfrak p)^{-s}|\le(N\mathfrak p)^{-\Re s}<1\).
This proves equivalence of primitive and imprimitive zero sets in
\(\Re s>0\), with the pole excluded. It also justifies deleting and
restoring any additional fixed Euler factors at \(S\).

For reflected sums, this scalar nonvanishing observation is not enough.
One must retain the \textbf{combined conductor}: the primitive conductor
times the redundant prime ideal product. B0.5 and B4 allocate
repeated-prime powers against the row powerful part; B3.4 deletes and
restores the redundant factors explicitly. This is why a
primitive-conductor-only estimate cannot replace those arguments. The
exact upstream finite-polynomial deletion identity and primitive
presentation are
\sourcecode{Centered\allowbreak{}Moment\allowbreak{}Natural\allowbreak{}Primitive.\allowbreak{}polynomial\_\allowbreak{}natural\_\allowbreak{}deletion}
and \sourcecode{exists\_\allowbreak{}natural\_\allowbreak{}primitive} in
\sourcecode{OAI.\allowbreak{}Number\allowbreak{}Theory.\allowbreak{}Dirichlet\allowbreak{}L.\allowbreak{}Moments.\allowbreak{}Natural\allowbreak{}Primitive}.

Here is the family argument that closes the proof, including its
quantifiers. Under \(\beta_*>b_0\), B0.7 fixes a positive exponent
\(\Delta\) from the common gap, geometry, moment mesh and power reserves
before fixing \(\eta\). For each target it then chooses its finite set
\(S\), finitely many profile and height orders, a sufficiently small
analysis-height exponent, and an eventual threshold. B0.50--B0.51 imply
a common power improvement for the Mellin principal term. Its Mellin
transform is holomorphic farther left, while on its original half-plane
it equals a nonzero Gaussian factor times the nonvanishing Euler
correction divided by \(L^S(s,\allowbreak{}\eta)\). The identity theorem rules out
zeros in that enlarged half-plane; a zero could not be cancelled by that
numerator. Finite factors in T2.1 do not change this conclusion.

Consequently every member of the entire Hecke family has no zero to the
right of \(\beta_*-\Delta'\), for one \(\Delta'>0\) independent of the
member (reduce the reserve if needed). This contradicts the definition
of the supremum; no zero attaining the supremum is required.
Equivalently use upstream
\sourcecode{Hecke\allowbreak{}Zero\allowbreak{}Supremum.\allowbreak{}exists\_\allowbreak{}zero\_\allowbreak{}near\_\allowbreak{}beta} with a positive
\(\varepsilon<\min(\Delta',\allowbreak{}\beta_*-1/2)\). Target-dependent constants
and thresholds are harmless; a target-dependent power saving would not
suffice. We conclude \(\beta_*\le b_0\), hence the Hecke assertion of T0
for arbitrary height. The classical \(\Re s\ge1\) case and the principal
pole are retained throughout.

\hypertarget{t3.-all-dirichlet-transfer-including-imprimitive-characters-and-special-points}{%
\section{T3. All-Dirichlet transfer, including imprimitive characters
and special
points}\label{t3.-all-dirichlet-transfer-including-imprimitive-characters-and-special-points}}

Fix \(q\ge1\) and a complex Dirichlet character \(\chi\) modulo \(q\).
Its norm lift has modulus \(q\mathcal O\), and on ideals prime to that
modulus is \[
\eta_\chi(\mathfrak a)=\chi(N\mathfrak a).
\] The field norm maps residue units modulo \(q\mathcal O\) to units
modulo \(q\). Global units of \(\mathcal O\) have norm one, so this is
an allowed finite-order Hecke character of trivial infinite type. It is
not assumed primitive. Its continuation is the ordinary Hecke ideal
\(L\)-function of T0, as proved by the upstream ideal-series bridge.

Let \(\chi_{-3}\) be the quadratic Dirichlet character of the field
\(K\), and regard \(\chi\chi_{-3}\) with the common modulus used by the
upstream \sourcecode{Short\allowbreak{}Draft\allowbreak{}Hecke\allowbreak{}Bridge.\allowbreak{}base\allowbreak{}Change\allowbreak{}Char}. For \(\Re s>1\),
\[
L_K(s,\eta_\chi)=L(s,\chi)L(s,\chi\chi_{-3}).
\tag{T3.1}
\] To verify every local factor: a rational prime split in \(K\) gives
two norm-\(p\) ideals and the square of \((1-\chi(p)p^{-s})^{-1}\); an
inert prime gives one norm-\(p^2\) ideal and
\((1-\chi(p)^2p^{-2s})^{-1}\), the product of the factors with signs
\(+\) and \(-\); the ramified prime 3 gives one factor, while the
quadratic character vanishes at 3. If \(p\mid q\), the original
character is zero there and both sides omit the corresponding factors.
Thus imprimitive moduli and the ramified prime are included, not hidden
in a primitive convention.

The analytic identity away from \(s=0,\allowbreak{}1\) is proved in the pinned source
by
\sourcecode{OAI.\allowbreak{}Seven\allowbreak{}Eighths.\allowbreak{}Hecke\allowbreak{}Dirichlet.\allowbreak{}LFunction\_\allowbreak{}eq\_\allowbreak{}dirichlet\_\allowbreak{}product}.
Its norm lift, modulus and coefficient identity are the declarations
\sourcecode{character}, \sourcecode{modulus}, and \sourcecode{ideal\allowbreak{}Coeff\_\allowbreak{}character}
in the same file. We use those exact objects.

Now suppose \(b_0<\Re s\) and \(\neg(\chi=1\land s=1)\).

\begin{itemize}
\tightlist
\item
  If \(\Re s\ge1\), use the classical all-Dirichlet nonvanishing theorem
  \sourcecode{Dirichlet\allowbreak{}Character.\allowbreak{}LFunction\_\allowbreak{}ne\_\allowbreak{}zero\_\allowbreak{}of\_\allowbreak{}one\_\allowbreak{}le\_\allowbreak{}re}, with
  the supplied pole exclusion converted by \sourcecode{not\_\allowbreak{}and\_\allowbreak{}or}. This
  includes \(s=1\) for nonprincipal characters. The principal case at
  \(s=1\) is excluded by hypothesis.
\item
  If \(\Re s<1\), positivity of \(b_0\) gives \(s\ne0\), and the strict
  upper bound gives \(s\ne1\). T0 and T2 give nonvanishing of
  \(L_K(s,\allowbreak{}\eta_\chi)\). Equation T3.1 then implies \(L(s,\allowbreak{}\chi)\ne0\):
  otherwise its product with the finite complex value of the companion
  function would vanish.
\end{itemize}

This proves the all-modulus, all-character, arbitrary-height theorem. In
particular, the value at \(s=1\) is never handled by a fallacious
zero-times-pole argument; the companion can be principal there. The
point \(s=0\) is outside the asserted half-plane. There is no
Euler-factor restriction added to the theorem's statement.

The same proof with \(b_0<\theta\) gives the rational statement \[
\begin{gathered}
\forall\{q:\mathbf N\}\,[\mathrm{NeZero}\ q]\ (\chi:\mathrm{DirichletCharacter}\ \mathbf C\ q)
\ \{s:\mathbf C\},\\
\theta<\Re s\to\neg(\chi=1\land s=1)\to\chi.LFunction(s)\ne0.
\end{gathered}
\] The public declarations are \sourcecode{Cycle25.\allowbreak{}dirichlet\_\allowbreak{}nonzero} and
\sourcecode{Cycle25.\allowbreak{}dirichlet\_\allowbreak{}nonzero\_\allowbreak{}catalogue} in
\sourcecode{Cycle25.\allowbreak{}Assembly.\allowbreak{}Final.\allowbreak{}Endpoint}.

\hypertarget{t4.-zeta-half-plane}{%
\section{T4. Zeta half-plane}\label{t4.-zeta-half-plane}}

For the principal character modulo one, upstream
\sourcecode{Dirichlet\allowbreak{}Character.\allowbreak{}LFunction\_\allowbreak{}mod\allowbreak{}One\_\allowbreak{}eq} identifies its
\(L\)-function with \sourcecode{riemann\allowbreak{}Zeta}. Thus \(\zeta(s)\ne0\) for
\(\Re s>b_0\), \(s\ne1\). The ordinary analytic zeta function has its
pole at one. If a Lean statement includes the assigned value there,
Mathlib's separate \sourcecode{riemann\allowbreak{}Zeta\_\allowbreak{}one\_\allowbreak{}ne\_\allowbreak{}zero} handles that
convention; it is not an assertion that the analytic pole is removable.

\hypertarget{background-sources-used-in-b0b4}{%
\section{Background sources used in
B0--B4}\label{background-sources-used-in-b0b4}}

The formal development uses
\href{https://github.com/openai/math/tree/fd4aeeb2ee4fc729c18d98444fed42fd0529eeeb}{OpenAI/math}
at revision \sourcecode{fd4aeeb2\allowbreak{}ee4fc729\allowbreak{}c18d9844\allowbreak{}4fed42fd\allowbreak{}0529eeeb}, and Akash
Levy's
\href{https://github.com/akashlevy/QuasiRiemannTracker/tree/2fd60c0926b66ea18d7436f5ed55250fd006ab5d}{weighted-numerator
development} at revision
\sourcecode{2fd60c09\allowbreak{}26b66ea1\allowbreak{}8d7436f5\allowbreak{}ed55250f\allowbreak{}d006ab5d}. Both upstream
attributions and Apache-2.0 notices are retained. The public endpoint
module is \sourcecode{Cycle25.\allowbreak{}Assembly.\allowbreak{}Final.\allowbreak{}Endpoint}; the coefficient,
reflection, moment and source adapters named above identify the analytic
dependencies. The following literature supplies the background and
normalization conventions:

\begin{itemize}
\tightlist
\item
  {[}DR{]} Alexander Dunn and Maksym Radziwiłł, \emph{Bias in cubic
  Gauss sums: Patterson's conjecture}, Annals of Mathematics 200 (2024),
  967--1057,
  \href{https://arxiv.org/html/2109.07463v3}{arXiv:2109.07463v3}. Only
  the unconditional theta/cusp material of §5 and Appendix A is used.
  Its conditional prime asymptotic is not used.
\item
  {[}GL{]} Leo Goldmakher and Benoît Louvel, \emph{A quadratic large
  sieve inequality over number fields}, Math. Proc. Cambridge Philos.
  Soc. 154 (2013), 193--212,
  \href{https://arxiv.org/pdf/1112.1642v2}{arXiv:1112.1642v2},
  Definition 1 and Theorem 1.1. B2 checks the quadratic family and
  masks; the corresponding sharp canonical sieve is also located in the
  pinned Lean source.
\item
  {[}HB{]} D. R. Heath-Brown, \emph{Kummer's conjecture for cubic Gauss
  sums}, Israel J. Math. 120 (2000), 97--124,
  \href{https://doi.org/10.1007/s11856-000-1273-y}{Theorem 2}. The
  squarefree cubic sieve is also stated in DR, Theorem 9.1, and proved
  in the pinned \sourcecode{Cubic\allowbreak{}Sieve/\allowbreak{}Sharp.\allowbreak{}lean}.
\item
  {[}GZ{]} Peng Gao and Liangyi Zhao, \emph{Moments of quadratic Hecke
  L-functions of imaginary quadratic number fields}, J. Number Theory
  209 (2020), 359--377,
  \href{https://arxiv.org/html/1707.00091v3}{arXiv:1707.00091v3},
  equation (1.1), for the ordinary primitive trivial-infinite-type
  functional equation. The formal proof uses the pinned Hecke primitive
  reflection declarations.
\item
  {[}TZ{]} Jesse Thorner and Asif Zaman, \emph{A unified and improved
  Chebotarev density theorem}, Algebra \& Number Theory 13 (2019),
  1039--1068,
  \href{https://arxiv.org/html/1803.02823v3}{arXiv:1803.02823v3},
  Theorem 1.1 for the fixed identity extension. Only a fixed-field
  annular prime-ideal count is needed; no uniformity over growing
  extensions is invoked.
\end{itemize}

\hypertarget{authorship-and-provenance}{%
\section{Authorship and provenance}\label{authorship-and-provenance}}

Hailey Collet initiated and directed the research, chose the targets and
sources, and decided what to retain, check and combine. The detailed
analytic derivations and much of the drafting were developed with AI
assistance, mainly GPT6-Pro and non-Pro Astra, with contributions from
GPT 6.1 Sol, Grok 4.7 and Gemini 3.1 Pro. AI assistance was also used
for Lean formalization. These model names are the author's account of
the tools used, rather than a per-message telemetry audit. Some research
prompts were AI-written and selected or adapted by the author.

The contribution is offered under Apache-2.0, with the OpenAI/math and
weighted-numerator source licenses and attributions retained. Source
attribution and proof verification are separate: the accompanying
reproduction record reports the checks completed for the submitted
revision.

\end{document}
